\providecommand{\dirfiles}[1]{\par\smallskip\noindent
  {\footnotesize\raggedright\emph{Files:} #1.\par}\smallskip}
\providecommand{\dirP}{\mathbb{P}}
\providecommand{\dirE}{\mathbb{E}}
\providecommand{\dirU}{\mathrm{U}}
\section{Small primes in residue classes: subset products, sign patterns and codes}\label{sec:prd}

Under \cref{src:thm:dgrh} the primes up to a fixed power of $\log q$ are equidistributed in
$(\Z/q\Z)^\times$ and in the values of every character modulo $q$, with an error that is small
uniformly in short dyadic ranges. This section uses that input in two ways. The first five
subsections study subset products: which residues are squarefree products of small primes after
arbitrary deletions (\cref{dir:sec:products}), simultaneous congruences for the factors
(\cref{dir:sec:factorizations}), disjoint products in one modulus (\cref{dir:sec:disjoint}), transition
profiles at every density (\cref{dir:sec:profiles}) and rainbow routing (\cref{dir:sec:routing}). The
last two treat the small primes as samples of characters: primes with a prescribed majority of
quadratic signs, and prime designs after erasures (\cref{dir:sec:signs}), and small primes as codes for
sparse character sums (\cref{dir:sec:codes}). The analytic input is the smoothed prime-sum estimate
\cref{rec:lem:gram} (and \cref{rec:lem:triangle} for the once-integrated kernel), with the exact
certificate argument of \cref{dir:lem:bump} for rational weights; the notation and the conventions
\stG, \stU\ are those of \cref{sec:dir}.

Booker and Pomerance asked for squarefree smooth representatives of every class modulo a prime
\cite{BookerPomerance2017Squarefree}. Klurman, Shparlinski and Ter\"av\"ainen observed under GRH that
unrestricted subset products of all primes up to $(\log q)^B$, $B>4$, cover every unit, and constructed
a common small set of primes for almost all moduli \cite{KlurmanShparlinskiTeravainen2026Generation};
Hughes obtains a first-order optimal size for that covering problem \cite{Hughes2026Few}. Related
results are due to Munsch and Shparlinski on smooth squarefree representatives
\cite{MunschShparlinski2020Smooth}, Martin and Parvardi on subproducts of small residues
\cite{MartinParvardi2022Subproducts}, Matom\"aki and Ter\"av\"ainen on products of a bounded number of
larger primes \cite{MatomakiTeravainen2023Products}, and Pach on fixed-cardinality subset sums in a
finite abelian group \cite{Pach2025Equidistribution}. Erd\H{o}s and R\'enyi proved near-entropy covering
by random unrestricted subset sums in finite abelian groups \cite{ErdosRenyi1965Probabilistic}, Ma and
Tang study the second-order term of that threshold \cite{MaTang2026Erdos}, and Lee and Oh study
zero-sum-free sequences in random subsets of cyclic groups \cite{LeeOh2023ZeroSum}. The results below
add an endpoint cutoff, exact numbers of factors, arbitrary retained sets and relative counts through a
bounded quotient; \cref{dir:thm:factorization} does not optimize its leading constant, while for one
tuple of moduli and prescribed proportions \cref{dir:thm:entropy} does.

\subsection{Squarefree products of small primes after deletions}\label{dir:sec:products}
\sources{squarefree_prime_products; orthogonal_nullity_and_prime_products; grh_entropy_products_2026-11-27}

A positive integer is $y$-smooth if all its prime factors are at most $y$. Booker and
Pomerance asked how small $y$ can be if every nonzero class modulo a prime $p$ must contain a
squarefree $y$-smooth integer, and suggested that GRH should give $y=O((\log p)^2)$
\cite[\S6]{BookerPomerance2017Squarefree}. Squarefreeness matters: a set of primes can generate
the unit group by multiplication while exhausting its elements before a squarefree product
reaches a given class. This subsection answers the question in a stronger form, with a
prescribed number of distinct prime factors, uniform counts, and an arbitrary retained subset of
the primes.

Let $q\ge3$, $G=(\Z/q\Z)^\times$, $L=\log q$, and for a cutoff $y$ let
$\mathcal P=\set{p\le y:\ p\text{ prime},\ p\nmid q}$ and $N=\abs{\mathcal P}$; distinct primes
in the same class modulo $q$ are distinct elements of $\mathcal P$. For
$S\subseteq\mathcal P$ with $m=\abs S$ and $1\le k<m$ let $X_k$ be the product, in $G$, of a
uniformly chosen $k$-element subset of $S$, so that $\binom mk\dirP(X_k=g)$ counts squarefree
integers with exactly $k$ prime factors in $S$ that are congruent to $g$. Put $h=\min(k,m-k)$,
and for $\chi\in\widehat G$ and nonempty $U\subseteq\mathcal P$ let
$b_\chi(U)=\abs{\abs U^{-1}\sum_{p\in U}\chi(p)}$. Characters vanish on nonunits.

\begin{lemma}[Small bias of the full pool; \stG]\label{dir:lem:bias}
Given $\eps>0$ there is $C$ such that for $y=CL^2$ and all sufficiently large $q$,
$\max_{\chi\ne1}b_\chi(\mathcal P)\le\eps$. For $y=L^A$ with $A>2$,
$\max_{\chi\ne1}b_\chi(\mathcal P)\le L^{1-A/2+o(1)}$. In both regimes $N\sim y/\log y$.
\end{lemma}

\begin{proof}
At most $L/\log2$ primes divide $q$, so $N\sim y/\log y$ by the prime number theorem. For
$y=L^A$ the sharp-cutoff estimate $\sum_{p\le x}\chi(p)\ll\sqrt x\log(qx)$ for $\chi\ne1$
\cite[Theorem~13.7, (13.21)]{MontgomeryVaughan2007Multiplicative}, which holds under
\cref{src:thm:dgrh} also for imprimitive characters, gives the second bound. For $y=CL^2$ use
the once-integrated estimate
\[
 \sum_{n\le t}(t-n)\Lambda(n)\chi(n)\ll t^{3/2}\log(2q)\qquad(t\ge2,\ \chi\ne1)
\]
(also \cite[(13.28)]{MontgomeryVaughan2007Multiplicative}): it is $t$ times \eqref{rec:eq:triangle} of
\cref{rec:lem:triangle} for the primitive character inducing $\chi$, whose zero sum is
$\ll\sqrt t\log(2q)$ under \cref{src:thm:dgrh} since $\sum_\rho\abs{\rho(\rho+1)}^{-1}\ll\log(2q)$, and passing to $\chi$
costs $O(tL\log t)$; no factor $\log t$ is lost on the zeros. Fix $0<\eta<1/2$ and let
$w(n)=((y-n)_+-((1-\eta)y-n)_+)/(\eta y)$, which is $1$ up to $(1-\eta)y$ and decreases
linearly to $0$ at $y$. The estimate bounds $\abs{\sum_nw(n)\Lambda(n)\chi(n)}$ by
$O(\sqrt y\log(2q)/\eta)$. Removing prime powers costs $O(\sqrt y\log^2y)$, restoring the last
$\eta y$ of the interval costs at most $(\eta+o(1))y$, and replacing $\log p$ by $\log y$ costs
$\pi(y)\log y-\vartheta(y)=o(y)$, uniformly in $\chi$. Dividing by $N\log y\sim y$ gives
$b_\chi(\mathcal P)\le\eta+O(1/(\eta\sqrt C))+o(1)$; choose $\eta$ and then $C$.
\end{proof}

The bias of the retained set $S$ is controlled by Fourier analysis on $G$ alone. Write
\[
 E_k(\chi)=\dirE\chi(X_k)=\binom mk^{-1}[z^k]\prod_{p\in S}\bigl(1+z\chi(p)\bigr).
\]

\begin{lemma}[Fixed-cardinality coefficients; \stU]\label{dir:lem:circle}
Let $z_1,\dots,z_m$ be complex numbers of absolute value one, $b=\abs{m^{-1}\sum_jz_j}$, and
$1\le k<m$. Then
\[
 \Bigl|\binom mk^{-1}\sum_{\abs T=k}\prod_{j\in T}z_j\Bigr|
 \le\min\set{m+1,\ 8\sqrt{1+h}}\exp\Bigl(-(1-b)\frac{k(m-k)}m\Bigr).
\]
\end{lemma}

\begin{proof}
Put $u=k/(m-k)$. For real $\theta$, $\abs{1+ue^{i\theta}z_j}^2/(1+u)^2
=1-2u(1-\Re(e^{i\theta}z_j))/(1+u)^2$; with $1-x\le e^{-x}$ and multiplication over $j$,
\[
 \prod_j\abs{1+ue^{i\theta}z_j}\le(1+u)^m\exp\Bigl(-\frac{um}{(1+u)^2}(1-b)\Bigr),
\]
and $um/(1+u)^2=k(m-k)/m$. Cauchy's estimate on the circle of radius $u$ gives the bound up to
the factor $(1+u)^m/(u^k\binom mk)$, the reciprocal of the probability that
$\mathrm{Bin}(m,k/m)$ takes the value $k$. This value is a mode, so its probability is at least
$1/(m+1)$. The variance is at most $h$, so by Chebyshev's inequality the interval of radius
$2\sqrt{1+h}$ about $k$ carries mass at least $3/4$ and contains at most $5\sqrt{1+h}$
integers; hence the modal probability is also at least $1/(8\sqrt{1+h})$.
\end{proof}

For $0<\tau<1$ let $H_\tau\le\widehat G$ be the subgroup generated by the characters with
$b_\chi(S)>1-\tau$, let $K=\bigcap_{\chi\in H_\tau}\ker\chi$ and $\mathcal Q=G/K$, so that
$\widehat{\mathcal Q}=H_\tau$ and $d:=\abs{\mathcal Q}=\abs{H_\tau}$. Let $\pi\colon G\to\mathcal Q$
be the quotient map and $\lambda_k(v)=\dirP(\pi(X_k)=v)$. Its values are computed by expanding
$\prod_{p\in S}(1+z[\pi(p)])$ in the group ring $\Z[\mathcal Q][z]$.

\begin{theorem}[Exact support and relative counts; \stG]\label{dir:thm:quotient}
Fix $0<a\le1$ and put $D=\floor{1/a}$. There are $\tau>0$ and $B,C>0$ depending only on $a$
such that the following holds for all sufficiently large $q$. Let $y=CL^2$, $m\ge aN$ and
$h\ge BL$. Then the subgroup $H_\tau$ is the same for all such $k$, $d\le D$, and
\begin{equation}\label{dir:eq:profile}
 \bigl|\varphi(q)\dirP(X_k=g)-d\,\lambda_k(\pi(g))\bigr|\le q^{-2}\qquad(g\in G).
\end{equation}
The support of $X_k$ is exactly $\pi^{-1}(\supp\lambda_k)$, and on it
$\dirP(X_k=g)=d\lambda_k(\pi(g))(1+o(1))/\varphi(q)$. Thus every fibre of $\pi$ that contains
one product contains every residue as a product, with asymptotically equal counts.
\end{theorem}

\begin{proof}
\emph{A quotient of bounded order.} Suppose $b_\chi(\mathcal P)\le\eps$ for every
$\chi\ne1$ and put $\alpha=m/N$. In the probability inner product on functions on
$\mathcal P$, distinct characters have inner products of absolute value at most $\eps$, so any
$t$ of them have a Gram matrix of largest eigenvalue at most $1+(t-1)\eps$. Bessel's inequality
for the indicator of $S$ shows: if $b_\chi(S)\ge u$ for $t$ distinct characters, then
\begin{equation}\label{dir:eq:gram}
 t\alpha^2u^2\le\alpha\bigl(1+(t-1)\eps\bigr),\qquad\text{so}\qquad
 t\le\frac{1-\eps}{\alpha u^2-\eps}\quad(\alpha u^2>\eps).
\end{equation}
On $S$ a character of bias greater than $1-\tau$ is within $\sqrt{2\tau}$, in normalized $L^2$,
of a unimodular constant, and so is its inverse. Telescoping a product shows that every word
of length at most $D$ in these generators and their inverses is within $D\sqrt{2\tau}$ of a
unimodular constant, hence has bias at least $1-D^2\tau$. Since $1/a<D+1$, choose $\tau>0$
with $1-D^2\tau>0$ and then $\eps>0$ with
$(1-\eps)/(a(1-D^2\tau)^2-\eps)<D+1$ and positive denominator. By \eqref{dir:eq:gram} the
ball of radius $D$ in the word metric of $H_\tau$ has at most $D$ elements. A word ball
gains a new element at each step until it is the whole group, so $\abs{H_\tau}\le D$; and
$b_\chi(S)\le1-\tau$ for every $\chi\notin H_\tau$. By \cref{dir:lem:bias} the hypothesis on
$\mathcal P$ holds for $y=CL^2$ with $C=C(a)$.

\emph{Fourier inversion.} The characters in $H_\tau$ factor through $\mathcal Q$ and contribute
exactly $d\lambda_k(\pi(g))$, so
\[
 \bigl|\varphi(q)\dirP(X_k=g)-d\lambda_k(\pi(g))\bigr|\le\sum_{\chi\notin H_\tau}\abs{E_k(\chi)}
 \le q(m+1)e^{-\tau h/2}
\]
by \cref{dir:lem:circle}, since $k(m-k)/m\ge h/2$. As $m\le CL^2$, the choice $B=10/\tau$
gives \eqref{dir:eq:profile}.

\emph{Reachable values are not rare.} For any abelian quotient of order $d$, every nonzero
value of $\lambda_k$ is at least
\[
 \Bigl(\sum_{i=0}^{d-1}\binom ki\binom{m-k}i\Bigr)^{-1}\ge\bigl(d\,m^{2(d-1)}\bigr)^{-1}.
\]
Fix a $k$-subset $T$ with the desired quotient value. For another $k$-subset $U$, pair
$U\setminus T$ bijectively with $T\setminus U$; the swaps along these pairs correct the value of
$U$ to that of $T$. Among $d$ or more swap ratios two partial products, the empty one included,
coincide, and the block between them can be removed; repeating, at most $d-1$ swaps remain.
So every $U$ lies within $d-1$ swaps of a $k$-subset with the target value, and the displayed
sum bounds the size of such a neighbourhood. Since $d\le D$ is fixed and $m$ is
polylogarithmic in $q$, a positive value of $\lambda_k$ is far larger than the error $q^{-2}$,
which gives the exact support and the relative asymptotic.
\end{proof}

For $r=1$ and $A=3$, \cref{dir:thm:disjoint-growing} below contains \eqref{dir:eq:profile} in relative form. The
theorem is kept because its proof (the Gram bound, the word ball and the swap argument) is used by
\cref{dir:thm:disjoint-growing,dir:thm:minority,dir:thm:typical,dir:thm:fibre,dir:thm:allmix,dir:thm:profiles}.

\begin{corollary}[Squarefree representatives with a prescribed number of factors; \stG]\label{dir:cor:squarefree}
Given $0<\delta<1/2$ there are $B,C>0$ such that for all sufficiently large $q$, with $y=CL^2$,
every $S\subseteq\mathcal P$ with $m\ge(1/2+\delta)N$ and every $BL\le k\le m-BL$,
\[
 \#\Bigl\{T\subseteq S:\ \abs T=k,\ \prod_{p\in T}p\equiv g\pmod q\Bigr\}
 =\frac{\binom mk}{\varphi(q)}\bigl(1+O(q^{-2})\bigr)\qquad(g\in G).
\]
In particular every unit modulo $q$ is a product of exactly $\ceil{B\log q}$ distinct primes
below $C(\log q)^2$, which is the smoothness scale proposed by Booker and Pomerance.
\end{corollary}

\begin{proof}
For $a>1/2$, $D=1$ in \cref{dir:thm:quotient}, so $d=1$ and $\lambda_k$ is the point mass. Take
$S=\mathcal P$ for the last assertion. For the finitely many small prime moduli the constant $C$
can be enlarged, since Dirichlet's theorem gives a prime in each nonzero class.
\end{proof}

\begin{remark}[The half and third thresholds; sharpness]\label{dir:rem:thresholds}
For $a>1/3$ the only other possibility is $d=2$, given by one nonprincipal quadratic character
$\chi$. If $S$ contains primes of both signs, a swap of two primes of opposite sign changes the
quotient value, so both values occur and every unit is represented in the range of
\cref{dir:thm:quotient}; if $S$ lies in one sign class, all products lie in the class fixed by $k$.
So for $a>1/3$ full coverage in the range of \cref{dir:thm:quotient} holds exactly when no
nonprincipal quadratic character is constant on $S$. If $s$ of the $m$ primes have sign $-1$, the two quotient masses are
$(1\pm E)/2$ with $E=\binom mk^{-1}[z^k](1+z)^{m-s}(1-z)^s$, so one prime of the opposite sign gives
coverage although the two counts can differ markedly; \cref{dir:thm:fibre} gives both counts with
relative errors, and \cref{dir:thm:disjoint-support} the coverage criterion for several products.
Conversely, for an integer $d\ge2$ take primes
$q\equiv1\pmod d$ and a quotient of $G$ of order $d$, and retain the most populous coset among
the primes of $\mathcal P$. This retains at least $N/d$ primes, confines all products of a fixed
length to one coset, and any quotient explaining this support exactly has order at least $d$.
In particular the majority hypothesis of \cref{dir:cor:squarefree} cannot be lowered to one
half: retaining the primes with $\chi(p)=1$ for the quadratic character modulo a prime $q$
keeps $(1/2+o(1))N$ primes, all of whose products lie in an index-two subgroup.
\end{remark}

For the polylogarithmic pool $y=(\log q)^B$, $B>2$, the same argument runs with explicit
constants and characters treated individually by order. Write
\[
 \rho=\max_{\chi\ne1}b_\chi(\mathcal P),
\]
which is $O_B(\log\log q/(\log q)^{B/2-1})$ under \cref{src:thm:dgrh}. This follows by
partial summation from the estimate
$\sum_{p\le x}\chi(p)\log p\ll\sqrt x\log(2x)\log(2qx)$ of Jao, Miller and Venkatesan
\cite[proof of Theorem~1.1, (2.6)]{JaoMillerVenkatesan2009Expander}, extended to imprimitive
characters by passing to the inducing character.

\begin{proposition}[Explicit constants for a polylogarithmic pool; \stG]\label{dir:prop:explicit}
Let $B>2$ and $y=(\log q)^B$, and let $A\subseteq\mathcal P$ have $M$ elements, $1\le k<M$ and
$h=\min(k,M-k)$.
\begin{enumerate}[label=\textup{(\alph*)}]
\item Fix $0<\eps<1/2$. For all sufficiently large $q$, depending on $B$ and $\eps$, if
$M\ge(1/2+\eps)N$ then
\[
 \max_{g\in G}\Bigl|\frac{\varphi(q)\,\#\set{T\subseteq A:\abs T=k,\ \prod_{p\in T}p\equiv g}}
 {\binom Mk}-1\Bigr|\le8\varphi(q)\sqrt{1+h}\,e^{-\eps h/2};
\]
in particular $h\ge(4/\eps)\log q$ gives relative error $O_\eps(q^{-1}\sqrt{\log q})$.
\item Fix $d\ge2$ and $0<\eps<1-1/d$. If $M\ge(1/d+\eps)N$, then for all sufficiently large $q$
and every character $\chi$ of order at least $d$,
\[
 \bigl|\dirE\chi(X_k)\bigr|\le8\sqrt{1+h}\exp\Bigl(-\frac{d^2\eps}{4C_d}h\Bigr),\qquad
 C_d=\frac{d^2(d^2-1)}6 .
\]
\end{enumerate}
The density thresholds $1/2$ in \textup{(a)} and $1/d$ in \textup{(b)} are sharp.
\end{proposition}

\begin{proof}
(a) For $\chi\ne1$, deleting $N-M$ unimodular summands gives
$\abs{\sum_{p\in A}\chi(p)}\le\rho N+N-M$, so $1-b_\chi(A)\ge(2M-(1+\rho)N)/M\ge\eps$ for large
$q$. Apply \cref{dir:lem:circle} with $k(M-k)/M\ge h/2$ and sum over $\chi\ne1$ by orthogonality.

(b) Choose $\theta$ so that the mean of $z_p=e^{-i\theta}\chi(p)$ over $A$ is real and
nonnegative, and put $F_d(z)=\abs{1+z+\dots+z^{d-1}}^2$ on the unit circle. Since $\chi^j$ is
nonprincipal for $1\le j<d$, positivity gives
$\sum_{p\in A}F_d(z_p)\le\sum_{p\in\mathcal P}F_d(z_p)\le N(d+d(d-1)\rho)$. Telescoping $1-z^j$
gives $1-\Re z^j\le j^2(1-\Re z)$, hence $F_d(z)\ge d^2-C_d(1-\Re z)$. So
$1-b_\chi(A)\ge[d^2M-N(d+d(d-1)\rho)]/(C_dM)\ge d^2\eps/(2C_d)$ for large $q$, and
\cref{dir:lem:circle} applies.

For sharpness use the construction of \cref{dir:rem:thresholds} with the kernel of a character $\chi$
of exact order $d$ modulo a prime $q\equiv1\pmod d$: averaging $1,\chi,\dots,\chi^{d-1}$ shows that the
retained density is $1/d+o(1)$, and every product has character value one.
\end{proof}

Once the retained density exceeds $1/3$ by a fixed amount, every nonprincipal character except
those of order two therefore has an exponentially small coefficient; this makes no independence
assumption on the retained primes. At the larger cutoff the required number of factors drops
below $\log q$.

\begin{theorem}[Short products; \stG]\label{dir:thm:short}
Let $A>2$, $y=L^A$, $0<a\le1$, $m\ge aN$, and $c_A=2(A-1)/(A(A-2))$, and let $H_\tau$ be the
subgroup of \cref{dir:thm:quotient}. For every fixed $c>c_A$ there is $\gamma>0$, independent of
$q$, $S$ and $k$, such that for all sufficiently large $q$
\[
 h\ge c\frac{L}{\log L}\quad\Longrightarrow\quad
 \sum_{\chi\notin H_\tau}\abs{E_k(\chi)}\le\exp\Bigl(-\gamma\frac L{\log L}\Bigr).
\]
Consequently the quotient law and the exact support of \cref{dir:thm:quotient} hold on this
scale, and for $k=\ceil{cL/\log L}$ every attainable class has a positive squarefree
representative at most $q^{Ac+o(1)}$.
\end{theorem}

\begin{proof}
By \cref{dir:lem:bias}, $N\sim L^A/(A\log L)$ and $\max_{\chi\ne1}b_\chi(\mathcal P)\le\eps_L
=L^{1-A/2+o(1)}$. Let $j=\floor{\log(q-1)/\log y}$. A product of $j$ primes of $\mathcal P$,
repetitions allowed, is less than $q$, so congruence of two such products is equality, and by
unique factorization each ordered $j$-tuple has at most $j!$ congruent partners. Orthogonality
gives
\begin{equation}\label{dir:eq:shortmoment}
 \sum_\chi b_\chi(\mathcal P)^{2j}\le\frac{\varphi(q)j!}{N^j}\le q^{1/A+o(1)},
\end{equation}
since $j=(1/A+o(1))L/\log L$ and $\log(j/N)=-(A-1)\log L+O(1)$. Fix $c'$ strictly between $c_A$
and $c$, and let $s$ be the least even integer at least $c'L/\log L$; as $c_A>2/A$, $s>2j$.
Multiplying the nonprincipal terms of \eqref{dir:eq:shortmoment} by $\eps_L^{s-2j}$ bounds
their $s$th moment by $q^{1/A-(A/2-1)(c'-2/A)+o(1)}=o(1)$, the exponent being negative by the
definition of $c_A$. For even $s$ the $s$th moment counts, with positive weights, ordered pairs
of $(s/2)$-tuples with congruent products, and each such pair from $S$ is one from
$\mathcal P$; so, including the principal term,
\begin{equation}\label{dir:eq:momenttransfer}
 \sum_\chi b_\chi(S)^s\le(N/m)^s\sum_\chi b_\chi(\mathcal P)^s\le a^{-s}(1+o(1)).
\end{equation}

Complementation gives $\abs{E_k(\chi)}=\abs{E_h(\chi)}$. If $h\ge BL$ with $B$ large, the
bound $q(m+1)e^{-\tau h/2}$ from the proof of \cref{dir:thm:quotient} suffices, so let
$cL/\log L\le h\le BL$. For $u\ge b_\chi(S)$ and $\abs z=r=h/(mu)$ one has
$\prod_{p\in S}\abs{1+z\chi(p)}\le\exp(mru+mr^2/2)$, by $1+x\le e^x$ applied to each squared
modulus. Cauchy's estimate and $\binom mh\ge(m/h)^h$ give
\begin{equation}\label{dir:eq:smallcoeff}
 \abs{E_h(\chi)}\le u^h\exp\Bigl(h+\frac{h^2}{2mu^2}\Bigr)\le e^{3h/2}u^h
 \qquad\text{if }u\ge u_0:=\sqrt{h/m}.
\end{equation}
Uniformly in this range $u_0=L^{-(A-1)/2+o(1)}$. The characters with $b_\chi(S)\le u_0$
contribute at most $q(e^{3/2}u_0)^h\le q^{1-c(A-1)/2+o(1)}$, which is exponentially small in
$L$ because $c_A>2/(A-1)$. For $u_0<b_\chi(S)\le\zeta$ use \eqref{dir:eq:smallcoeff} with
$u=b_\chi(S)$ and \eqref{dir:eq:momenttransfer}: the contribution is at most
$(1+o(1))e^{3h/2}\zeta^{h-s}a^{-s}$, where $h-s\ge\kappa h$ for a fixed $\kappa>0$; choosing
$0<\zeta<1-\tau$ small in terms of $a,c,c'$ makes this at most $e^{-h}$. By \eqref{dir:eq:gram}
only $O_{a,\zeta}(1)$ characters have $b_\chi(S)>\zeta$, the full-pool coherence being
$\eps_L=o(1)$; for those outside $H_\tau$ the circle bound gives $(m+1)e^{-\tau h/2}$ each, and
$\log m=O(\log L)$. The minimum positive quotient mass is an inverse power of $m$, so the error
suffices for exact support. Finally $y^k=\exp(Ak\log L)=q^{Ac+o(1)}$.
\end{proof}

For $a>1/2$ the quotient is trivial, so for each fixed $\eta>0$ every unit has a squarefree
representative at most $q^{2(A-1)/(A-2)+\eta}$ composed of primes from any retained fraction at
least $a$ of the primes up to $(\log q)^A$; $A=4$ gives the exponent $3+\eta$. The order
$L/\log L$ of the number of factors is correct: for prime $q$ and $k=(b+o(1))L/\log L$,
Stirling's formula gives $\binom Nk=q^{b(A-1)+o(1)}$, so even the full pool has fewer than
$q-1$ subsets of size $k$ when $b<1/(A-1)$. For typical residues the constant $1/(A-1)$ is attained
for almost all prime moduli (\cref{dir:thm:typical,dir:thm:typicalsharp}); for the pointwise law the
optimal leading constant is not determined.

For prime moduli the large sieve replaces \cref{src:thm:dgrh} outside a sparse exceptional set,
which can be bounded before the retained primes are chosen.

\begin{theorem}[An unconditional version for almost all prime moduli; \stU]\label{dir:thm:almostall}
Fix $0<a\le1$ and $b>0$. There are $B=B(a)>0$ and $C=C(a,b)>0$ with the following property.
Let $y=C(\log Q_0)^2$. For all sufficiently large $Q_0$, apart from at most
$Q_0\exp(-b\log Q_0/\log\log Q_0)$ primes $q\in[Q_0,2Q_0]$, the conclusions of
\cref{dir:thm:quotient} hold for every $S$ of density at least $a$ and every $h\ge B\log q$.
\end{theorem}

\begin{proof}
Write $Q=Q_0$, $L=\log Q$, $\ell=\log L$ and $N=\pi(y)$; for large $Q$ all primes up to $y$ are
eligible. Fix $\xi>0$ small in terms of $a$, and call a prime $q\in[Q,2Q]$ exceptional if some
nonprincipal $\chi$ modulo $q$ has $\abs{\sum_{p\le y}\chi(p)}>\xi N$. Let $j=\floor{2L/\log y}$
and let $a_n$ count ordered $j$-tuples of primes up to $y$ with product $n$; then
$0\le a_n\le j!$ and $\sum_na_n^2\le j!N^j$. Since $y^j\le Q^2$, the multiplicative large sieve
\cite[Theorem~25.15]{Koukoulopoulos2019Distribution}, with an absolute implied constant,
gives
\[
 \sum_{d\le2Q}\frac d{\varphi(d)}\sideset{}{^*}\sum_{\chi\bmod d}
 \Bigl|\sum_{n\le y^j}a_n\chi(n)\Bigr|^2\ll(Q^2+y^j)\sum_na_n^2,
\]
the star restricting to primitive characters. Every nonprincipal character modulo a prime is
primitive, and each exceptional prime contributes at least $(\xi N)^{2j}$, so the number $E$ of
exceptional primes satisfies $E\ll Q^2j!/(\xi^2N)^j$. With $N\sim y/\log y$ and
$j=2L/(2\ell+\log C)+O(1)$ one has $\log(N/j)=\ell+\log(C/2)+o(1)$, and Stirling's formula
gives
\[
 \log\Bigl(Q^2\frac{j!}{(\xi^2N)^j}\Bigr)
 =L-\Bigl(\tfrac12\log C+1+2\log\xi-\log2+o(1)\Bigr)\frac L\ell+O(1).
\]
Choose $C$ so that the coefficient without $o(1)$ exceeds $b$. For every other modulus all
full-pool nonprincipal biases are at most $\xi$; taking $\xi$ at most the tolerance $\eps$ in
the proof of \cref{dir:thm:quotient}, the Gram bound, the circle estimate and the swap argument
give all the conclusions with $B=10/\tau$.
\end{proof}

\begin{corollary}[Short products for almost every prime modulus; \stU]\label{dir:cor:almostshort}
Fix $A>2$, $0<a\le1$ and $c>2/(A-1)$, and let $y=(\log Q_0)^A$. There is $\eta>0$ such that
the conclusions of \cref{dir:thm:short}, with $c$ in place of the constant in the length
threshold, hold for all but $O(Q_0^{1-\eta})$ primes $q\in[Q_0,2Q_0]$.
\end{corollary}

\begin{proof}
Put $L=\log Q_0$ and choose $2/(A-1)<c'<\min(c,4/A)$, possible since $A>2$. Take $s=2j$ even
with $s=c'L/\log L+O(1)$; then $y^j\le Q_0^2$ for large $Q_0$, and the large sieve of the
preceding proof gives
\[
 \sum_{\substack{Q_0\le q\le2Q_0\\ q\text{ prime}}}\ \sum_{\chi\ne1}b_\chi(\mathcal P)^s
 \ll Q_0^2\frac{j!}{N^j}=Q_0^{2-c'(A-1)/2+o(1)},
\]
with exponent below one. Choose $\eta>0$ with right side $O(Q_0^{1-2\eta})$. Outside
$O(Q_0^{1-\eta})$ prime moduli the inner sum is at most $Q_0^{-\eta}$, so the moment including
the principal character is $1+o(1)$ and each nonprincipal bias is at most
$Q_0^{-\eta/s}=o(1)$. These are the two full-pool estimates used in the proof of
\cref{dir:thm:short}; with this $s$, \eqref{dir:eq:momenttransfer} and \eqref{dir:eq:smallcoeff}
apply, the low-bias range is controlled because $c>2/(A-1)$, the middle range because $c>c'$,
and the remaining characters by the Gram and circle bounds. Replacing $\log q$ by $L+O(1)$
does not affect these strict inequalities.
\end{proof}

\paragraph{Typical residues.}
Measuring the error in mean square over the residues, instead of pointwise, halves the required
number of factors, and for almost all prime moduli the resulting constant is sharp. In this paragraph
$X\to\infty$, $q$ ranges over $[X,2X]$, and $L=\log X$, $\ell=\log L$, $y=L^A$ with $A>2$; the pool
$\mathcal P$, the retained set $S$ with $m=\abs S\ge aN$, the product $X_k$, $h=\min(k,m-k)$, the
subgroup $H_\tau$, the quotient map $\pi$ and the quotient law $\lambda_k$ are as above. Put
$\mu_k(g)=\dirP(X_k=g)$ and $F_k(g)=d\lambda_k(\pi(g))$, the density of the law obtained by making each
fibre of $\pi$ uniform without changing its mass.

\begin{theorem}[Mean-square quotient law; \stG\ and \stU]\label{dir:thm:typical}
Fix $A>2$, $0<a\le1$ and $c>0$, and assume one of the following.
\begin{enumerate}[label=\textup{(\roman*)}]
\item \textup{(\stG)} $c>c_G=(A-1)/(A(A-2))$, and every integer $q\in[X,2X]$ is allowed.
\item \textup{(\stU)} $c>c_*=1/(A-1)$; there is $\eta>0$ such that all but $O(X^{1-\eta})$ primes
$q\in[X,2X]$ are allowed, the exceptional set depending on $A,c,X$ but not on $S$ or $k$.
\end{enumerate}
There are $\tau,\gamma>0$, depending only on $A,a,c$, such that for all large $X$, all allowed $q$, all
$S$ with $m\ge aN$ and all $1\le k<m$ with $h\ge cL/\ell$, the subgroup $H_\tau$ is independent of $k$,
$d=\abs{H_\tau}\le\floor{1/a}$, and
\begin{equation}\label{dir:eq:meansquare}
 \frac1{\varphi(q)}\sum_{g\in G}\bigl|\varphi(q)\mu_k(g)-F_k(g)\bigr|^2\le e^{-\gamma h}.
\end{equation}
In particular the total variation distance from the fibre-uniform law is at most $\tfrac12e^{-\gamma h/2}$;
unreachable fibres have mass exactly zero; and in every reachable fibre all but a fraction at most
$e^{-\gamma h/4}$ of its elements satisfy
$\mu_k(g)=F_k(g)(1+\theta_g)/\varphi(q)$ with $\abs{\theta_g}\le e^{-\gamma h/4}$.
\end{theorem}

\begin{proof}
\emph{The full pool.} Choose $c'$ strictly between the threshold and $c$, in case (ii) also with
$c'<2/A$, and put $t=\ceil{c'L/\ell}$. We show that, uniformly in the allowed $q$,
\begin{equation}\label{dir:eq:typicalmoment}
 \sum_{\chi\ne1}b_\chi(\mathcal P)^{2t}\le X^{-\eta},\qquad\rho:=\max_{\chi\ne1}b_\chi(\mathcal P)=o(1),
\end{equation}
for some $\eta>0$. In case (i), \cref{dir:lem:bias} gives $\rho\le L^{1-A/2+o(1)}$, and the collision count
\eqref{dir:eq:shortmoment}, with $\varphi(q)\le2X$, gives $\sum_\chi b_\chi(\mathcal P)^{2j}\le X^{1/A+o(1)}$
for $j=\floor{\log(q-1)/\log y}=(1/A+o(1))L/\ell$. Since $c_G>1/A$ we have $t>j$, so the nonprincipal part
of the $2t$th moment is at most $\rho^{2(t-j)}X^{1/A+o(1)}\le X^{(A-1)/A-(A-2)c'+o(1)}$, with negative
exponent by the choice of $c'$. In case (ii) the pool is common to all primes $q\in[X,2X]$ once $q>y$.
Let $a_v$ count the ordered $t$-tuples of pool primes with product $v$; then $a_v\le t!$,
$\sum_va_v^2\le t!N^t$ and $v\le y^t=X^{Ac'+o(1)}<X^2$, so the multiplicative large sieve, as in the
proof of \cref{dir:thm:almostall}, gives
\[
 \sum_{\substack{X\le q\le2X\\ q\text{ prime}}}\ \sum_{\chi\ne1}b_\chi(\mathcal P)^{2t}\ll X^2\frac{t!}{N^t}
 =X^{2-c'(A-1)+o(1)}.
\]
Choose $\eta>0$ with $2\eta<c'(A-1)-1$. Outside $O(X^{1-\eta})$ primes the inner sum is at most
$X^{-\eta}$, and then $\rho\le X^{-\eta/(2t)}=o(1)$.

\emph{A quadratic Fourier bound.} The Gram argument in the proof of \cref{dir:thm:quotient} needs only
$\rho=o(1)$; it gives $\tau=\tau(a)>0$ with $d\le\floor{1/a}$ and $b_\chi(S)\le1-\tau$ for
$\chi\notin H_\tau$. We claim: for every fixed retained fraction $a'>0$ and gap $\tau'>0$ there is
$\gamma>0$ such that, for every $S$ with $\abs S\ge a'N$ and every set $\mathcal J$ of characters with
$b_\chi(S)\le1-\tau'$ on $\mathcal J$,
\begin{equation}\label{dir:eq:quad}
 \sum_{\chi\in\mathcal J}\abs{E_k(\chi)}^2\le e^{-\gamma h}\qquad(h\ge cL/\ell).
\end{equation}
By complementation $\abs{E_k(\chi)}=\abs{E_h(\chi)}$. If $h\ge BL$ with $B=4/\tau'$, \cref{dir:lem:circle}
bounds the sum by $2X(m+1)^2e^{-\tau'h}\le e^{-\tau'h/2}$. Otherwise $cL/\ell\le h<BL$; then
$\log(h/m)=-(A-1)\ell+O(\log\ell)$ and $t/h\le1-\delta$ for a fixed $\delta>0$. Put $u_0=\sqrt{h/m}$ and fix
$0<z<1-\tau'$ small. Characters with $b_\chi(S)\le u_0$ contribute, by \eqref{dir:eq:smallcoeff}, at
most $2Xe^{3h}(h/m)^h=\exp(L-(A-1+o(1))h\ell)$, which is at most $e^{-\delta_1L}$ for a fixed
$\delta_1>0$ because $c(A-1)>1$ in both cases ($c_G>c_*$). Characters with $u_0<b_\chi(S)\le z$
contribute, by \eqref{dir:eq:smallcoeff} and the moment transfer \eqref{dir:eq:momenttransfer} applied
with the principal term included and $\sum_\chi b_\chi(\mathcal P)^{2t}\le2$, at most
$2e^{3h}z^{2(h-t)}(a')^{-2t}$, which is at most $2e^{-2h}$ once
$3+2\log(1/a')+2\delta\log z<-2$. By \eqref{dir:eq:gram} only $O_{a',z}(1)$ characters have
$b_\chi(S)>z$, and those in $\mathcal J$ contribute $O((m+1)^2e^{-\tau'h})$ by \cref{dir:lem:circle}.
Since $h\gg L/\ell$ and $\log m=O(\ell)$, the three bounds give \eqref{dir:eq:quad}.

\emph{Conclusion.} Fourier inversion with the normalized counting measure on $G$ gives the identity
\[
 \frac1{\varphi(q)}\sum_g\abs{\varphi(q)\mu_k(g)-F_k(g)}^2=\sum_{\chi\notin H_\tau}\abs{E_k(\chi)}^2,
\]
since the terms with $\chi\in H_\tau$ are the inverse transform of the lifted quotient law; so
\eqref{dir:eq:quad} with $\mathcal J=\widehat G\setminus H_\tau$ proves \eqref{dir:eq:meansquare}, and
Cauchy--Schwarz gives the variation bound. On a reachable fibre $F_k\ge m^{-2(d-1)}$ by the swap argument
in the proof of \cref{dir:thm:quotient}; the fibre has $\varphi(q)/d$ elements, so its conditional mean
square error is at most $de^{-\gamma h}$, and Markov's inequality with tolerance $e^{-\gamma h/4}$ bounds
the exceptional fraction by $dm^{4(d-1)}e^{-\gamma h/2}\le e^{-\gamma h/4}$ for large $X$.
\end{proof}

\begin{corollary}[Typical squarefree representatives; \stG\ and \stU]\label{dir:cor:typical}
If $a>1/2$, then in either regime of \cref{dir:thm:typical}, with $k=\ceil{cL/\ell}$, all but a fraction
$e^{-\gamma k/4}$ of the units $g$ modulo $q$ satisfy
\[
 \#\Bigl\{T\subseteq S:\ \abs T=k,\ \prod_{p\in T}p\equiv g\pmod q\Bigr\}
 =\frac{\binom mk}{\varphi(q)}\bigl(1+O(e^{-\gamma k/4})\bigr).
\]
So almost every unit has a squarefree representative with all prime factors in $S$ and of size at
most $X^{\vartheta+o(1)}$, where for every fixed $\eps>0$ one may take $\vartheta=(A-1)/(A-2)+\eps$ under
\cref{src:thm:dgrh} for all moduli, and $\vartheta=A/(A-1)+\eps$ unconditionally for almost all prime moduli.
\end{corollary}

\begin{proof}
For $a>1/2$ the quotient is trivial; $k=o(m)$, so $h=k$; and the product is at most
$y^k=X^{Ac+o(1)}$, with $c$ chosen just above the threshold.
\end{proof}

For $A=4$ this gives representatives below $q^{3/2+\eps}$ for almost all units under
\cref{src:thm:dgrh}, robust under any deletion of fewer than half the primes below $(\log q)^4$,
against the pointwise exponent $3+\eps$ after \cref{dir:thm:short}; for $A=3,5,10$ the typical constants
$(c_*,c_G)$ are $(1/2,2/3)$, $(1/4,4/15)$ and $(1/9,9/80)$, against $4/3$, $8/15$, $9/40$ pointwise.

\begin{theorem}[Sharpness at the entropy threshold; \stU]\label{dir:thm:typicalsharp}
Let $0<c<c_*=1/(A-1)$ and let $q\in[X,2X]$ be prime. For every $S\subseteq\mathcal P$ and every $k$ with
$h\le cL/\ell$, the fraction of units attained by $X_k$ tends to zero, and the total variation distance of
$X_k$ from every law that is uniform on the fibres of a quotient of bounded order tends to one,
uniformly in $q$, $S$ and $k$.
\end{theorem}

\begin{proof}
Complementation maps the support bijectively onto an inverted translate, so the support has at most
$\binom mh\le\binom N{h_0}$ elements with $h_0=\floor{cL/\ell}<N/2$. By $\binom N{h_0}\le(eN/h_0)^{h_0}$ and
$N\sim L^A/(A\ell)$ this is $X^{c(A-1)+o(1)}=o(X)$, while $\varphi(q)\asymp X$. A law uniform on the
fibres of a quotient of order $d=O(1)$ gives each unit mass at most $d/\varphi(q)$, hence mass $o(1)$ to
the attained set.
\end{proof}

Thus $c_*$ is the sharp first-order constant for typical residues for almost all prime moduli, even
without deletions. The same normalization controls many disjoint products at once, even when they
consume any fixed fraction less than one of the retained primes.

\begin{theorem}[Many disjoint short products; \stG\ and \stU]\label{dir:thm:typicaljoint}
In either regime of \cref{dir:thm:typical}, with $a>1/2$, fix $0<\beta<1$. For every allowed $q$ and $S$
with $m\ge aN$, every $r\ge1$ and positive integers $k_1,\dots,k_r$ with $k_i\ge cL/\ell$ and
$\sum_ik_i\le(1-\beta)m$, let $T_1,\dots,T_r\subseteq S$ be a uniformly chosen ordered family of pairwise
disjoint subsets with $\abs{T_i}=k_i$, and $Y_i=\prod_{p\in T_i}p\bmod q$. There is $\nu>0$, depending
only on the fixed parameters, with
\[
 d_{\mathrm{TV}}\bigl(\mathcal L(Y_1,\dots,Y_r),\ \mathrm{Unif}(G^r)\bigr)\le e^{-\nu L/\ell}.
\]
Hence all but a fraction $e^{-\nu L/(3\ell)}$ of target tuples have, up to relative error
$e^{-\nu L/(3\ell)}$, exactly the expected number
$\varphi(q)^{-r}\prod_i\binom{m-\sum_{j<i}k_j}{k_i}$ of representations by pairwise disjoint prime
sets. With $k_i=\ceil{cL/\ell}$ the number of blocks may be as large as $\floor{(1-\beta)m/k_1}$, of
order $L^{A-1}$.
\end{theorem}

\begin{proof}
Choose the family sequentially: $T_i$ is a uniform $k_i$-subset of the remaining pool $S_i$, of size
$m_i=m-\sum_{j<i}k_j\ge\beta m$. Unconditionally $S_i$ is a uniform $m_i$-subset of $S$. Deleting from the
full pool gives $b_\chi(S)\le(1-a+\rho)/a\le1-\kappa$ for $\chi\ne1$, with $\kappa=(2a-1)/(2a)$. Hoeffding's
inequality for sampling without replacement \cite[\S6]{Hoeffding1963Probability}, applied to real and
imaginary parts, and a union bound over at most $2X$ characters show that, outside an event of
probability $8X\exp(-\beta aN\kappa^2/16)$, every nonprincipal bias on $S_i$ is at most $1-\kappa/2$. For
such a good pool, \eqref{dir:eq:quad} with retained fraction $\beta a$, gap $\kappa/2$ and $\mathcal J$ all
nonprincipal characters, the moment input still being that of the full pool, shows, with the constant
$\gamma'$ of \eqref{dir:eq:quad} for these parameters, that the conditional law of $Y_i$ given the history
is within $\tfrac12e^{-\gamma'cL/(2\ell)}$ of uniform; for a bad pool use
the bound one. Compare the hybrid laws whose first $i$ coordinates are the actual products and whose
others are independent uniform: consecutive hybrids differ in total variation by at most the expected
conditional distance, so
\[
 d_{\mathrm{TV}}(\mathcal L(Y),\mathrm{Unif}(G^r))\le r\bigl(\tfrac12e^{-\gamma'cL/(2\ell)}+8Xe^{-\beta aN\kappa^2/16}\bigr),
\]
which is at most $e^{-\nu L/\ell}$ since $r\le m\ell/(cL)=O(L^{A-1})$ and $N\asymp L^A/\ell$. Finally
$\dirE\abs{\varphi(q)^r\dirP(Y=g)-1}=2d_{\mathrm{TV}}$ for uniform $g$, and Markov's inequality gives the
counting statement.
\end{proof}

For equal block lengths below $c_*L/\ell$, \cref{dir:thm:typicalsharp} shows that already the first
coordinate occupies a vanishing fraction of $G$ for prime moduli, so the block-length threshold is sharp
to first order in the unconditional regime. On the modulus $11$ with pool $2,\dots,31$ and blocks
$3,3,2$, the $25\,200$ ordered families reach $999$ of the $1000$ target triples, with exact joint
distance $9431/63000$ from uniform.
\dirfiles{\artifact{grh_entropy_products_2026-11-27/finite_checks.json}}

\begin{corollary}[One short list of primes for many moduli; \stG\ and \stU]\label{dir:cor:list}
Fix $0<\delta<1/2$ and $b>0$. There are $C,D>0$ and, for each sufficiently large $Q$, a list
$U$ of $M=\ceil{D\log Q}$ distinct primes at most $C(\log Q)^2$ with the following property,
where $U_q=\set{p\in U:p\nmid q}$ and $k=\floor{M/4}$: for every $S\subseteq U_q$ with
$\abs S=m\ge(1/2+\delta)\abs{U_q}$ and every unit $g$ modulo $q$, the number of $k$-subsets of
$S$ with product congruent to $g$ is $\binom mk(1+O(Q^{-2}))/\varphi(q)$. Under
\cref{src:thm:dgrh} \textup{(\stG)} one list works for all integers $q\in[Q,2Q]$.
Unconditionally \textup{(\stU)} one list works for all primes $q\in[Q,2Q]$ outside a set of at
most $Q\exp(-b\log Q/\log\log Q)$ exceptions.
\end{corollary}

\begin{proof}
Put $y=C(\log Q)^2$ and $\gamma=\delta/8$. For all relevant $q$, the mean of each nonprincipal
$\chi$ over all primes up to $y$ has modulus at most $\gamma$: under \cref{src:thm:dgrh} by the
proof of \cref{dir:lem:bias}, uniformly since $\log q\le\log(2Q)$, and unconditionally by the
exceptional-set estimate in the proof of \cref{dir:thm:almostall} with $\xi=\gamma$. Choose $U$
uniformly among the $M$-subsets of these primes. Hoeffding's inequality for sampling without
replacement \cite[\S6]{Hoeffding1963Probability}, applied to real and imaginary parts, bounds the
probability that one sample mean differs from its population mean by more than $\gamma$ by
$4\exp(-M\gamma^2/4)$. There are $O(Q^2)$ characters with moduli up to $2Q$, so a large $D$
gives a list with all these sample biases at most $2\gamma$. Since $\omega(q)!\le q$ gives
$\omega(q)=O(\log Q/\log\log Q)$, one has $\abs{U_q}=M-o(M)$, hence $b_\chi(U_q)\le3\gamma$, and
deletion gives $b_\chi(S)\le(1/2-\delta+3\gamma)/(1/2+\delta)=:1-\kappa<1$. For large $Q$,
$h\ge M/5$, and \cref{dir:lem:circle} bounds the sum over nonprincipal characters by
$2Q(M+1)e^{-\kappa M/10}$, which is $O(Q^{-2})$ for large $D$.
\end{proof}

The conditional part is complemented by \cref{dir:cor:onelist}, which under weaker
hypotheses (all $q\le Q$, deletions counted against the whole list, every power saving and a
range of $k$) controls the length of the list only up to a factor two.

Even with unrestricted subset sizes, covering all units modulo a prime $q\in[Q,2Q]$ needs
$2^{\abs U}\ge q-1$, so the order $\log Q$ of the common list is optimal. Earlier work on subset
products of small primes is surveyed at the start of \cref{sec:prd}.

\begin{example}\label{dir:ex:products}
\begin{enumerate}[label=\textup{(\arabic*)}]
\item For $q=101$ and $y=251$ there are $53$ eligible primes. Retain the $26$ with quadratic
character $+1$ and the nine smallest with character $-1$. The $\binom{35}{14}=2\,319\,959\,400$
products of fourteen distinct retained primes reach all $100$ units, with between
$23\,197\,877$ and $23\,201\,277$ representations each; the $26$ residues alone reach only $50$
classes, with between $0$ and $193\,232$ representations.
\item Exhaustive enumeration over all retained sets of the given size $m$, with representation
counts computed both directly and by a recurrence in the residue and the number of primes
used:
\begin{center}\small
\begin{tabular}{rrrrrrr}
\toprule
$q$&$y$&$N$&$m$&$k$&sets missing a unit / all sets&minimum count\\
\midrule
$7$&$47$&$14$&$8$&$3$&$1/3003$&$0$\\
$7$&$47$&$14$&$9$&$3$&$0/2002$&$5$\\
$15$&$37$&$10$&$6$&$3$&$21/210$&$0$\\
$15$&$37$&$10$&$7$&$3$&$0/120$&$2$\\
\bottomrule
\end{tabular}
\end{center}
For $q=15$ exactly five of the ten primes have character $+1$ under the quadratic character
modulo $3$; retaining them gives ten triples of positive sign, and adding two primes of negative
sign gives fifteen positive and twenty negative triples, as the quotient polynomial predicts.
\item Modulo $13$, retain the thirteen primes up to $211$ lying in the subgroup of cubes,
together with the prime $2$. Products of five of them occupy eight of the twelve units, with
counts $\binom{13}5=1287$, $\binom{13}4=715$ and $0$ in the three cubic classes, although both
quadratic signs occur in the retained set.
\end{enumerate}
\dirfiles{\artifact{squarefree_prime_products/results.json},
\artifact{orthogonal_nullity_and_prime_products/checks/prime_products.txt}}
\end{example}




\subsection{Simultaneous congruences for the factors of a squarefree product}\label{dir:sec:factorizations}
\sources{squarefree_factorizations; squarefree_prime_products; disjoint-prime-products; entropy_optimal_prime_factorizations_2026-11-27; grh_growing_factorizations}

Suppose that a squarefree integer is to be split into three factors $F_1,F_2,F_3$, with $F_1$
prescribed modulo $12$, $F_2$ modulo $15$ and $F_3$ modulo $20$, and with a prescribed number of
prime divisors in each. Every two of these moduli share a prime factor, but the three have gcd
one. One set of small primes makes every triple of units occur with almost the same number of
representations; with the moduli $6,9,12$ exactly the triples whose product is congruent to the
whole product modulo $3$ occur. The following theorem treats both cases, with up to
$O(\log Q/\log\log Q)$ factors, and with the moduli and the deletions chosen after the primes.

Write $\dirU(q)=(\Z/q\Z)^\times$, including the trivial group $\dirU(1)$. For a finite set $T$
of $m$ primes and positive integers $k_1+\dots+k_r=m$, an ordered partition
$T=T_1\sqcup\dots\sqcup T_r$ with $\abs{T_j}=k_j$ determines the factors
$F_j=\prod_{p\in T_j}p$ of $P_T=\prod_{p\in T}p$; there are $\binom m{k_1,\dots,k_r}$ such
ordered factorizations.

\begin{theorem}[Simultaneous factorization; \stG]\label{dir:thm:factorization}
Fix $0<\eps<1/2$ and $A>0$. There are $b,c,C,K>0$, depending only on $\eps$ and $A$, such that for
every sufficiently large $Q$ and every integer $2\le R\le c\log Q/\log\log Q$ there is a set $S$
of $n\le CR\log(2Q)$ distinct primes in $[K\log^2(2Q),2K\log^2(2Q)]$ with the following
property. Let $2\le r\le R$, let $1\le q_1,\dots,q_r\le Q$, let $D\subseteq S$ with
$\abs D\le(1/2-\eps)n$, and put
\[
 T=\set{p\in S\setminus D:\ p\nmid q_1\cdots q_r},\qquad m=\abs T,\qquad
 g=\gcd(q_1,\dots,q_r).
\]
Let $k_1+\dots+k_r=m$ with $\min_jk_j\ge b\log(2Q)$, and for $u_j\in\dirU(q_j)$ let
$N(\mathbf u)$ be the number of ordered factorizations with these quotas and
$F_j\equiv u_j\pmod{q_j}$ for all $j$. If $\prod_ju_j\not\equiv P_T\pmod g$ then
$N(\mathbf u)=0$; otherwise
\begin{equation}\label{dir:eq:factorization}
 N(\mathbf u)=\binom m{k_1,\dots,k_r}\frac{\varphi(g)}{\prod_{j=1}^r\varphi(q_j)}
 (1+\theta_{\mathbf u}),\qquad\abs{\theta_{\mathbf u}}\le Q^{-A}.
\end{equation}
The balanced quotas $k_j\in\set{\floor{m/r},\ceil{m/r}}$ are always admissible.
\end{theorem}

Since $Q^{-A}<1$, \eqref{dir:eq:factorization} gives the exact support. The moduli may repeat,
divide one another or share large factors. When all $q_j$ are equal, the first $r-1$ residues
are free and the last is forced; when $g=1$ every tuple occurs. For $r=2$ the problem reduces,
by the Chinese remainder theorem, to one subset product modulo $\lcm(q_1,q_2)$, and the joint
problem begins with three factors. \Cref{dir:thm:minority} below, at retained density $a=1/2+\eps/2$,
contains the theorem up to the constants and with its own prime window $[x,2x]$,
$x=K_\delta\log^2(2Q^2)$; the direct proof given here is shorter and makes the quota explicit,
$b=2(A+6)/\eps$.

\paragraph{A bound for one prescribed partition.}
Let $q_j\le Q$, put $H=\prod_j\dirU(q_j)$ and $M=\lcm(q_1,\dots,q_r)$, and let $T$ be any set
of $m$ labelled integers coprime to $M$, the partition being uniform with the given positive quotas.
For a character $\chi_j$ of $\dirU(q_j)$ write $\widetilde\chi_j$ for its pullback to $\dirU(M)$,
and suppose that $\Delta\in[0,m]$ satisfies
\begin{equation}\label{dir:eq:gap}
 \Bigl|\sum_{p\in T}\widetilde\chi_i(p)\overline{\widetilde\chi_j(p)}\Bigr|\le m-\Delta
\end{equation}
for every pair of labels with unequal pullbacks. Put $k_*=\min_jk_j$.

\begin{lemma}[Fourier coefficient of a partition; \stU]\label{dir:lem:partition}
For a tuple of characters, merge labels with equal pullbacks into $h$ classes $C$ with
characters $z_C$ and masses $w_C=\sum_{j\in C}k_j/m$, and put $V=\prod_j\chi_j(F_j)$. Then
\[
 \abs{\dirE V}\le(m+1)^{h-1}\exp\Bigl(-\Delta\sum_{C<C'}w_Cw_{C'}\Bigr).
\]
If $C_0$ has the largest mass and $t$ labels lie outside $C_0$, then
$\abs{\dirE V}\le\bigl[(m+1)\exp(-\Delta k_*/(2m))\bigr]^t$.
\end{lemma}

\begin{proof}
After merging, every partition into classes of sizes $mw_C$ has the same number of refinements,
and $V$ is constant on them, so the merged partition is uniform. Assign each position a class
independently with probabilities $w_C$ and condition on the counts being exactly $mw_C$. This
count vector is a mode of the multinomial law: moving one count from an overfull class $i$ to
an underfull class $j$ multiplies the mass by $v_iw_j/((v_j+1)w_i)\ge1$. Since there are at most
$(m+1)^{h-1}$ count vectors, the conditioning event has probability at least $(m+1)^{1-h}$. The
unconditioned numerator is a coefficient of $\prod_{p\in T}\bigl(\sum_Cw_CZ_Cz_C(p)\bigr)$, bounded
by the maximum of the product on the torus $Z_C=e^{i\tau_C}$. For unimodular $v_C$,
\[
 \Bigl|\sum_Cw_Cv_C\Bigr|^2=1-2\sum_{C<C'}w_Cw_{C'}\bigl(1-\Re(v_C\overline{v_{C'}})\bigr);
\]
apply $1-y\le e^{-y}$, take square roots and multiply over $p$, using \eqref{dir:eq:gap} for each
pair uniformly in the phases. Finally
$\sum_{C<C'}w_Cw_{C'}=\frac12(1-\sum_Cw_C^2)\ge\frac12(1-w_{C_0})\ge tk_*/(2m)$ and $h-1\le t$.
\end{proof}

\begin{proposition}[Pointwise estimate for mixed moduli; \stU]\label{dir:prop:mixed}
Put $B=(m+1)\exp(-\Delta k_*/(2m))$. For a tuple $\mathbf u$ with $\prod_ju_j\equiv P_T\pmod g$,
the quantity
\[
 E(\mathbf u)=\frac{\abs H}{\varphi(g)}\,\dirP(F_j\equiv u_j\bmod q_j\text{ for all }j)-1
\]
satisfies
\[
 \abs{E(\mathbf u)}\le\frac Q{\varphi(g)}\bigl((1+QB)^r-1\bigr).
\]
\end{proposition}

\begin{proof}
The reductions $\dirU(M)\to\dirU(q_j)$ are onto, and, prime by prime, the intersection of the $r$
inflated character groups consists of the characters inflated from $\dirU(g)$, since characters
factoring through $\ell^a$ form nested subgroups and the minimum of the exponents
$v_\ell(q_j)$ is $v_\ell(g)$. In Fourier inversion on $H$ a common character $\chi$ of $\dirU(g)$
contributes $\chi(P_T/\prod_ju_j)$, so the $\abs H/\varphi(g)$ common tuples contribute
$\varphi(g)/\abs H$ at a compatible $\mathbf u$. For every other tuple fix a largest class. For a
given set of $t\ge1$ labels outside it, the common character of the remaining labels has at most
$Q$ choices, read off from any one remaining modulus, and each of the $t$ other characters has at
most $Q$ choices; \cref{dir:lem:partition} bounds each term by $B^t$. Summing over the
exceptional labels and dividing by $\varphi(g)$ gives the bound.
\end{proof}

Charging a tuple by the number of labels outside its largest class is what keeps the quota of
order $\log Q$ for growing $r$; counting all $\abs H$ tuples at the worst single-label decay
would cost a further factor of order $r$.

\begin{lemma}[One set of primes for all characters; \stG]\label{dir:lem:oneset}
For each fixed $0<\delta<1$ there are $K_\delta,C_\delta,c_\delta>0$ such that for all large $Q$,
with $x=K_\delta\log^2(2Q)$ and any $L$ with
$C_\delta\log(2Q)\le L\le c_\delta\log^2(2Q)/\log\log(3Q)$, there is a set $S$ of distinct primes
in $[x,2x]$ with $L\le n=\abs S\le2L$ and
\[
 \Bigl|\sum_{p\in S}\chi(p)\Bigr|\le\delta n\qquad
 (2\le h\le Q^2,\ \chi\bmod h\text{ nonprincipal}).
\]
\end{lemma}

\begin{proof}
Fix a nonnegative smooth function $W$ supported in $(1,2)$ with $\int W=1$, and put
$w_p=(\log p)W(p/x)$; then $\sum_pw_p=x+o(x)$. For a primitive nonprincipal $\psi$ of conductor
$f$, \cref{rec:lem:gram} with $\xi=1$ gives $\sum_pw_p\psi(p)=O_W(\sqrt x\log(2f))$, and passing from $\psi$
to an induced character modulo $h$ costs at most $\norm{W}_\infty\log h$. Hence
$\abs{\sum_pw_p\chi(p)}/\sum_pw_p$ is at most a constant times $\log(2Q)/\sqrt x$, which is at most
$\delta/4$ once $K_\delta$ and then $Q$ are large.

Put $\mu=3L/2$ and select each prime independently with probability $\mu w_p/\sum w$; these
probabilities are $O(L\log(2x)/x)\le1$ for small $c_\delta$. For each admissible $\chi$ the sum
$Z_\chi$ of $\chi$ over the selected primes has $\abs{\dirE Z_\chi}\le\delta\mu/4$, and
Bernstein's inequality for the real and imaginary parts gives
$\dirP(\abs{Z_\chi-\dirE Z_\chi}>\delta\mu/4)\le4\exp(-c\delta^2\mu)$ and
$\dirP(\abs{n-\mu}>\mu/3)\le2\exp(-c'\mu)$ with absolute $c,c'>0$. There are at most $Q^4$
characters, so a large $C_\delta$ makes the total failure probability less than one. A successful
selection has $L\le n\le2L$ and $\abs{Z_\chi}\le\delta\mu/2\le\delta n$.
\end{proof}

A modulus $q\le Q$ is divisible by at most $\log Q/\log x$ primes of $S$, so even the most
divisible moduli remove only $o(n)$ primes.

\begin{proof}[Proof of \cref{dir:thm:factorization}]
Take $\delta=\eps/2$ in \cref{dir:lem:oneset}, $b=2(A+6)/\eps$,
$C_0\ge\max\set{C_\delta,4(b+1)}$ and $L=C_0R\log(2Q)$; a small $c$ keeps $L$ in the allowed
range. The number $v$ of primes of $S$ dividing some $q_j$ is at most $r\log Q/\log x$, so
$v/n\le1/(C_0\log x)\le\eps/4$ for large $Q$, uniformly in $r$ and the moduli, and
$m\ge(1/2+\eps)n-v\ge(1/2+3\eps/4)n$; in particular balanced quotas exceed $b\log(2Q)$. The ratio
of two distinct inflated characters comes from a nonprincipal character modulo
$\lcm(q_i,q_j)\le Q^2$, and deleting the $n-m$ unavailable primes changes its sum by at most
$n-m$, so
\[
 m-\Bigl|\sum_{p\in T}\widetilde\chi_i(p)\overline{\widetilde\chi_j(p)}\Bigr|\ge2m-(1+\delta)n\ge\eps n ,
\]
which is \eqref{dir:eq:gap} with $\Delta=\eps n$; only pairwise ratios enter, although $M$ may be
much larger than $Q^2$. For large $Q$, $m+1\le Q$ and $r\le Q$, so
$QB\le Q^2\exp(-\eps k_*/2)\le Q^{-A-4}$, $(1+QB)^r-1\le2rQB$, and the bound of
\cref{dir:prop:mixed} is at most $2Q^{-A-2}\le Q^{-A}$. Multiply by the multinomial count.
\end{proof}

The case of two factors, one of them with modulus one, counts subsets with a prescribed product.

\begin{corollary}[One short list for one product; \stG]\label{dir:cor:onelist}
Fix $0<\delta<1/2$ and $A>0$. There are $b,C,K>0$ such that for all large $Q$ some set $S$ of
$n\le2C\log(2Q)$ distinct primes in $[K\log^2(2Q),2K\log^2(2Q)]$ has the following property. For
every $q\le Q$, every $S'\subseteq S_q=\set{p\in S:p\nmid q}$ with
$\abs{S_q\setminus S'}\le(1/2-\delta)n$, every $k$ with $b\log(2Q)\le k\le\abs{S'}-b\log(2Q)$ and
every unit $u$ modulo $q$,
\[
 \#\Bigl\{T\subseteq S':\ \abs T=k,\ \prod_{p\in T}p\equiv u\pmod q\Bigr\}
 =\frac{\binom{\abs{S'}}k}{\varphi(q)}(1+\theta),\qquad\abs\theta\le Q^{-A}.
\]
This range of $k$ contains $\floor{n/4}$.
\end{corollary}

\begin{proof}
Apply \cref{dir:thm:factorization} with $R=r=2$, $q_1=q$, $q_2=1$ and $D=S_q\setminus S'$; then
$g=1$, $T=S'$, and $N(u,1)$ is the displayed count. The proof of the theorem gives
$n\ge4(b+1)R\log(2Q)$ and $\abs{S'}\ge(1/2+3\delta/4)n$, so $\floor{n/4}$ is admissible.
\end{proof}

Compared with the conditional part of \cref{dir:cor:list}, with $2Q$ in place of $Q$: the
hypothesis $\abs{S'}\ge(1/2+\delta)\abs{S_q}$ implies $\abs{S_q\setminus S'}\le(1/2-\delta)n$, the
error $Q^{-A}$ is any power saving, all moduli up to $Q$ are allowed and every quota in a range;
but the length of the list is fixed there exactly and here only up to a factor two, so neither
statement implies the other.

\paragraph{The number of primes and of factors.}
Both the order $R\log Q$ of the size of $S$ and the order $\log Q/\log\log Q$ of $R$ are
sharp.

\begin{lemma}[A sample of signs; \stU]\label{dir:lem:signs}
Let a population of $m$ signs have absolute sum at least $\beta m$, $0<\beta\le1$, and let $Y_k$
be the product of a uniformly chosen $k$-subset. Then $\abs{\dirE Y_k}\ge(\beta/2)^k$ for
$0\le k\le\beta^2m/4$.
\end{lemma}

\begin{proof}
Switch all signs if necessary so that the sum $s$ is at least $\beta m$, and put $a_j=\dirE Y_j$.
With $m_\pm$ signs of each kind and $f(z)=(1+z)^{m_+}(1-z)^{m_-}$, whose $z^j$ coefficient is
$\binom mja_j$, the identity $(1-z^2)f'=(s-mz)f$ gives $(m-j)a_{j+1}=sa_j-ja_{j-1}$ with
$a_0=1$, $a_1=s/m$. If $a_j\ge(\beta/2)a_{j-1}>0$ and $j\le\beta^2m/4$, then
$(m-j)a_{j+1}\ge(\beta m-2j/\beta)a_j\ge(\beta m/2)a_j$, and induction keeps every ratio at least
$\beta/2$.
\end{proof}

\begin{theorem}[Necessity of logarithmic quotas; \stU]\label{dir:thm:factorlower}
Fix $0<\lambda<1/2$ and $A>0$, and put $\beta=\lambda/2$. Suppose that a set of $n$ distinct
primes has the following property at level $Q$: for $r\ge4/\beta^2$ labelled factors, after
every deletion of at most $\lambda n$ primes and removal of the prime $3$ if present, the
factorization law with balanced quotas and moduli $q_1=\dots=q_r=3$ has pointwise relative error
at most $Q^{-A}$ on the compatible tuples, the balanced profile being required to exist. Then,
for all large $n$,
\[
 n\ge\frac A{\log(2/\beta)}\,r\log Q .
\]
\end{theorem}

\begin{proof}
Remove $3$ and delete up to $\floor{\lambda n}$ primes of the minority sign modulo $3$. Either
the surviving signs are constant, or their imbalance has grown by $\floor{\lambda n}\ge\beta n$;
in both cases it is at least $\beta m$. A smallest block has size $k=\floor{m/r}\le\beta^2m/4$,
and its primes form a uniform $k$-subset. Since $r\ge2$, the uniform law on compatible tuples puts
half its tuples on each sign of this block, so summing the pointwise error against the sign gives
$\abs{\dirE Y_k}\le Q^{-A}$. By \cref{dir:lem:signs}, $(\beta/2)^k\le Q^{-A}$, so
$n\ge rk\ge Ar\log Q/\log(2/\beta)$.
\end{proof}

For bounded $r$ coverage alone needs $n\gg\log Q$: take one modulus a prime in $(Q/2,Q]$ and the
others equal to one, and note that the first factor has at most $2^n$ values. So the order
$R\log Q$ in \cref{dir:thm:factorization} is sharp for its balanced-profile guarantee. Primes up
to $C'\log^2Q$ number $O_{C'}(\log^2Q/\log\log Q)$, so the order $\log Q/\log\log Q$ of the
number of factors is sharp at this height and accuracy; the least possible height for a fixed
number of factors is not determined. The deletion threshold is sharp already for two factors of
positive size with modulus $3$: deleting the smaller sign class costs at most $\floor{n/2}$
primes, fixes the sign of each factor, and changing both target signs gives a compatible but
unattained tuple.

\paragraph{The exact leading constant for coverage.}
For a fixed tuple of moduli and fixed proportions of factors, and asking only that every compatible
residue vector occur, the least number of primes is determined to first order. Every coarsening of
the factors imposes a counting obstruction, and the largest of these is also sufficient. Fix $r\ge2$
and positive $\alpha_1,\dots,\alpha_r$ with sum one, let $q_1,\dots,q_r\ge1$, put
$M=\lcm_jq_j$, $g=\gcd_jq_j$ and $\ell_M=\log(2M)$, and for a set partition $\pi$ of $[r]$ into
$\abs\pi\ge2$ blocks put
\[
 g_J=\gcd_{j\in J}q_j,\quad\alpha_J=\sum_{j\in J}\alpha_j,\quad
 N_\pi=\frac{\prod_{J\in\pi}\varphi(g_J)}{\varphi(g)},\quad
 \mathsf H_\pi(\alpha)=-\sum_{J\in\pi}\alpha_J\log\alpha_J,
\]
and
\[
 T(\mathbf q,\alpha)=\max_{\abs\pi\ge2}\frac{\log N_\pi}{\mathsf H_\pi(\alpha)} .
\]
A set $S$ of $n$ distinct primes coprime to $M$ is \emph{universal} at a quota vector $k$ with
$\sum_jk_j=n$ if every tuple $(u_j)\in\prod_j\dirU(q_j)$ with $\prod_ju_j\equiv P_S\pmod g$ is realized
by an ordered factorization with these quotas.

\begin{theorem}[Sharp entropy threshold for coverage; \stG, the lower bound \stU]\label{dir:thm:entropy}
Fix $r$, $\alpha$ and $A>0$. There are $B_0>1$, $C>0$ and $M_0$, depending only on $r,\alpha,A$, such
that for every tuple with $M\ge M_0$ and every real $B\in[B_0,\ell_M^A]$ there is a set $S$ of
\[
 n=\Bigl\lceil T(\mathbf q,\alpha)+C\frac{\ell_M}{\log B}\Bigr\rceil
\]
distinct primes coprime to $M$, all in $[B\ell_M^2,2B\ell_M^2]$, universal simultaneously at every integer
quota vector with $\sum_jk_j=n$ and $\abs{k_j-n\alpha_j}\le1$. Conversely, if some set of $n$ distinct
primes coprime to $M$, of any size, is universal at one such quota vector, then
\[
 n\ge\max_{\abs\pi\ge2}\frac{\log N_\pi+\frac{\abs\pi-1}2\log n}{\mathsf H_\pi(\alpha)}-O_{r,\alpha}(1).
\]
Moreover $T(\mathbf q,\alpha)\asymp_{r,\alpha}\ell_M$. Hence if $B\to\infty$ with $B\le\ell_M^A$, the least
universal size at these rounded proportions is $(1+o(1))T(\mathbf q,\alpha)$; and for each $\eps>0$ a fixed
$B=B_{r,\alpha,\eps}$ gives $n\le(1+\eps)T(\mathbf q,\alpha)+1$ with primes of size $O_{r,\alpha,\eps}(\ell_M^2)$.
\end{theorem}

The statement concerns coverage for one tuple of moduli; it asserts neither comparable numbers of
representations nor one list for different tuples, which \cref{dir:thm:factorization} provides with an
unspecified constant. With $B=\ell_M^\delta$ the additive error is $O(\ell_M/\log\ell_M)$ at prime height
$O(\ell_M^{2+\delta})$, and with $B=\log\ell_M$ it is $O(\ell_M/\log\log\ell_M)$ at height
$O(\ell_M^2\log\ell_M)$; the lower bound exceeds $T$ by order $\log\ell_M$.

\begin{corollary}[Equal moduli and balanced factors; \stG]\label{dir:cor:entropyequal}
If $q_j=q$ and $\alpha_j=1/r$ for all $j$, then $T=(r-1)\log\varphi(q)/\log r$. For $r=2$ a single set of
$(1+o(1))\log_2\varphi(q)$ primes in a dyadic interval at height $B\log^2q$, for any slowly growing $B$,
has the property that its subsets of half its size, rounded either way, represent every unit modulo
$q$; even unrestricted subset products need $\log_2\varphi(q)$ primes.
\end{corollary}

\begin{proof}
Here $N_\pi=\varphi(q)^{\abs\pi-1}$. Among partitions into $h$ blocks the entropy is least for block
sizes $r-h+1,1,\dots,1$, by convexity of $t\log t$, so the minimum is $f(h-1)$ with
$f(t)=\log r-\frac{r-t}r\log(r-t)$. This function is concave with $f(0)=0$ and $f(r-1)=\log r$, so
$f(t)\ge t\log r/(r-1)$, with equality for the partition into singletons.
\end{proof}

Every coarsening matters. For $r=3$, $q_1=q_2=q_3=101$ and quotas $(6,6,1)$ there are $12\,012$ ordered
assignments for $10\,000$ compatible targets, yet the third factor has only $13$ possible values; the
partition $\set{\set{1,2},\set{3}}$ detects this. For $\alpha=(0.49,0.49,0.02)$ and a common modulus,
separating the third label gives $T\ge\log\varphi(q)/\mathsf H(0.02,0.98)$, which exceeds
$2\log\varphi(q)/\mathsf H(0.49,0.49,0.02)$. For coprime moduli the compatibility condition is void, and
for $r=2$ the constant is $\log\varphi(M)/\mathsf H(\alpha_1,\alpha_2)$.

Let $H=\prod_j\dirU(q_j)$, let $\mathcal D$ be the kernel of the product map $H\to\dirU(g)$, and
$d=\abs{\mathcal D}=\prod_j\varphi(q_j)/\varphi(g)$. For any labelled list of units modulo $M$ the compatible
targets form a coset of $\mathcal D$. As in the proof of \cref{dir:prop:mixed}, for every block $J$ the
intersection of the inflated character groups of $\dirU(q_j)$, $j\in J$, is the group inflated from
$\dirU(g_J)$; and $\widehat{\mathcal D}$ is $\widehat H$ modulo the diagonal group of characters inflated
from $\dirU(g)$, which acts freely on lifts.

\begin{proposition}[Every coarsening imposes a counting bound; \stU]\label{dir:prop:entropycount}
If $S$ is universal at $k$, then $\binom n{(k_J)_{J\in\pi}}\ge N_\pi$ for every partition $\pi$ with
$\abs\pi\ge2$, where $k_J=\sum_{j\in J}k_j$.
\end{proposition}

\begin{proof}
Multiply the factors in each block $J$ and reduce modulo $g_J$. The map from compatible fine targets to
compatible coarse targets is onto: given coarse units $v_J$ with product $P_S\bmod g$, lift each $v_J$ to
one label of its block and put the other fine targets of the block equal to one. The coarse targets form
a coset of size $N_\pi$, and each coarse assignment of the $n$ primes with quotas $k_J$ gives one coarse
target.
\end{proof}

With $k_j=n\alpha_j+O(1)$, Stirling's formula gives
$\log\binom n{(k_J)}=n\mathsf H_\pi(\alpha)-\frac{\abs\pi-1}2\log n+O_{r,\alpha}(1)$ uniformly over the finitely
many partitions, which is the lower bound of \cref{dir:thm:entropy}. For the scale, the partition into
singletons has $N_\pi=d\ge\max_j\varphi(q_j)$; as $M\le\prod_jq_j$ some $q_j$ has $\log q_j\ge\log M/r$, and
$\varphi(q)\ge\sqrt{q/2}$, so $\log d\gg_r\ell_M$, while $\log N_\pi\le r\ell_M$ always.

For the upper bound let $\mu$ be a probability measure on a finite labelled pool of units modulo $M$,
with
\[
 b=\max_{\chi\ne1}\bigl|\dirE_{x\sim\mu}\chi(x)\bigr|,\qquad\omega=\max_x\mu(x),
\]
the maximum in $\omega$ being over labels. Draw $X_1,\dots,X_m$ independently from $\mu$, assign the $m$
positions by a uniform ordered partition with positive quotas $\ell=(\ell_j)$, let $\nu_{X,\ell}$ be the law
of the resulting target on the compatible coset, and let $\delta(X,\ell)=1-\abs{\supp\nu_{X,\ell}}/d$ be the
missing fraction.

\begin{lemma}[Collisions stratified by coarsenings; \stU]\label{dir:lem:collision}
For every quota vector $\ell$,
\[
 \dirE_X\delta(X,\ell)\le\sum_{\abs\pi\ge2}N_\pi\frac{(1+(\abs\pi-1)b)^m}{\binom m{(\ell_J)_{J\in\pi}}} .
\]
If $m\omega\le1/2$, the $m$ draws have distinct labels with probability at least $e^{-m^2\omega}$, so for a
finite family $\mathcal L$ of quotas the expected sum of the missing fractions, conditional on distinct
labels, is at most $e^{m^2\omega}$ times the sum of these bounds over $\ell\in\mathcal L$.
\end{lemma}

\begin{proof}
Translate the law to $\mathcal D$ and put $f=d\nu$. A missing element contributes one to
$d^{-1}\sum_z\abs{f(z)-1}^2$, so Parseval gives $\delta\le\sum_{\eta\ne1}\abs{\widehat\nu(\eta)}^2$ over
$\widehat{\mathcal D}$. Lift $\eta$ to $(\chi_1,\dots,\chi_r)\in\widehat H$ and let $\pi$ be the partition of
the labels by equality of the inflations $\widetilde\chi_j$; it is invariant under the diagonal action, a
nontrivial $\eta$ has $\abs\pi\ge2$, and at most $N_\pi$ characters $\eta$ have a given pattern $\pi$,
since a common character of a block descends to its gcd. Up to a unimodular factor
$\widehat\nu(\eta)=\dirE_A\prod_i\widetilde\chi_{A(i)}(X_i)$ for a uniform assignment $A$. Expand
$\abs{\widehat\nu(\eta)}^2$ with two independent assignments $A,A'$ and average over $X$: a position where
the two inflated characters agree contributes one, and otherwise a factor of modulus at most $b$. After
merging labels into the blocks of $\pi$ both assignments are uniform with block quotas $(\ell_J)$. For a
fixed first assignment, the sum of $b^{\#\text{differing positions}}$ over all unrestricted second
assignments is $(1+(\abs\pi-1)b)^m$, and the prescribed quotas only decrease it; divide by the multinomial
number. For the second assertion multiply the conditional probabilities $1-i\omega$, $i<m$, and use
$\log(1-u)\ge-2u$ for $u\le1/2$.
\end{proof}

A nearly full support is repaired by fresh pairs of primes. At any stage remove the used labels and
renormalize; after at most $n$ removals with $n\omega<1$ the bias and largest weight of the remaining law
satisfy
\begin{equation}\label{dir:eq:residual}
 b_R\le\frac{b+n\omega}{1-n\omega},\qquad\omega_R\le\frac{\omega}{1-n\omega}.
\end{equation}
Put $s=2(r-1)$ and $\mathbf a=(1,\dots,1,r-1)$. A repair block is a set of fresh pairs $(A_j,B_j)$,
$j<r$: the baseline places $A_j$ in factor $j$ and all $B_j$ in factor $r$, the alternative swaps every
pair; both add $\mathbf a$ to the quotas, and their targets differ by $V\in\mathcal D$ with
$V_j=B_j/A_j$ for $j<r$ and $V_r=\prod_{j<r}A_j/B_j$.

\begin{lemma}[A common repair step; \stU]\label{dir:lem:pairrepair}
Suppose finitely many quota profiles each have an attained subset of the same compatible coset, with
missing fractions $\delta_\ell$ of sum $\Delta$. If $b_R\le\beta$, $\binom s2\omega_R\le1/2$ and
$\Delta\le\beta^2$, there is a block of $s$ distinct fresh labels after whose addition the missing fractions at
the incremented quotas have sum at most $4\beta^2\Delta$.
\end{lemma}

\begin{proof}
For independent draws $A_j,B_j$ from the residual law and $\eta\in\widehat{\mathcal D}$ with lift
$(\chi_j)$, $\dirE\eta(V)=\prod_{j<r}\abs{\dirE(\widetilde\chi_j\overline{\widetilde\chi_r})}^2$, which lies in
$[0,\beta^2]$ for $\eta\ne1$, since then some ratio is nonprincipal. For a missing set $U$ of density
$\delta$, translated to $\mathcal D$, orthogonality gives
$\dirE\abs{U\cap VU}/d=\sum_\eta\abs{\widehat{\mathbf 1}_U(\eta)}^2\dirE\eta(V)\le\delta^2+\beta^2(\delta-\delta^2)$.
The new attained set contains both translates of the old one, so after the baseline translation the new
missing set lies in $U\cap VU$. The $s$ draws are distinct with probability at least $1/2$, and
conditioning costs at most a factor two, so some distinct block gives
$\Delta_{\mathrm{new}}\le2(\Delta^2+\beta^2\Delta)\le4\beta^2\Delta$ for the whole family at once.
\end{proof}

\begin{proposition}[Finite completion criterion; \stU]\label{dir:prop:completion}
Let $n=m+st$ with $m,t\ge1$ and $0<\beta\le1/4$. Suppose the right side of \eqref{dir:eq:residual} for
$b_R$ is at most $\beta$, $n\omega\le1/2$, $\binom s2\omega/(1-n\omega)\le1/2$, the conditional bound of
\cref{dir:lem:collision} for a family $\mathcal L$ is at most $\beta^2$, and
$t\ge(\log d+1)/\log(1/(4\beta^2))$. Then some $n$-element pool of distinct labels is universal at every
quota $\ell+t\mathbf a$, $\ell\in\mathcal L$.
\end{proposition}

\begin{proof}
Choose $m$ distinct labels with total missing fraction at most $\beta^2$ by \cref{dir:lem:collision} and
apply \cref{dir:lem:pairrepair} $t$ times with fresh labels; the hypotheses persist. The final total is
at most $\beta^2(4\beta^2)^t<1/d$, and each missing fraction is a multiple of $1/d$.
\end{proof}

\begin{lemma}[A small-bias dyadic prime measure; \stG]\label{dir:lem:dyadic}
Fix $A>0$ and let $y=B\ell_M^2$ with $1\le B\le\ell_M^A$. For large $M$ the primes in $[y,2y]$ coprime to
$M$ carry a probability measure with $b\le KB^{-1/2}$ and $\omega\le K_A\log\ell_M/(B\ell_M^2)$, with $K$
absolute and $K_A$ depending only on $A$.
\end{lemma}

\begin{proof}
Give each such prime mass proportional to $w_y(p)\log p$, where
$w_y(v)=\frac2y\bigl((2y-v)_+-2(3y/2-v)_++(y-v)_+\bigr)$ is the triangle of height one on $[y,2y]$; the
normalizer is $y/2+o(y)$. The once-integrated estimate in the proof of \cref{dir:lem:bias}, which comes
from \cref{rec:lem:triangle}, at the three points $2y,3y/2,y$ gives $\abs{\sum_vw_y(v)\Lambda(v)\chi(v)}\ll\sqrt y\,\ell_M$ for $\chi\ne1$;
higher prime powers cost $O(\sqrt y\log^2(2y))=o(\sqrt y\,\ell_M)$ uniformly for $B\le\ell_M^A$. Each mass is
at most $3\log(2y)/y$.
\end{proof}

\begin{theorem}[Uniform entropy sufficiency; \stG]\label{dir:thm:entropyuniform}
Fix $r\ge2$, $A>0$, $0<a_0<1/r$ and $0<c_0<c_1$. There are $C_2,B_0,M_0$ such that for $M\ge M_0$,
$B_0\le B\le\ell_M^A$ and every integer $c_0\ell_M\le n\le c_1\ell_M$, one set of $n$ distinct primes coprime to
$M$ in $[B\ell_M^2,2B\ell_M^2]$ is universal at every quota vector $k$ with $\sum_jk_j=n$, $k_j\ge a_0n$ and
\begin{equation}\label{dir:eq:entropyslack}
 n\mathsf H_\pi(k/n)\ge\log N_\pi+C_2\frac{\ell_M}{\log B}\qquad(\abs\pi\ge2).
\end{equation}
\end{theorem}

\begin{proof}
Use the measure of \cref{dir:lem:dyadic}, $\beta=4\max(1,K)/\sqrt B$, and enlarge $B_0$ so that
$\beta\le1/4$ and $\log(1/(4\beta^2))\ge\frac12\log B$. Take $t=\ceil{(\log d+1)/\log(1/(4\beta^2))}$,
which is $O_r(\ell_M/\log B)$, and $m=n-st$; enlarging $B_0$ gives $st\le a_0n/2$. For each eligible $k$ use
the base quota $\ell=k-t\mathbf a$; along the segment from $k$ to $\ell$ every coordinate stays at least
$a_0n/2$, so the homogeneous entropy $F_\pi(u)=U\log U-\sum_Ju_J\log u_J$, whose partial derivatives are
$\log(U/u_J)\le\log(2/a_0)$ there, satisfies $F_\pi(\ell)\ge F_\pi(k)-st\log(2/a_0)$. With
$\log\binom m{(\ell_J)}\ge F_\pi(\ell)-O_r(\log(n+1))$ and \eqref{dir:eq:entropyslack}, a large $C_2$ gives
$\log\binom m{(\ell_J)}-\log N_\pi\ge3\ell_M/\log B$ for all eligible $k$ and $\pi$. There are at most
$(n+1)^r$ base quotas. Now $n\omega=o(1)$, $m^2\omega=O(\log\ell_M/B)$, the hypotheses on $b_R$ and $\omega_R$
hold, and after enlarging $B_0$ so that $(r-1)c_1K/\sqrt B\le1/\log B$ one has
$m\log(1+(\abs\pi-1)b)\le\ell_M/\log B$; so the conditional collision bound is at most
$\exp(-2\ell_M/\log B+O(\log\ell_M))\le\beta^2$. \Cref{dir:prop:completion} gives the set.
\end{proof}

\begin{proof}[Proof of the upper bound in \cref{dir:thm:entropy}]
Choose $a_0<\min_j\alpha_j/2$ and $c_0,c_1$ with $2c_0\ell_M\le T\le c_1\ell_M/2$, let
$h_*=\min_\pi\mathsf H_\pi(\alpha)>0$ and $C>(C_2+1)/h_*$, and enlarge $B_0$ so that $C/\log B_0<c_1/4$. For
$n$ as in the theorem every rounded profile has $k_j\ge a_0n$ and
$n\mathsf H_\pi(k/n)=n\mathsf H_\pi(\alpha)+O_{r,\alpha}(1)\ge\log N_\pi+Ch_*\ell_M/\log B-O(1)$, which is
\eqref{dir:eq:entropyslack}; apply \cref{dir:thm:entropyuniform}.
\end{proof}

The argument is existential, with an exhaustive finite procedure; support can be computed by dynamic
programming on quota counts and residue vectors. Exact computations of that kind find all $144$, $96$
and $2304$ compatible targets for the moduli $(13,13,13)$, $(12,18,30)$, $(15,21,35)$ with quotas
$(3,3,3)$, $(3,3,3)$, $(4,4,4)$.
\dirfiles{\artifact{entropy_optimal_prime_factorizations_2026-11-27/experiment_results.json}}

\paragraph{Smaller retained densities.}
At any positive retained density there is still an exact support theorem, with the extra obstruction
living in a quotient of bounded index, and it holds for a number of factors growing as fast as in
\cref{dir:thm:factorization}. This is the factorization analogue of \cref{dir:thm:quotient}, for a
selected list of primes instead of all primes up to $C(\log q)^2$. Occupied quotient classes are shown
not to be rare by an integer-table repair.

\begin{lemma}[Table repair; \stU]\label{dir:lem:repair}
Let $m$ labelled items carry colours in a finite abelian group $\Gamma$ of order $d$ and exponent
$e$, and let $Y$ and $X'$ be two allocations of the items into $k$ labelled blocks with the same block
sizes. Then $X'$ is within $6edk$ changed labels of an allocation with the same block sizes in which
every block has the same product of colours as in $Y$. Consequently, if the allocations with given
block sizes are equally likely, every attained vector of block colour products has probability at
least $(1+mk)^{-6edk}$.
\end{lemma}
\begin{proof}
Record an allocation by its $k\times d$ table of colour counts per block; $X'$ and $Y$ have equal row
and column sums. Let $\sigma_{ij}\in\set{0,\dots,e-1}$ be the residue of $Y_{ij}$ modulo $e$, take the
margins of $(Y-\sigma)/e$ as target margins, and start from
$Z_{ij}=\max\set{0,\floor{(X'_{ij}-\sigma_{ij})/e}}$, so that $\norm{\sigma+eZ-X'}_1\le ekd$ and the
total absolute row and column discrepancies of $Z$ from the target margins are at most $kd$ each.
Delete units from overfull rows and then from overfull columns, at cost at most $kd$ per stage, since
row deletions only decrease column excesses; the remaining row deficit equals the column deficit and
is at most $3kd$, and adding units at intersections of deficient rows and columns fills it. With the
repaired table $Z_*$, the table $X''=\sigma+eZ_*$ has the margins of $X'$ and entries congruent to those
of $Y$ modulo $e$, hence the colour products of $Y$, and $\norm{X''-X'}_1\le6ekd$. Moving individual
items of each colour from surplus to deficient blocks realizes $X''$ with at most $6ekd$ changed
labels. An allocation has at most $(1+mk)^{6edk}$ allocations within that many changes (pad a list of
edits with empty moves), and the neighbourhoods of the allocations with the target products cover all
allocations.
\end{proof}

\begin{remark}[A cycle-compression bound]\label{dir:rem:cycle}
The number of changed labels in \cref{dir:lem:repair} can also be bounded by $d^2(e-1)k$, which is zero
for $d=1$ and smaller than $6edk$ when $d(e-1)<6e$; the mass bound becomes $(1+mk)^{-d^2(e-1)k}$. Let
$A_0,B_0$ be the $d\times k$ count tables of $X'$ and $Y$. The integer matrix $B_0-A_0$ has zero row and
column sums, so it is a balanced circulation on the bipartite graph of colours and blocks (a positive
entry an arc from colour to block, a negative one an arc back). Peeling off simple directed cycles writes
it as a sum of at most $dk$ cycles with positive integer coefficients $w_i$, each cycle changing at most
$d$ placements per unit coefficient. Adding the cycles to $A_0$ with coefficients $w_i\bmod e$ gives a
table $C_0$ with the margins of $A_0$, entries between those of $A_0$ and $B_0$, and $C_0\equiv B_0\pmod e$,
hence the colour products of $Y$; moving items of each colour from surplus to deficit cells of
$A_0-C_0$ changes at most $d\sum_i(w_i\bmod e)\le d^2(e-1)k$ labels. A cycle of length six in a fixture
with $900\,003$ original moves is reduced to three moves.
\dirfiles{\artifact{grh_growing_factorizations/checks.txt}}
\end{remark}

\begin{theorem}[Factorizations at any positive density; \stG]\label{dir:thm:minority}
Fix $0<a\le1$ and $A>0$, and put $D=\floor{1/a}$. There are $B,K_1,K_2,c>0$, depending only on $a$
and $A$, such that, with $L=\log Q$, for every sufficiently large $Q$ and every integer
$2\le R\le cL/\log L$ there is one list $S$ of between $K_1RL$ and $2K_1RL+2$ distinct primes in
$[x,2x]$, where $x\asymp_aL^2$ and $x\le K_2L^2$, with the following property. Let $2\le r\le R$ and
$1\le q_1,\dots,q_r\le Q$, put $M=\lcm_jq_j$ and $g=\gcd_jq_j$, let $U$ be the primes of $S$ coprime to
$M$, and let $T\subseteq U$ with $m=\abs T\ge a\abs U$. Let $\mu$ be the factorization law of $T$ in
$H=\prod_j\dirU(q_j)$ for quotas $k_j$ with $k_*=\min_jk_j\ge BL$; balanced such quotas exist for
every such $T$. There is a character subgroup $E\subseteq\widehat H$, independent of the quotas,
containing the diagonal group $\Delta_g\simeq\widehat{\dirU(g)}$, with $[E:\Delta_g]\le D^{r-1}$. With the
evaluation map $\varpi\colon H\to\widehat E$ and $\nu=\varpi_*\mu$, a vector $u\in H$ occurs exactly
when $\nu(\varpi(u))>0$, and then
\[
 \mu(u)=\frac{\abs E}{\abs H}\,\nu(\varpi(u))(1+\theta_u),\qquad\abs{\theta_u}\le Q^{-A}.
\]
At most $[E:\Delta_g]$ projections are occupied, so at least a fraction $[E:\Delta_g]^{-1}\ge D^{1-r}$ of
the compatible tuples occur, and the index bound and this fraction are sharp. For $a>1/2$ the only
constraint is $\prod_ju_j\equiv P_T\pmod g$.
\end{theorem}

\begin{proof}
\emph{The list.} Choose $\tau,\eta>0$ with $\eta<1$, $1-D^2\tau>0$ and
$(D+1)a(1-D^2\tau)^2>1+D\eta$, and apply \cref{dir:lem:oneset} with $\delta=\eta/4$ and $Q^{2D}$ in
place of $Q$: for every admissible length $N_*$ there is a list of between $N_*$ and $2N_*$ primes in
$[x,2x]$, $x=K_\delta\log^2(2Q^{2D})$, on which every nonprincipal character of modulus at most
$Q^{4D}$ has mean of modulus at most $\delta$. Take $B\ge\max\set{1,(2/\tau)(A+36D^2+13)}$, then $K_1$
so large that $N_*=\ceil{K_1RL}$ is admissible and $aK_1/2\ge2(B+1)$, then $c\le1$ so small that
$N_*$ stays admissible for $R\le cL/\log L$. At most $rL/\log x=o(\abs S)$ primes of the list divide
$M$, so every such character has mean of modulus at most $\eta$ on $U$, $m\ge aK_1RL/2$, balanced
quotas exceed $BL$, and for large $Q$, since $r\le L/\log L$,
\begin{equation}\label{dir:eq:mixedsize}
 r\log(m+1)\le3L .
\end{equation}

\emph{The subgroup.} Inflate block characters to $\dirU(M)$, writing $\widetilde\psi_j$, and let $J$ be
generated by the ratios $\widetilde\psi_i\overline{\widetilde\psi_j}$ of bias greater than $1-\tau$ on $T$.
The quotient of two distinct words of length at most $D$ in these generators involves at most $4D$
block characters, so it is induced by a nonprincipal character of modulus at most $Q^{4D}$, and the
induced and inflated values agree on $U$; this bounds the Gram entries by $\eta$ even though $M$ may
exceed every fixed power of $Q$. Each word has bias at least $1-D^2\tau$ on $T$, and the Gram and
word-ball argument of \cref{dir:thm:quotient} with $\abs T/\abs U\ge a$ gives $d=\abs J\le D$. Let $E$
consist of the tuples whose inflated characters lie in one coset of $J$. Taking ratios with
$\widetilde\psi_1$ gives a homomorphism $E\to J^{r-1}$ whose kernel consists of equal inflated characters;
a common character has conductor dividing every $q_j$, so the kernel is $\Delta_g$, and
$[E:\Delta_g]\le D^{r-1}$.

\emph{Fourier decay.} For a tuple outside $E$ let $t\ge1$ be the number of labels outside a coset of
largest combined quota. Labels in different cosets have ratios of bias at most $1-\tau$ on $T$; in
\cref{dir:lem:partition} keep the $r$ labels for the conditioning factor and discard the pairs inside
a coset, so that the coefficient is at most $(m+1)^re^{-\tau tk_*/2}$. For each $t$ there are at most
$rQ\binom rtd^rQ^t$ such tuples: a character in the dominant coset from one block, the outside labels
and their characters, and at most $d$ choices for every inside label. As $rQe^{-\tau k_*/2}\le1/2$ for
large $Q$, \eqref{dir:eq:mixedsize}, $d\le m+1$ and $r\le Q$ give
\[
 \sum_{\boldsymbol\psi\notin E}\abs{\widehat\mu(\boldsymbol\psi)}\le(m+1)^rrQd^r\bigl[(1+Qe^{-\tau k_*/2})^r-1\bigr]
 \le2Q^{10-\tau B/2},
\]
while Fourier inversion over $E$ contributes exactly $\abs E\nu(\varpi(u))$ to $\abs H\mu(u)$.

\emph{Occupied projections.} Colour each prime by its value in $\widehat J$. For a tuple in $E$ with
$\delta_j=\widetilde\psi_j/\widetilde\psi_1\in J$, its value on factors $F_j$ is
$\widetilde\psi_1(P_T)\prod_j\delta_j(F_j)$, so the $r$ colour products of a witness determine its
$E$-projection. By \cref{dir:lem:repair} with $k=r$ and $e\le d\le D$, every occupied projection has
mass at least $(1+mr)^{-6D^2r}\ge Q^{-36D^2}$ by \eqref{dir:eq:mixedsize}, which exceeds $Q^A$ times the
Fourier error by the choice of $B$. All occupied projections restrict to the same character of
$\Delta_g$, fixed by $P_T$, which bounds their number and gives the case $D=1$.

\emph{Sharpness.} Let all moduli equal a prime $\ell\equiv1\pmod D$ and retain, as in
\cref{dir:rem:thresholds}, the most populous coset in $U$ of the subgroup of index $D$ of $\dirU(\ell)$, at
least $\abs U/D\ge a\abs U$ primes. Every
character of the cyclic quotient of order $D$ has bias one on this set, so $J$ is that character
group, $E/\Delta_g\to J^{r-1}$ is onto and $[E:\Delta_g]=D^{r-1}$; each quota fixes its label's quotient
residue, and compatibility imposes one relation among the $r$ residues, so the support is a $D^{1-r}$
fraction of the compatible fibre.
\end{proof}

The growing range is sharp in order for this size of primes and the balanced-quota guarantee: even
without deletions, $r$ blocks of size $BL$ use $rBL$ distinct primes, while $[1,2K_2L^2]$ contains only
$O(L^2/\log L)$ primes. For $a>1/2$ the theorem contains \cref{dir:thm:factorization} up to the
constants, which are explicit there.

The projected law $\nu$ can be computed before any full tuple is tested. Start from the
allocation placing every prime in label $1$, with residue vector $(P_T,1,\dots,1)$; for each
prime $p$ and label $j$ let $v_{p,j}$ be the projected change caused by moving $p$ to label $j$,
which lies in the kernel of $\widehat E\to\widehat{\Delta_g}$, of order $s$. Processing one prime
at a time and storing integer multiplicities indexed by the current quota counts of the first
$r-1$ labels and the current class gives at most $s\prod_{j<r}(k_j+1)$ states per layer, and a
finite procedure for the support criterion.

\begin{example}\label{dir:ex:factorizations}
\begin{enumerate}[label=\textup{(\arabic*)}]
\item Take the fifteen primes from $101$ to $173$ and three factors of five primes each, so that
there are $\binom{15}{5,5,5}=756\,756$ ordered assignments. For the moduli $(12,15,20)$, with
$g=1$, all $256$ unit triples occur, with $2712$ to $3196$ representations and maximum relative
deviation $8.257\%$ from the uniform count; for $(6,9,12)$, with $g=3$, exactly the $24$
compatible triples of the $48$ occur, with $30\,426$ to $32\,502$ representations and maximum
deviation $3.506\%$.
\item The projected law can be nonuniform on its occupied classes. Give three factors the
modulus $5$ and quotas $(2,3,4)$, and use the primes $11,19,29,31,41,59,61,71,13$, of which only
$13$ is a nonresidue modulo $5$. The total product is $2\bmod5$, so the four quadratic-sign
triples with an odd number of minus signs are compatible; exactly one factor contains $13$, so
$(-,+,+),(+,-,+),(+,+,-)$ have probabilities $2/9,3/9,4/9$ and $(-,-,-)$ has probability zero.
Enumeration of all $1260$ factorizations finds all four compatible residue triples in each
occupied sign class, and the dynamic program above gives the same sign law.
\end{enumerate}
\dirfiles{\artifact{squarefree_factorizations/check_output.txt},
\artifact{squarefree_factorizations/check_quotient_output.txt},
\artifact{squarefree_factorizations/test_lower_output.txt}}
\end{example}

\subsection{Disjoint products of small primes in one modulus}\label{dir:sec:disjoint}
\sources{disjoint-prime-products; grh_quadratic_fibre; grh_quarter_density}

How can a few extra primes change which products are possible? If every retained prime is a
quadratic residue modulo an odd prime except one nonresidue, that prime can go into either of two
products or stay unused, so each product separately can have either sign but the two can never both
be negative; a second nonresidue removes this joint obstruction, although the four sign patterns
still have different probabilities. This subsection describes every attainable tuple of several
disjoint products at retained density at least one quarter, with sharp coverage thresholds, and above
quarter density an exact criterion for uniform counts. At any fixed density the support is decided by
a zero-sum completion problem in a bounded quotient, solved in polynomial time once the quotient is
known; at density one quarter the two quotients of order four need, respectively, three quadratic
projections and four linear inequalities with an additional parity cut, and for every odd number of
products a single joint pattern can be missing although every proper marginal is full. It also answers the question of \cref{dir:sec:products} whether the distance from a
quadratic fibre can replace a fixed density margin above one half: the right quantity is the number of
minority primes times the relative quota, and for one product the counts on both sign classes follow a
binomial law with relative error $O(1/h)$, which is sharp (\cref{dir:thm:fibre}).

Keep the notation of \cref{dir:sec:products}: $q\ge3$, $G=(\Z/q\Z)^\times$, $L=\log q$, $y=CL^2$,
$\mathcal P$ the primes up to $y$ not dividing $q$, $N=\abs{\mathcal P}$, $S\subseteq\mathcal P$ with
$m=\abs S\ge aN$. Fix $r\ge1$ and choose a uniform ordered partition $S=S_0\sqcup S_1\sqcup\dots\sqcup S_r$
with prescribed sizes $n_0,\dots,n_r$, $\sum_in_i=m$; block $0$ is a reserve. Put
\[
 X=(X_1,\dots,X_r),\qquad X_j=\prod_{p\in S_j}p\bmod q,\qquad h=\min_{0\le i\le r}n_i,\qquad t=h/m .
\]
By unique factorization the $m!/\prod_in_i!$ equally likely partitions are the ordered tuples of
pairwise coprime squarefree integers with these numbers of prime factors from $S$. For $r=1$ this is
the product $X_k$ of \cref{dir:sec:products} with $k=n_1$. The constants $B,C$ depend only on $a,r$
(and on $A$ where it occurs), and the conclusions hold for all large $q$ and all admissible $S$ and
quotas with $h\ge BL$.

\begin{theorem}[A growing number of products at any positive density; \stG]\label{dir:thm:disjoint-growing}
Fix $0<a\le1$ and $A>0$, and put $D=\floor{1/a}$. There are $C=C(a)$ and $B=B(a,A)$ such that the
following holds for all large $q$, every $S$ with $m\ge aN$ and every $r\ge1$ with $h\ge BL$. The
subgroup $H_\tau$ of \cref{dir:thm:quotient} has order $d\le D$; with the evaluation quotient
$\pi\colon G\to\widehat{H_\tau}$ and $\nu$ the joint law of the projected products, the support of $X$ is
exactly the inverse image of $\supp\nu$, and at every occupied $g$
\[
 \dirP(X=g)=\frac{d^r}{\abs G^r}\,\nu(\pi^r(g))(1+\theta_g),\qquad\abs{\theta_g}\le q^{-A}.
\]
Balanced quotas allow $r\asymp_{a,A}L/\log L$ at every permitted retained set, and this order is
maximal for quotas of size $BL$.
\end{theorem}

\begin{proof}
The Gram argument of \cref{dir:thm:quotient} works with any fixed $D$ and gives $\tau=\tau(a)$, then
$C$, with $d\le D$. Put $k=r+1$; since $kh\le m\le N$ and $N\sim CL^2/\log(CL^2)$, there is $K=3(C+1)$
with $k\log(m+1)\le KL$. The Fourier and repair steps of the proof of \cref{dir:thm:minority} apply
with $k$ labels, the reserve label carrying the trivial character, $J=H_\tau$, and $q$ in place of
$Q$: a tuple outside $H_\tau^r$ has $t_*\ge1$ labels outside a coset of $H_\tau$ of largest weight,
because the reserve character is trivial, its coefficient is at most $(m+1)^ke^{-\tau t_*h/2}$, there
are at most $q\binom k{t_*}d^kq^{t_*}$ such tuples, and for $B\ge8/\tau$ the total error is at most
$2q^{2K+3-\tau B/2}$; orthogonality on $H_\tau^r$ gives the main term $d^r\nu(\pi^r(g))$ on the
normalized scale. Colouring primes by their values in $\widehat{H_\tau}$, \cref{dir:lem:repair} gives
$\nu(\pi^r(g))\ge(1+mk)^{-6D^2k}\ge q^{-12D^2K}$ at every occupied point, and
$B\ge(2/\tau)(A+12D^2K+2K+6)$ gives the relative error. Finally $kh\le m$ gives $r=O_{a,A}(L/\log L)$,
balanced quotas with $k=\floor{aN/(2BL)}$ attain this order, and $r$ blocks of size $BL$ need $rBL$
primes, of which there are only $O(L^2/\log L)$.
\end{proof}

For $r=1$ this sharpens the relative asymptotic of \cref{dir:thm:quotient} to $1+O(q^{-A})$ when $B$
may depend on $A$.

\paragraph{Zero-sum completion in the quotient.}
By \cref{dir:thm:disjoint-growing} the support of $X$ is that of the projected allocation problem in
$\Gamma=\widehat{H_\tau}$, a group of order $d\le\floor{1/a}$, written additively. This is a finite
problem about coloured items. Let $m$ labelled items carry colours in a finite abelian group $\Gamma$ of
order $d$, to be partitioned into a reserve block $0$ and recorded blocks $1,\dots,r$ with quotas
$n_0,\dots,n_r$, $\sum_jn_j=m$. Translating all colours by a fixed $b\in\Gamma$ makes $b$ the baseline
$0$; the corrected target of block $j$ is its original target minus $n_jb$, and $c_a$ is the number of
items of translated colour $a$.

\begin{lemma}[Short witnesses; \stU]\label{dir:lem:shortwitness}
Every multiset of nonzero colours summing to $z$ contains a submultiset of length at most $d-1$
summing to $z$; for $z=0$ the empty multiset suffices.
\end{lemma}

\begin{proof}
Delete nonempty zero-sum submultisets while one exists; this preserves the sum. If the remainder had
length at least $d$, two of its first $d+1$ partial sums, the empty one included, would coincide, and
the terms between them would form a nonempty zero-sum subsequence.
\end{proof}

\begin{theorem}[Quota-free reduction; \stU]\label{dir:thm:packing}
Assume
\[
 n_j\ge d-1\quad(1\le j\le r),\qquad n_0\ge M_\Gamma:=\sum_{a\ne0}(\ord a-1),\qquad c_0\ge(d-1)(r+1).
\]
Then the following are equivalent:
\begin{enumerate}[label=\textup{(\alph*)}]
\item some partition of all items with these quotas has recorded colour sums equal to the corrected
targets $z_1,\dots,z_r$;
\item there are pairwise disjoint multisets of nonzero-colour items, the $j$th summing to $z_j$, with no
quotas imposed and unused items allowed;
\item the multisets of \textup{(b)} can be chosen of length at most $d-1$.
\end{enumerate}
The passage from \textup{(c)} to \textup{(a)} is constructive; no search over complete partitions is needed.
\end{theorem}

\begin{proof}
Deleting the zero-colour items of each recorded block of (a) gives (b), and \cref{dir:lem:shortwitness}
gives (c). Given (c), place each witness in its recorded block, which it fits since $n_j\ge d-1$. For each
nonzero colour $a$ split its unused items into packets of $\ord a$ items, each of sum zero, and one
leftover group of fewer than $\ord a$ items; all leftover groups fit in the reserve since $n_0\ge M_\Gamma$.
Distribute the full packets, in any order, to any block with enough remaining capacity. If a packet of
size $s\le d$ does not fit, every block has remaining capacity at most $s-1$, so the total remaining
capacity is at most $(s-1)(r+1)\le(d-1)(r+1)$; but it equals the number of unplaced items, which include all
$c_0$ zero-colour items and the $s>0$ items of the pending packet, contradicting $c_0\ge(d-1)(r+1)$. Fill
the remaining places with zero-colour items; no recorded sum changes. For $d=1$ there are no nonzero
colours.
\end{proof}

\begin{lemma}[Direct matches first; \stU]\label{dir:lem:direct}
For $a\ne0$ let $t_a$ be the number of recorded blocks with corrected target $a$. If \textup{(b)} of
\cref{dir:thm:packing} is feasible, it is feasible with $\min(c_a,t_a)$ singleton witnesses of colour $a$
for every $a\ne0$. The remaining witnesses then serve only targets of colours with $t_a>c_a$ and use only
items of colours with $c_a>t_a$.
\end{lemma}

\begin{proof}
Take the witnesses of zero targets empty and, among all feasible collections, one with the largest number
of singleton witnesses. Suppose fewer than $\min(c_a,t_a)$ singletons have colour $a$. Then some target $a$
has a non-singleton witness $W$, and some item $x$ of colour $a$ is not a singleton witness. If $x$ is
unused or lies in $W$, replace $W$ by $\set x$; if $x$ lies in another non-singleton witness $V$, replace $W$
by $\set x$ and $V$ by $(V\setminus\set x)\cup W$, which has the sum of $V$ because $W$ sums to $a$. The
witnesses stay disjoint, no singleton is lost and one is gained, a contradiction. After all direct
matches, a colour with $c_a\le t_a$ has no items left and one with $c_a\ge t_a$ no targets left.
\end{proof}

\begin{corollary}[Exact support at any positive density; \stG]\label{dir:cor:anydensity}
In the setting of \cref{dir:thm:disjoint-growing}, colour each $p\in S$ by $\pi(p)\in\Gamma$, let $b$ be a
most frequent colour, and for $g\in G^r$ put $z_j=\pi(g_j)-n_jb$. For all large $q$, the target $g$ lies in
the support of $X$ if and only if condition \textup{(b)} of \cref{dir:thm:packing} holds for
$(z_1,\dots,z_r)$, and then $\dirP(X=g)$ is given by \cref{dir:thm:disjoint-growing}. For fixed $d$, given
the colour counts and the targets, this is decided in $O_d\bigl(r(1+(d-1)r)^{d-1}\bigr)$ operations on
small integers, and a witness allocation is reconstructed from stored predecessors and the proof of
\cref{dir:thm:packing}.
\end{corollary}

\begin{proof}
A most frequent colour has $c_0\ge m/d\ge(r+1)h/d$, every quota is at least $h\ge BL$, and
$M_\Gamma\le d(d-1)$ with $d\le\floor{1/a}$, so the hypotheses of \cref{dir:thm:packing} hold for large $q$.
The support of $X$ is the inverse image of $\supp\nu$, and $\pi^r(g)\in\supp\nu$ is condition (a) for the
colouring by $\pi$. For the algorithm list, for each $z\in\Gamma$, the vectors $e\in\Z_{\ge0}^{\Gamma\setminus\set0}$
with $\sum_ae_a\le d-1$ and $\sum_ae_aa=z$, only $e=0$ for $z=0$; cap each nonzero supply at $(d-1)r$, the
most that short witnesses can use; and block by block add the candidate vectors to the reachable vectors of
used counts, discarding those above the capped supply. There are at most $(1+(d-1)r)^{d-1}$ states, and the
final set is nonempty exactly when (c) holds.
\end{proof}

The procedure presupposes the quotient; it does not compute $H_\tau$ from the retained primes.

\paragraph{Quotients of order at most four.}
For cyclic $\Gamma$ label the colours by their natural residues, and in $C_2\times C_2$ label the nonzero
elements $1,2,3$ with $1+2=3$. Write $c_1,c_2,c_3$ for the supplies of the nonzero translated colours and
$\ell_1,\ell_2,\ell_3$ for the numbers of recorded blocks with corrected targets $1,2,3$ (only the indices
that occur in $\Gamma$).

\begin{theorem}[Quota-free criteria in groups of order at most four; \stU]\label{dir:thm:smallquotients}
Condition \textup{(b)} of \cref{dir:thm:packing} holds if and only if:
\begin{enumerate}[label=\textup{(\alph*)}]
\item always for the trivial group; $c_1\ge\ell_1$ for $C_2$; and for $C_3$
\[
 2c_1+c_2\ge2\ell_1+\ell_2,\qquad c_1+2c_2\ge\ell_1+2\ell_2 ;
\]
\item for $C_2\times C_2$,
\begin{equation}\label{dir:eq:klein}
 c_1+c_2\ge\ell_1+\ell_2,\qquad c_1+c_3\ge\ell_1+\ell_3,\qquad c_2+c_3\ge\ell_2+\ell_3 ;
\end{equation}
\item for $C_4$, the four inequalities
\begin{equation}\label{dir:eq:cyclic4}
\begin{gathered}
 c_1+c_3\ge\ell_1+\ell_3,\qquad 3c_1+2c_2+c_3\ge3\ell_1+2\ell_2+\ell_3,\\
 c_1+2c_2+3c_3\ge\ell_1+2\ell_2+3\ell_3,\qquad c_1+2c_2+c_3\ge\ell_1+2\ell_2+\ell_3,
\end{gathered}
\end{equation}
and, only when $\ell_2>c_2$, $\ell_1\le c_1$ and $\ell_3\le c_3$, the parity cut
\begin{equation}\label{dir:eq:paritycut}
 \floor{(c_1-\ell_1)/2}+\floor{(c_3-\ell_3)/2}\ge\ell_2-c_2 .
\end{equation}
\end{enumerate}
\end{theorem}

\begin{proof}
By \cref{dir:lem:shortwitness} the witnesses may be taken zero-sum free, and by \cref{dir:lem:direct} with
all direct matches made.

(a) For $C_2$ each nonzero target needs a nonzero item. In $C_3$ a zero-sum-free witness of sum $1$ is
$\set1$ or $\set{2,2}$, and of sum $2$ is $\set2$ or $\set{1,1}$; with the weights $(2,1)$ on the colours
$(1,2)$ every witness of sum $1$ weighs at least two and every witness of sum $2$ at least one, which gives
the first inequality, and the weights $(1,2)$ give the second. Conversely both colours cannot be deficient,
by either inequality. If $c_1<\ell_1$, each of the $\ell_1-c_1$ unmatched targets $1$ needs two items of
colour $2$, and the first inequality is exactly $c_2-\ell_2\ge2(\ell_1-c_1)$; the case $c_2<\ell_2$ is
symmetric, and otherwise direct matching suffices.

(b) A zero-sum-free witness of a nonzero sum in $C_2\times C_2$ is the singleton of that colour or the pair of
the other two nonzero colours, since equal colours cancel and the three nonzero elements sum to zero. Each
pair sum, say $c_1+c_2$, counts the items with nonzero image under the homomorphism to $C_2$ with kernel
$\set{0,3}$, and each target $1$ or $2$ needs such an item; this gives necessity. Two deficient colours
violate the inequality of their pair. If, say, $c_1<\ell_1$, the first two inequalities give
$c_2-\ell_2\ge\ell_1-c_1$ and $c_3-\ell_3\ge\ell_1-c_1$, and one surplus $2$ with one surplus $3$ serves each
unmatched target $1$; the other cases are symmetric.

(c) As $1+3$ and $2+2$ are zero sums, the zero-sum-free multisets of nonzero sum in $C_4$ are
\begin{center}\small
\begin{tabular}{@{}cl@{}}
\toprule
target & minimal witnesses\\
\midrule
$1$ & $\set1,\ \set{2,3},\ \set{3,3,3}$\\
$2$ & $\set2,\ \set{1,1},\ \set{3,3}$\\
$3$ & $\set3,\ \set{1,2},\ \set{1,1,1}$\\
\bottomrule
\end{tabular}
\end{center}
\emph{Necessity.} For each of the cost vectors $(1,0,1)$, $(3,2,1)$, $(1,2,3)$, $(1,2,1)$ on the colours
$1,2,3$, every witness in the table costs at least the cost of its sum; summing over disjoint witnesses
gives \eqref{dir:eq:cyclic4}. In the active case of the parity cut, direct matching serves every odd target
and uses all items of colour $2$, and each of the $\ell_2-c_2$ remaining targets $2$ needs a pair of equal odd
colours from the surplus; there are at most the left side of \eqref{dir:eq:paritycut} such pairs.

\emph{Sufficiency.} After direct matching put $U=c_1-\ell_1$, $V=c_2-\ell_2$, $W=c_3-\ell_3$. If none is
negative, direct matching suffices. If only $U<0$, put $D=-U$; the first two inequalities give $W\ge D$ and
$2V+W\ge3D$, so $x=\min(V,D)$ unmatched targets $1$ can be served by $\set{2,3}$ and the other $D-x$ by
$\set{3,3,3}$, using $3D-2x\le W$ items of colour $3$; only $W<0$ is symmetric. If only $V<0$, the parity cut
supplies $-V$ disjoint pairs of equal odd colours. Negative $U$ and $W$ together violate the first
inequality. If $V,W<0$, the third inequality gives $U\ge2(-V)+3(-W)$, and the deficient targets $2$ and $3$ are
served by $\set{1,1}$ and $\set{1,1,1}$; if $U,V<0$, the second gives $W\ge3(-U)+2(-V)$, and $\set{3,3,3}$ and
$\set{3,3}$ serve. All allocations respect the supplies, and \cref{dir:thm:packing} completes them to the
quotas.
\end{proof}

The Klein case is detected by its three quadratic projections, whereas the parity cut in $C_4$ is a genuinely
integral condition. In its active case put $D=\ell_2-c_2$, $U=c_1-\ell_1\ge0$, $W=c_3-\ell_3\ge0$; the fourth
inequality of \eqref{dir:eq:cyclic4} is $U+W\ge2D$, and given it the cut fails exactly when $U,W$ are both odd
and $U+W=2D$. So all the linear tests can pass at an unreachable target, already for $(c_1,c_2,c_3)=(1,0,1)$
and $(\ell_1,\ell_2,\ell_3)=(0,1,0)$.

\begin{theorem}[Coverage capacities; \stU]\label{dir:thm:capacity}
All corrected vectors in $\Gamma^r$ satisfy \textup{(b)} of \cref{dir:thm:packing} if and only if $r\le R_\Gamma$,
where
\begin{center}\small
\begin{tabular}{@{}ll@{}}
\toprule
$\Gamma$ & $R_\Gamma$\\
\midrule
$C_2$ & $c_1$\\
$C_3$ & $\floor{\min(2c_1+c_2,\,c_1+2c_2)/2}$\\
$C_2\times C_2$ & $\min(c_1+c_2,\,c_1+c_3,\,c_2+c_3)$\\
$C_4$ & $\min\bigl\{c_1+c_3,\ \floor{(3c_1+2c_2+c_3)/3},\ \floor{(c_1+2c_2+3c_3)/3},\ c_2+\floor{c_1/2}+\floor{c_3/2}\bigr\}$\\
\bottomrule
\end{tabular}
\end{center}
and $R_\Gamma=\infty$ for the trivial group. The constant corrected target vectors alone decide this.
\end{theorem}

\begin{proof}
For $C_2$, $C_3$ and $C_2\times C_2$ maximize the right sides of the linear tests over histograms with total at
most $r$; a maximizer is a constant nonzero vector, and the bounds are the ones listed. For $C_4$ the constant
vectors with entries $1$ and $3$ force $c_1+c_3\ge r$, $3c_1+2c_2+c_3\ge3r$ and $c_1+2c_2+3c_3\ge3r$, and the
constant vector with entries $2$ forces $c_2+\floor{c_1/2}+\floor{c_3/2}\ge r$ by the parity cut. Conversely
these four bounds give the first three inequalities of \eqref{dir:eq:cyclic4} for every histogram, since
$\ell_1+\ell_2+\ell_3\le r$, and the fourth since $c_1+2c_2+c_3\ge2r\ge\ell_1+2\ell_2+\ell_3$; in the active case
of the cut,
\[
 \floor{(c_1-\ell_1)/2}+\floor{(c_3-\ell_3)/2}\ge\floor{c_1/2}+\floor{c_3/2}-\ell_1-\ell_3\ge r-c_2-\ell_1-\ell_3\ge\ell_2-c_2 .
\]
\cref{dir:thm:smallquotients} completes the proof.
\end{proof}

At retained density $m\ge N/4$, for large $q$ and $h\ge BL$, \cref{dir:cor:anydensity} with $a=1/4$ and
\cref{dir:thm:smallquotients} decide the support: the quotient $\Gamma=\widehat{H_\tau}$ has order $d\le4$ and depends only on $S$, and with $b$
and $z_j$ as in \cref{dir:cor:anydensity} a target $g\in G^r$ occurs if and only if the histogram of
$(z_1,\dots,z_r)$ passes the test of \cref{dir:thm:smallquotients} for $\Gamma$. The same test can be run
without determining $\Gamma$.

\begin{theorem}[Exact support at quarter density; \stG]\label{dir:thm:quarter}
Fix $A>0$. There are $B,C>0$ such that for all large $q$, every $S\subseteq\mathcal P$ with $m\ge N/4$, every
$r\ge1$ and all quotas with $h\ge BL$ the following hold. A target $g\in G^r$ occurs if and only if, for
every character $\chi$ of exact order $2$, $3$ or $4$, the test of \cref{dir:thm:smallquotients} for the cyclic
group $\chi(G)$ passes for the colouring $p\mapsto\chi(p)$ of $S$, with a most frequent value $b_\chi$ of $\chi$ on
$S$ as baseline and corrected targets $\chi(g_j)b_\chi^{-n_j}$; at every occupied $g$ the relative count is that
of \cref{dir:thm:disjoint-growing}. Balanced quotas allow $r\asymp L/\log L$.
\end{theorem}

\begin{proof}
Necessity follows by projection: an allocation realizing $g$ maps under $\chi$ to an allocation of the
$\chi$-coloured primes with the same quotas, and the implication from (a) to (b) in \cref{dir:thm:packing}
needs no hypotheses. For sufficiency use the quotient test above: if $\Gamma$ is cyclic a generator of $H_\tau$ is among the tested
characters, its most frequent value is the image of a most frequent colour, and its test is the quotient
test. If $\Gamma=C_2\times C_2$, each inequality of \eqref{dir:eq:klein} is the test of one of the three
quadratic characters in $H_\tau$ with the image of $b$ as baseline. If that image is not a most frequent
value of the character, or is one of two tied values, at least $m/2\ge r$ primes are nonzero in that
projection and the inequality holds trivially; otherwise it is the tested condition. The growth of $r$ is
that of \cref{dir:thm:disjoint-growing}.
\end{proof}

The theorem uses only characters of order at most four, but enumerating them for composite $q$ is not
claimed to take time polynomial in $\log q$.

For a character $\chi$ of order $2$ or $3$ choose a most frequent value $b$ of $\chi$ on $S$, put
$s_\chi=\#\set{p\in S:\chi(p)\ne b}$, and for a target $g\in G^r$ call $w_j=\chi(g_j)b^{-n_j}$ its
corrected values; let $\omega=e^{2\pi i/3}$. Above quarter density the tests take the following closed form.

\begin{corollary}[Exact support above quarter density; \stG]\label{dir:thm:disjoint-support}
Let $1/4<a\le1$. A target $g\in G^r$ occurs if and only if, for every character of order
two, the number of $j$ with $w_j=-1$ is at most $s_\chi$, and for every character of order three,
with $u=\#\set{p\in S:\chi(p)=b\omega}$, $v=\#\set{p\in S:\chi(p)=b\omega^{-1}}$ and $i,j$ the numbers of
corrected entries equal to $\omega$ and $\omega^{-1}$,
\[
 2u+v\ge2i+j,\qquad u+2v\ge i+2j .
\]
In particular $X$ covers all of $G^r$ if and only if
\[
 s_\chi\ge r\quad(\operatorname{ord}\chi=2),\qquad u+v+\min(u,v)\ge2r\quad(\operatorname{ord}\chi=3).
\]
For $a>1/3$ the quadratic tests alone suffice, and ties for $b$ may be broken arbitrarily. These
tests and thresholds hold for fixed $r$ and throughout the growing range of \cref{dir:thm:disjoint-growing}.
\end{corollary}

\begin{proof}
For $a>1/4$ the Gram argument gives $d\le3$, and $d\le2$ for $a>1/3$, so $\Gamma$ is cyclic of order at most
three. The displayed tests are those of \cref{dir:thm:smallquotients}(a) for the colouring by $\chi$, and the
argument of \cref{dir:thm:quarter} (necessity by projection, sufficiency through a generator of $H_\tau$)
applies with \cref{dir:cor:anydensity} at density $a$. Since $\min(2u+v,u+2v)=u+v+\min(u,v)$, the coverage
thresholds are those of \cref{dir:thm:capacity}, again necessary for every character by projection and
sufficient for a generator.
\end{proof}

For $a>1/3$ and $s=\min_{\operatorname{ord}\chi=2}s_\chi$, a minimizing character generates $H_\tau$ when
$s<r$, since its bias is $1-2s/m>1-\tau$, so
\[
 \abs{\supp X}=\Bigl(\frac{\abs G}2\Bigr)^r\sum_{j=0}^{\min(r,s)}\binom rj ;
\]
for $s=r-1$ every proper marginal has full support while the joint law misses one sign class. For
$r=1$ the corollary is the criterion of \cref{dir:rem:thresholds}. In the same way, when $\Gamma$ has order
four, each corrected quotient vector is a translate of one quotient target and has $(\abs G/4)^r$ lifts, all
attained, so
\[
 \abs{\supp X}=\Bigl(\frac{\abs G}4\Bigr)^r\sum_{(\ell_1,\ell_2,\ell_3)\in\mathcal H_r}
 \frac{r!}{\ell_1!\,\ell_2!\,\ell_3!\,(r-\ell_1-\ell_2-\ell_3)!},
\]
with $\mathcal H_r$ the set of histograms with $\ell_1+\ell_2+\ell_3\le r$ passing the test of
\cref{dir:thm:smallquotients}; the sum has at most $\binom{r+3}3$ terms.

\paragraph{Sharpness at quarter retention.}
\begin{proposition}[One missing pattern at every odd arity; \stU]\label{dir:prop:missing}
Let $\Gamma=C_4$, $r$ odd, $c_1=c_3=r$ and $c_2=0$. Exactly one corrected vector in $\Gamma^r$ fails
\textup{(b)} of \cref{dir:thm:packing}, the vector with every entry $2$. Every proper subset of the
coordinates of any vector extends to a vector satisfying \textup{(b)}, and the four inequalities
\eqref{dir:eq:cyclic4} hold at the missing vector.
\end{proposition}

\begin{proof}
The four linear inequalities hold for every histogram with total at most $r$, so only the parity cut can fail,
and only for $\ell_2>0$. Since $r$ is odd,
$\floor{(r-\ell_1)/2}+\floor{(r-\ell_3)/2}=r-1-\floor{\ell_1/2}-\floor{\ell_3/2}$. If $\ell_1+\ell_3\ge1$ then
$\ell_1+\ell_3-\floor{\ell_1/2}-\floor{\ell_3/2}\ge1$, so this is at least $r-\ell_1-\ell_3\ge\ell_2$; if
$\ell_1=\ell_3=0$ the cut holds if and only if $\ell_2\le r-1$. So only $(0,r,0)$ fails, and setting an
unrecorded coordinate to zero extends a proper marginal.
\end{proof}

\begin{corollary}[An arithmetic instance of the parity cut; \stG]\label{dir:cor:paritycut}
There is $c>0$ such that for all large primes $q\equiv1\pmod4$ and every odd $r\le cL/\log L$ some retained set
of exactly $\ceil{N/4}$ primes and some quotas with $h\ge BL$ have the following property: every proper
marginal of $X$ has full support, and exactly $(\abs G/4)^r$ tuples, forming one Cartesian product of quartic
residue classes, are missing, although all four inequalities \eqref{dir:eq:cyclic4} hold at their projected
target.
\end{corollary}

\begin{proof}
Take a character $\theta$ of order four. By Fourier inversion and \cref{dir:lem:bias} each of its four value
classes in $\mathcal P$ has at least $(1-3\eps)N/4$ primes, and a largest one, of value $b$, at least
$\ceil{N/4}$. Retain $\ceil{N/4}-2r$ primes of value $b$ and $r$ primes of each of the values $bi$ and $bi^3$,
none of value $-b$. For large $q$ these choices exist, $2r/m\to0$, $b$ is the unique most frequent value, and
$\theta$ has bias $1-2r/m>1-\tau$ on $S$, so $H_\tau$ contains the cyclic group of order four generated by
$\theta$ and, having order at most four, equals it. Balanced quotas with $h\ge BL$ exist in the stated range of
$r$. \Cref{dir:prop:missing} and the quotient test before \cref{dir:thm:quarter} (\cref{dir:cor:anydensity} with
\cref{dir:thm:smallquotients}) give the assertions; in a proper marginal a residue tuple is
extended by giving an omitted coordinate corrected colour zero.
\end{proof}

Thus at the threshold itself the obstruction is arithmetic and is not detected by any of the linear
inequalities. A cruder failure already defeats the tests of \cref{dir:thm:disjoint-support}. Take primes
$q\equiv5\pmod{12}$: there is one quadratic character, none of order three, and a character $\theta$ of order
four. Retain a most frequent quartic class $c$, with at least $N/4$ primes, and exactly $r$ primes of the class
$ic$, which has $\gg N$ primes by expanding its indicator in $1,\theta,\theta^2,\theta^3$ and \cref{dir:lem:bias}.
The quadratic character has $r$ minority primes, so every test of \cref{dir:thm:disjoint-support} passes, but
prescribing $\theta(g_j)c^{-n_j}=-1$ in every recorded block needs two exceptional primes per block; in the
notation of \cref{dir:thm:smallquotients} this is $c_1=\ell_2=r$, which violates the third inequality of
\eqref{dir:eq:cyclic4}. So the tests of order two and three cease to suffice at $a=1/4$.

For the probabilities on the support put
\[
 D_q=\max_{g\in G^r}\bigl\lvert\abs G^r\dirP(X=g)-1\bigr\rvert,\qquad B_2(t)=1-2t,\qquad B_3(t)=\sqrt{1-3t(1-t)},
\]
and $\mathcal M_q=\max_\chi B_{\operatorname{ord}\chi}(t)^{s_\chi}$ over the characters of orders two and
three, with $0^0=1$ and an empty maximum zero.

\begin{theorem}[Pointwise mixing; \stG]\label{dir:thm:disjoint-mixing}
Let $1/4<a\le1$ and fix $r$. Uniformly in the retained primes and the quotas,
\[
 \mathcal M_q-O_{a,r}\Bigl(\frac{\log(2h)}h\Bigr)\le D_q\le(3^r-1)\mathcal M_q+O_{a,r}\Bigl(\frac{\log(2h)}h\Bigr),
\]
so $D_q\to0$ along an admissible sequence if and only if $\mathcal M_q\to0$. For $a>1/3$ the quadratic
characters suffice and $3^r-1$ may be replaced by $2^r-1$. For $r=1$ and $a>1/3$, \cref{dir:thm:fibre}
below gives the sharper relative law with error $O(1/h+q^{-A})$.
\end{theorem}

For $r\ge2$ one has $t\le1/3$, and $-\log B_d(t)$ is comparable to $t$ there, so uniformity is
equivalent to $t\min_\chi s_\chi\to\infty$. For one product the quadratic exception at $t=1/2$ is real:
half of the partitions put a single exceptional prime in the product, and its sign is balanced. A
cubic exception has $B_3(1/2)=1/2$, so boundedly many such primes can keep a detectable bias.

The proof compares the projected law with independent placement. For the generator of $H_\tau$, of
order $d\in\set{2,3}$, write the exceptional values as $b\xi_1,\dots,b\xi_s$, $s=s_\chi$, $\xi_\ell\ne1$,
put each of the $s$ labelled exceptions independently into box $i$ with probability $\alpha_i=n_i/m$,
multiply the $\xi_\ell$ in each recorded box and multiply block $j$ by $b^{n_j}$; call the law
$\widetilde\nu$.

\begin{lemma}[Independent placement; \stU]\label{dir:lem:disjoint-iid}
For fixed $r$, every Fourier coefficient of $\nu$ differs from that of $\widetilde\nu$ by
$O_{a,r}(\log(2h)/h)$.
\end{lemma}
\begin{proof}
In the actual partition the exceptions occupy $s$ distinct slots chosen without replacement. After
$\ell$ draws with box occupancies $H_i$ the next draw has probabilities $p_i=(n_i-H_i)/(m-\ell)$;
relative entropy is at most $\sum(p_i-\alpha_i)^2/\alpha_i$, the hypergeometric variance is
$\ell\alpha_i(1-\alpha_i)(m-\ell)/(m-1)$, and the chain rule bounds the divergence of the label
sequence from independence by $r\sum_{\ell<s}\ell/((m-1)(m-\ell))\ll_r(s/m)^2$, as $s\le2m/3$. Pinsker's
inequality gives a difference $O_r(s/m)$ for every Fourier coefficient. Alternatively, in a nontrivial
coefficient some exponent differs from that of a block of largest quota, and the corresponding ratio in
\cref{dir:lem:partition} is $\chi^{\pm1}$, with gap $2s/m$ for $d=2$ and at least $s/(2m)$ for $d=3$,
since $1-b_\chi^2=3((m-s)s+uv)/m^2\ge s/m$; so the actual coefficient is at most
$(m+1)^{r+1}e^{-hs/(2(r+1)m)}$. The independent coefficient has modulus $\abs z^s$ with
$z=\sum_i\alpha_i\zeta^{k_i}$, $\zeta=e^{2\pi i/d}$: for $d=2$, $\abs z\le1-2t$, and for $d=3$,
$\abs z^2=1-3(P_0P_1+P_1P_2+P_2P_0)\le1-3t(1-t)$ with $P_c$ the weights of the three phases; both are
attained by isolating a smallest block. So it is at most $e^{-cts}$. As $m\ll_{a,r}h^2$, the entropy
bound suffices when $s/m\le K_r\log(2h)/h$, and otherwise both coefficients are $O(h^{-1})$.
\end{proof}

\begin{proof}[Proof of \cref{dir:thm:disjoint-mixing}]
A character outside $H_\tau$ has $s_\psi\ge\tau m/2$, so $B_{\operatorname{ord}\psi}(t)^{s_\psi}=O(e^{-c\tau h})$.
Inside $H_\tau$ the largest nontrivial independent coefficient is exactly $B_d(t)^s$, conjugate cubic
generators having the same $s$, so by \cref{dir:lem:disjoint-iid} the largest nontrivial coefficient of
the projected law is $\mathcal M_q$ up to the stated error. A nontrivial coefficient is the average of
$\abs G^r\dirP(X=g)-1$ against a unit phase, so it is at most $D_q$; conversely Fourier inversion bounds
$D_q$ by the $d^r-1\le3^r-1$ nontrivial coefficients inside the quotient and the $q^{-2}$ bound for all
others.
\end{proof}

The comparison with independent placement also holds in total variation, uniformly in the number of
blocks, when the exceptions are few. Let $\Gamma$ have any fixed order, translate the colours by a most
frequent one, let $s=m-c_0$ be the number of exceptional primes, and let $\nu_{\mathrm{corr}}$ be the law of
the corrected projected block sums. Let $\widetilde\nu$ be the law obtained by putting each exceptional prime
independently into block $j$ with probability $n_j/m$, ignoring the baseline primes, and summing the
exceptional colours of each recorded block.

\begin{proposition}[Sparse exceptions, uniformly in $r$; \stU]\label{dir:prop:sparseexc}
With $(m)_s=m(m-1)\cdots(m-s+1)$,
\[
 d_{\mathrm{TV}}(\nu_{\mathrm{corr}},\widetilde\nu)\le1-\frac{(m)_s}{m^s}\le\frac{s(s-1)}{2m}
\]
for every $r$ and all quotas.
\end{proposition}

\begin{proof}
Regard block $j$ as $n_j$ slots. Under a uniform allocation the slots of the $s$ labelled exceptions have the
law of $s$ independent uniform slots conditioned to be distinct. The total variation distance between a law
and its conditioning on an event of positive probability is the probability of the complement, here at most
$\binom s2/m$. Mapping slot positions to block sums does not increase total variation.
\end{proof}

The Fourier coefficient of $\widetilde\nu$ at $\boldsymbol\psi\in\widehat\Gamma^r$ is exactly
\[
 \prod_{a\ne0}\Bigl(\frac{n_0}m+\sum_{j=1}^r\frac{n_j}m\psi_j(a)\Bigr)^{c_a},
\]
so for $s=o(\sqrt m)$ the colour counts $c_a$ and the relative quotas alone determine the leading profile of
the projected law, without the identities of the primes. Since the relative error of
\cref{dir:thm:disjoint-growing} gives total variation at most $q^{-A}/2$ between the law of $X$ and the uniform
lifting of its quotient law, the uniform lifting of the independent-exception model is within
$q^{-A}/2+s(s-1)/(2m)$ of the law of $X$, after undoing the translation. In \cref{dir:cor:paritycut},
$s=2r=O(L/\log L)$ and $m\asymp L^2/\log L$, so this error is $O(1/\log L)$.

\paragraph{One product: the quadratic-fibre law.}
For one product above density one third the single possibly large mode can be computed exactly and
approximated with an error proportional to its own size, which gives relative counts on both sign
classes, including a class whose mass tends to zero. For $m\ge2$, $0\le h\le m$ and $0\le s\le m$ let
\[
 P_{m,s}(h)=\binom ms^{-1}\sum_{j=0}^s(-1)^j\binom hj\binom{m-h}{s-j}=\dirE(-1)^J,
\]
where $J$ is the number of marked elements in a uniform $h$-subset of an $m$-set with $s$ marked
elements; by symmetry $J$ is also the number of successes in $s$ draws from a population with $h$
successes. It is the coefficient $E$ of \cref{dir:rem:thresholds}.

\begin{lemma}[Uniform parity decay; \stU]\label{dir:lem:paritydecay}
There is an absolute $K$ with $\abs{P_{m,s}(h)}\le K\exp(-hs/(4m))$ for $1\le h\le m/2$ and $0\le s\le m/2$.
\end{lemma}

\begin{proof}
Assume $m\ge8$. Pair $\rho=\min(s,\floor{m/4})$ marked elements with distinct unmarked ones. Conditional on
the selection outside the pairs and on the number selected in each pair, the orientations of the pairs
with exactly one selected member are uniform and independent, so if any pair splits the conditional mean
of $(-1)^J$ vanishes; hence $\abs{P_{m,s}(h)}\le\dirP(\text{no split})$. Select each element
independently with probability $u=h/m$ and condition on the total $B_m=h$. Without conditioning,
no split has probability $(1-2u(1-u))^\rho\le e^{-u\rho}$; given no split, the $m-2\rho\ge m/2$ unpaired
elements are still independent, so $\dirP(B_m=h\mid\text{no split})\le\sup_j\dirP(\mathrm{Bin}(m-2\rho,u)=j)\ll
h^{-1/2}$ by Fourier inversion, while $\dirP(B_m=h)\gg h^{-1/2}$ since $h$ is a mode of variance at most $h$.
Bayes' formula and $\rho\ge s/4$ finish the proof.
\end{proof}

\begin{theorem}[Localized binomial approximation; \stU]\label{dir:thm:parity}
For $2\le h\le m/2$ and $0\le s\le m/2$, with $t=h/m$ and $\lambda=ts$,
\[
 \bigl|P_{m,s}(h)-(1-2t)^s\bigr|\ll\frac{\lambda^2e^{-\lambda/16}}h\ll\frac{\min(1,\lambda)}h,
\]
with absolute constants; the difference vanishes for $s=0,1$.
\end{theorem}

\begin{proof}
Let $J_i$ count successes in the first $i$ draws without replacement from the population with $h$
successes and $F_i=\dirE(-1)^{J_i}=P_{m,i}(h)$. The next draw succeeds with conditional probability
$(h-J_i)/(m-i)$, so
\[
 F_{i+1}=(1-2t)F_i+\frac2{m-i}\dirE\bigl[(-1)^{J_i}(J_i-it)\bigr],\qquad
 \dirE\bigl[J_i(-1)^{J_i}\bigr]=-it\,P_{m-1,i-1}(h-1),
\]
the identity by exchangeability and conditioning on a specified successful draw. By
\cref{dir:lem:paritydecay}, $\abs{F_i}\ll e^{-ti/4}$ and $\abs{P_{m-1,i-1}(h-1)}\ll e^{-ti/16}$, since
$(h-1)(i-1)/(m-1)\ge hi/(4m)$ for $i,h\ge2$, and the case $i=1$ is trivial. So the correction term is at
most $K_1(ti/m)e^{-ti/16}$, and iterating from $F_0=1$ with $0\le1-2t\le e^{-2t}$ gives
$\abs{F_s-(1-2t)^s}\le K_1(t/m)e^{-t(s-1)/16}\sum_{i<s}i\ll(\lambda^2/h)e^{-\lambda/16}$, as $ts^2/m=\lambda^2/h$.
\end{proof}

If $h\to\infty$, $h/m\to0$ and $hs/m\to\lambda_0\in(0,\infty)$, then
$h\bigl(P_{m,s}(h)-(1-2t)^s\bigr)\to-2\lambda_0^2e^{-2\lambda_0}$; this is the case $F=C_2$ of
\cref{dir:cor:cp} below, \eqref{dir:eq:paritylimit}. So the order $1/h$ of the error cannot be improved,
even in the sparse critical window. For one product,
with no reserve block, $h=\min(k,m-k)$ and $t=h/m$; let $s$ be the least number of retained primes outside
one fibre of a nonprincipal quadratic character, attained by a pair $(\chi_*,\sigma)$ (ties allowed), put
$\lambda=ts$ and $\epsilon_k=\sigma^k(-1)^{s[k>m/2]}$, and let
$R(g)=\#\set{T\subseteq S:\abs T=k,\ \prod_{p\in T}p\equiv g\pmod q}$.

\begin{theorem}[Quadratic-fibre law; \stG]\label{dir:thm:fibre}
Fix $1/3<a\le1$ and $A>0$. There are $B,C>0$, depending only on $a,A$, such that for all large $q$, with
$y=CL^2$, all $S\subseteq\mathcal P$ with $m\ge aN$ and all $k$ with $h\ge BL$:
\begin{enumerate}[label=\textup{(\alph*)}]
\item if $s=0$, then $R(g)=0$ when $\epsilon_k\chi_*(g)=-1$, and $\varphi(q)R(g)/\binom mk=2+O(q^{-A})$
otherwise;
\item if $s\ge1$, every unit occurs and, uniformly in $q,S,k,g$,
\[
 \frac{\varphi(q)R(g)}{\binom mk}=\bigl(1+\epsilon_k\chi_*(g)B_2(t)^s\bigr)\Bigl(1+O\Bigl(\frac1h+q^{-A}\Bigr)\Bigr);
\]
\item in all cases, with an absolute $c>0$,
$\varphi(q)R(g)/\binom mk=1+\epsilon_k\chi_*(g)B_2(t)^s+O\bigl(\lambda^2e^{-c\lambda}/h+q^{-A-2}\bigr)$.
\end{enumerate}
\end{theorem}

\begin{proof}
For $a>1/3$ the Gram argument of \cref{dir:thm:quotient} gives $\abs{H_\tau}\le2$, so at most one
nonprincipal character has bias above $1-\tau$, and it is quadratic. A quadratic character $\chi$ has bias
$1-2s_\chi/m$, so $\chi_*$ maximizes the quadratic bias, and every nonprincipal character other than
$\chi_*$ has bias at most $1-\tau$. With $B=2(A+5)/\tau$, \cref{dir:lem:circle} and Fourier inversion give
\[
 \frac{\varphi(q)R(g)}{\binom mk}=1+\chi_*(g)E_k(\chi_*)+\mathcal E_g,\qquad\abs{\mathcal E_g}\le q(m+1)e^{-\tau h/2}\le q^{-A-2}.
\]
For $k\le m/2$ each product has the majority sign raised to the power $k$ times the parity of its minority
primes, and for $k>m/2$ the selected minority count is $s$ minus the omitted one, so
$E_k(\chi_*)=\epsilon_kP_{m,s}(h)$; \cref{dir:thm:parity} gives (c). If $s=0$ every product has sign
$\epsilon_k$. If $s\ge1$, the smaller profile is $1-B_2(t)^s\ge1-e^{-2\lambda}\ge(1-e^{-2})\min(1,\lambda)$,
so the finite-sampling error is $O(1/h)$ relative to either profile, and since $\lambda\ge h/m$ dividing
$\mathcal E_g$ by the profile costs at most $O(m/h)=O(L^2)$.
\end{proof}

\begin{corollary}[Equidistribution criterion and critical window; \stG]\label{dir:cor:fibre}
In the setting of \cref{dir:thm:fibre}, $D_q=\max_g\abs{\varphi(q)R(g)/\binom mk-1}=B_2(t)^s+O(1/h+q^{-A})$, and
the total variation distance from uniform is $\frac12B_2(t)^s+O(1/h+q^{-A})$. So $D_q\to0$ if and only
if $B_2(t)^s\to0$; when $h=o(m)$ this means $hs/m\to\infty$. If $h=o(m)$ and $hs/m\to\lambda_0\in(0,\infty)$,
then uniformly in $g$
\[
 \frac{\varphi(q)R(g)}{\binom mk}=1+\epsilon_k\chi_*(g)\Bigl(B_2(t)^s-\frac{2\lambda_0^2e^{-2\lambda_0}}h\Bigr)+o(1/h),
\]
so $D_q\to e^{-2\lambda_0}$ and $h(D_q-B_2(t)^s)\to-2\lambda_0^2e^{-2\lambda_0}$. If $s\ge1$ and $hs/m\to0$,
the normalized count on the minority class is asymptotic to $2hs/m$; if $h/m\to v\in(0,1/2]$ and $s$ is
fixed, then $D_q\to(1-2v)^s$, so nearly balanced quotas mix with a single minority prime.
\end{corollary}

\begin{proof}
Both sign classes are nonempty of equal size; apply (c), the reverse triangle inequality and averaging.
If $t\to0$ then $\log(1-2t)=-2t+O(t^2)$. In the critical window the residual in the proof of
\cref{dir:thm:fibre} is $o(1/h)$ since $h\le m=O(L^2)$; apply the limit \eqref{dir:eq:paritylimit}.
\end{proof}

The criterion replaces a fixed density margin above one half by the number of minority primes in a
typical sample: at sparse quotas $h\asymp\log q$ one has $m\asymp(\log q)^2/\log\log q$, so the crossover
minority count is of order $\log q/\log\log q$. By \cref{dir:lem:paritydecay} the exceptional mode is
$O(q^{-A})$ once $hs/m\ge4(A+1)\log q$, but $hs/m\to\infty$ alone gives only $o(1)$: with
$hs/m=\log\log h+o(1)$ one has $B_2(t)^s\asymp(\log h)^{-2}$, so no fixed power of $q$ bounds the
deviation. For $y=L^F$ with fixed $F>2$ and $a>1/3$ the one-mode reduction in the proof of
\cref{dir:thm:fibre} holds for $h\ge cL/\log L$ with $c>2(F-1)/(F(F-2))$, the remainder being replaced
by the bound $\exp(-\gamma L/\log L)$ of \cref{dir:thm:short}; so $D_q=\abs{P_{m,s}(h)}+O(\exp(-\gamma L/\log L))$,
and \cref{dir:lem:disjoint-iid} gives $D_q=B_2(t)^s+O(\log(2h)/h)$ since $m$ is polynomial in $h$.

\begin{proposition}[Realized transitions and the density floor; \stG]\label{dir:prop:fibresharp}
\begin{enumerate}[label=\textup{(\alph*)}]
\item For prime $q$, $y=CL^2$ with $C$ large and $h=k=\ceil{BL}$, every value $e^{-2\lambda_0}$,
$\lambda_0\ge0$, is the limit of $D_q$ along retained sets of density at least $1/2$: the two normalized
class counts tend to $1\pm e^{-2\lambda_0}$.
\item At retained density $1/3$ no criterion using only distances from quadratic fibres suffices: for
primes $q\equiv1\pmod6$ there are retained sets of $\ceil{N/3}$ primes with $s$ a fixed positive fraction
of $m$ all of whose $k$-products lie in one fibre of a cubic character, a third of the units.
\end{enumerate}
\end{proposition}

\begin{proof}
(a) For prime $q$ there is one quadratic character; let $S_0$ be its larger sign class in $\mathcal P$,
of size $m_0\ge N/2$, the other class having at least $(1-\eps)N/2$ elements by \cref{dir:lem:bias}. Add
$s=\floor{\lambda_0m_0/h}$ primes of the other class; then $hs/m\to\lambda_0$ and \cref{dir:cor:fibre}
applies. For $\lambda_0=0$ either add none, or add one prime, which gives full support with limiting
deviation one; adding $\floor{(m_0/h)\sqrt h}$ primes gives limit zero.
(b) Use the construction of \cref{dir:rem:thresholds} for a character $\eta$ of order three: $S$ is a set of
$\ceil{N/3}$ primes in the largest fibre of $\eta$, so every $k$-product lies in the fibre $\eta(g)=\eta(S)^k$. With
$\chi$ the quadratic character, the six characters generated by $\eta$ and $\chi$ have bias at most $\eps<1/10$
on $\mathcal P$, so each joint fibre has at most $(1+5\eps)N/6$ primes and either sign class of $\chi$ in $S$ has
at least $(1-5\eps)N/6$; thus $s$ stays a fixed positive fraction of $m$.
\end{proof}

Exact rational checks of the parity identities over all $44\,280$ triples with $m\le80$ give
$\max h\abs{P_{m,s}(h)-(1-2t)^s}=2/3$, and at $m=h^2$, $h=640$, the quantity $h(P_{m,s}(h)-(1-2t)^s)$ is
within $7\cdot10^{-4}$ of the limiting correction for $\lambda_0=0.5,1,2$.
\dirfiles{\artifact{grh_quadratic_fibre/verification.txt}, \artifact{grh_quadratic_fibre/critical_window.csv}}

\begin{example}\label{dir:ex:disjoint}
Modulo $11$ keep the smallest $18-s$ residue primes and $s$ nonresidue primes up to $199$, excluding
$11$; modulo $13$ use the smallest requested numbers of primes up to $239$ in the three classes
$p^4\equiv1,3,9$, with counts $15-u-v,u,v$. With two recorded blocks and the reserve:
\begin{center}\small
\begin{tabular}{@{}llllr@{}}
\toprule
modulus & class data & quotas (reserve first) & occupied classes & residue pairs\\
\midrule
$11$ & $s=0$ & $6,6,6$ & $1/4$ & $25/100$\\
 & $s=1$ & $6,6,6$ & $3/4$ & $75/100$\\
 & $s=2$ & $6,6,6$ & $4/4$ & $100/100$\\
$13$ & $(u,v)=(0,0)$ & $5,5,5$ & $1/9$ & $16/144$\\
 & $(u,v)=(1,1)$ & $5,5,5$ & $7/9$ & $112/144$\\
 & $(u,v)=(2,0)$ & $5,5,5$ & $6/9$ & $96/144$\\
 & $(u,v)=(2,1)$ & $5,5,5$ & $9/9$ & $144/144$\\
\bottomrule
\end{tabular}
\end{center}
For $s=2$ modulo $11$ there are $17\,153\,136$ allocations; the sign class $++$ gets $5\,045\,040$ and
each other class $4\,036\,032$. Modulo $13$, two exceptions of opposite types give strictly larger
support than two of the same type. Urn checks cover $2\,464$ binary cases, $15\,072$ binary Fourier
coefficients and $104\,505$ cubic target vectors; the table repair is checked on $28\,183$ pairs of
tables, and for moduli $15,21$ with quotas $7,7$ exactly the $48$ of $96$ unit pairs with
$g_1g_2\equiv2\pmod3$ occur, with $65$ to $80$ representations each.
For $C_4$ and $C_2\times C_2$ the criteria of \cref{dir:thm:smallquotients,dir:thm:capacity} agree with an
allocation oracle that sends one item at a time to a block or the reserve and uses neither the
inequalities nor the witness table, in $479\,682$ support and $10\,206$ coverage comparisons with every
supply from $0$ to $8$ and $1$ to $7$ recorded blocks; $5\,280$ fixed-quota targets were also checked
against enumeration of colour-count tables and the dynamic program of \cref{dir:cor:anydensity},
the construction of \cref{dir:thm:packing} in $2\,564$ comparisons with the fixed-quota oracle, and
\cref{dir:prop:missing} for odd $r\le23$.
\dirfiles{\artifact{disjoint-prime-products/checks/results.txt},
\artifact{disjoint-prime-products/checks/repair_results.txt},
\artifact{grh_quarter_density/verification.txt}}
\end{example}

\subsection{Sparse transition profiles at every density}\label{dir:sec:profiles}
\sources{grh_extensions_2026-11-27_3}

For one product, \cref{dir:thm:fibre} describes the approach to uniformity above density one third,
where at most one quadratic character survives in the bounded quotient. Below that density the quotient
of \cref{dir:thm:quotient} may be cyclic of order four or a product of quadratic quotients, and several
fibres interact. In the sparse range $BL\le h=o(m)$ this subsection gives a necessary and sufficient
mixing criterion at every fixed retained density, shows that every subsequential limit of the normalized
residue profile is, after a deterministic phase correction, the lift of a compound Poisson law on a
quotient of order at most $\floor{1/a}$, and shows that every such law occurs. The probabilistic input is
an expansion of a fixed-size sample from a population with few defects, with a signed correction of
order $1/h$.

Keep the notation of \cref{dir:sec:products} with $y=CL^2$, fix $0<a\le1$, put $D=\floor{1/a}$, and let
$S\subseteq\mathcal P$ have $m\ge aN$ elements. For $1\le k<m$ put $h=\min(k,m-k)$, $t=h/m$ and
$U_S=\prod_{p\in S}p\bmod q$, and let
\[
 Y=X_k\quad(k\le m/2),\qquad Y=U_SX_k^{-1}\quad(k>m/2),
\]
which has the law of the product of a uniform $h$-subset of $S$; translation and inversion preserve the
distance from uniform. Let $R(g)=\binom mk\dirP(X_k=g)$ and
$D_q=\max_{g\in G}\abs{\varphi(q)R(g)/\binom mk-1}$, as in \cref{dir:cor:fibre}. For a nonprincipal
character $\chi$ modulo $q$ put
\[
 s_\chi(S)=m-\max_\zeta\#\set{p\in S:\chi(p)=\zeta},\qquad s_D(S)=\min_{1<\ord\chi\le D}s_\chi(S),
\]
the number of retained primes outside a largest fibre of $\chi$, and its least value over the
characters of order at most $D$; an empty minimum is $+\infty$. For characters of order two or three
$s_\chi(S)$ is the number $s_\chi$ of \cref{dir:thm:disjoint-support}.

\paragraph{Compound Poisson laws.} Let $F$ be a finite abelian group, written multiplicatively with
identity $e$. Signed measures on $F$ form an algebra under convolution $*$, with the norm
$\norm{\mu}_1=\sum_{x\in F}\abs{\mu(x)}$; convolution is submultiplicative, a probability measure has
norm one, and the total variation distance is half the norm of the difference. For rates $\lambda_x\ge0$,
$x\in F\setminus\set e$, let $J_x$ be independent Poisson variables of means $\lambda_x$ and
\begin{equation}\label{dir:eq:cp}
 P_\Lambda=\mathcal L\Bigl(\prod_{x\ne e}x^{J_x}\Bigr)=\exp_*(\Lambda),\qquad
 \Lambda=\sum_{x\ne e}\lambda_x(\delta_x-\delta_e),
\end{equation}
where $\mathcal L$ denotes the law and the exponential is taken in the convolution algebra. Its Fourier
coefficients are
\begin{equation}\label{dir:eq:cpfourier}
 \widehat{P_\Lambda}(\chi)=\exp\Bigl(\sum_{x\ne e}\lambda_x\bigl(\chi(x)-1\bigr)\Bigr)\ne0 .
\end{equation}
So a compound Poisson law on a nontrivial group is never uniform. Its support is the subgroup generated by
the jumps of positive rate, since in a finite group the semigroup generated by a set is the subgroup it
generates.

\paragraph{Two finite sampling estimates.} The first extends \cref{dir:lem:paritydecay} to roots of
unity of any bounded order. It has no growing factor in front of the exponential, so the minority scale
$hs/m$ may tend to infinity arbitrarily slowly.

\begin{lemma}[Split-pair decay; \stU]\label{dir:lem:splitpair}
Let $2\le d\le D$, let $z_1,\dots,z_m$ be roots of unity of order dividing $d$, put
$s=m-\max_\zeta\#\set{j:z_j=\zeta}$, and let $1\le h\le m/2$. If $E_h$ is the mean of $\prod_{j\in T}z_j$
over a uniform $h$-subset $T$ of $\set{1,\dots,m}$, then
\begin{equation}\label{dir:eq:splitpair}
 \abs{E_h}\le20\exp\Bigl(-c_D\frac{hs}m\Bigr),\qquad c_D=\frac{1-\cos(\pi/D)}{\max(D,8)} .
\end{equation}
\end{lemma}

\begin{proof}
If $m<8$ then $hs/m<4$ and $c_D\le1/8$, so the right side exceeds one. Let $m\ge8$. A largest value
class has $n_{\max}\ge m/D$ elements. Pair
\[
 \rho=\min\bigl(n_{\max},s,\floor{m/4}\bigr)\ge\frac s{\max(D,8)}
\]
of its labels with distinct labels of other values; at least $m/2$ labels remain unpaired. Conditional on
the number of selected labels in every pair and on the selected unpaired labels, the orientations of the
split pairs (those with exactly one selected member) are independent fair choices. For distinct $d$th
roots of unity $z,z'$ one has $\abs{(z+z')/2}\le\cos(\pi/d)\le\kappa:=\cos(\pi/D)$. Hence
$\abs{E_h}\le\dirE(\kappa^Z\mid B_m=h)$, where the subset is generated by independent Bernoulli$(t)$
selection of all labels, $t=h/m$, conditioned on the total $B_m=h$, and $Z$ is the number of split pairs.

Under independent selection $\dirE\kappa^Z=(1-2(1-\kappa)t(1-t))^\rho$. Tilt the independent law by the
nonnegative weight $\kappa^Z$; this is well defined also for $\kappa=0$, since the event of no split pair
has positive probability. Under the tilt the unpaired labels remain independent Bernoulli$(t)$,
independent of the pairs, and
\[
 \dirE(\kappa^Z\mid B_m=h)=\dirE\kappa^Z\,\frac{\dirP_{\mathrm{tilt}}(B_m=h)}{\dirP(B_m=h)} .
\]
The denominator is at least $3/(20\sqrt h)$: the integer mean $h$ is a mode, and by Chebyshev's inequality
at least $3/4$ of the mass lies in an interval containing at most $5\sqrt h$ integers. With
$n=m-2\rho\ge m/2$, Fourier inversion gives
\[
 \sup_j\dirP\bigl(\mathrm{Bin}(n,t)=j\bigr)\le\frac1{2\pi}\int_{-\pi}^{\pi}e^{-2nt(1-t)\sin^2(\theta/2)}\,d\theta
 \le\frac{\sqrt\pi}{2\sqrt{2nt(1-t)}}\le\frac{\sqrt\pi}{\sqrt{2h}} .
\]
The tilted total is the sum of this unpaired binomial and an independent variable, so its atoms obey the
same bound, and the ratio of atoms is less than $20$. Finally $t\le1/2$ gives
$(1-2(1-\kappa)t(1-t))^\rho\le e^{-(1-\kappa)t\rho}\le e^{-c_Dts}$.
\end{proof}

For $d=2$ the pairing argument of \cref{dir:lem:paritydecay} gives the better exponent $hs/(4m)$.

For the second estimate consider a population of $m-s$ labels equal to $e$ and $s$ labelled defects
$x_1,\dots,x_s\ne e$ in $F$, and let $\mu_{m,h}$ be the law of the product over a uniform $h$-subset.
Assume $1\le h\le m/2$ and $0\le s\le m/2$, and put $t=h/m$ and $\lambda=ts$. Let
\begin{equation}\label{dir:eq:nu}
 \nu(v)=\mathop{\ast}_{i=1}^s\bigl(\delta_e+v(\delta_{x_i}-\delta_e)\bigr),\qquad0\le v\le1,
\end{equation}
a polynomial in $v$ with values in signed measures. Thus $\nu(t)$ includes each defect independently with
probability $t$, whereas $\mu_{m,h}$ respects the fixed quota; derivatives of $\nu$ are derivatives in $v$.

\begin{theorem}[Uniform second-order correction; \stU]\label{dir:thm:fpc}
In this setting
\begin{equation}\label{dir:eq:fpcfirst}
 \norm{\mu_{m,h}-\nu(t)}_1\le\frac{4\lambda^2e^{2\lambda}}h ,
\end{equation}
and more precisely
\begin{equation}\label{dir:eq:fpcsecond}
 \Bigl\|\mu_{m,h}-\nu(t)+\Bigl(\frac1{2h}-\frac1{2m}\Bigr)t^2\nu''(t)\Bigr\|_1\le\frac{\mathcal R(\lambda)}{h^2},
 \qquad
 \mathcal R(\lambda)=2e^{2\lambda}\bigl((2\lambda)^4+6(2\lambda)^3+7(2\lambda)^2+2\lambda\bigr).
\end{equation}
In particular the remainder is $O_M(h^{-2})$ uniformly for $\lambda\le M$, independently of the finite
group and of the defect types. The correction satisfies $\norm{t^2\nu''(t)}_1\le4t^2s(s-1)\le4\lambda^2$.
\end{theorem}

\begin{proof}
Put $d_i=\delta_{x_i}-\delta_e$, and let $e_j$ be the sum of all $j$-fold convolutions of distinct
$d_i$, with $e_0=\delta_e$. A specified $j$-set of labels is contained in a uniform $h$-subset with
probability $(h)_j/(m)_j$, where $(n)_j=n(n-1)\cdots(n-j+1)$ and $(h)_j=0$ for $j>h$. Expanding the
selected factors gives
\begin{equation}\label{dir:eq:factorial}
 \mu_{m,h}=\sum_{j=0}^s\frac{(h)_j}{(m)_j}e_j,\qquad\nu(t)=\sum_{j=0}^st^je_j,\qquad
 \norm{e_j}_1\le\binom sj2^j .
\end{equation}
Set $r_j=(h)_j/((m)_jt^j)$, so $r_j=0$ for $j>h$. We claim
\begin{equation}\label{dir:eq:ratios}
 0\le1-r_j\le\frac{j(j-1)}h,\qquad
 \Bigl|r_j-1+\frac{j(j-1)}2\Bigl(\frac1h-\frac1m\Bigr)\Bigr|\le\frac{2j^4}{h^2}.
\end{equation}
For $j\le h$ write $r_j=\prod_{\ell=0}^{j-1}(1-a_\ell)$ with $a_\ell=(\ell/h-\ell/m)/(1-\ell/m)$. Since
$j\le s\le m/2$, one has $0\le a_\ell\le1$ and $a_\ell\le2\ell/h$. The union bound and its second-order
Bonferroni refinement give
\[
 0\le1-r_j\le\sum_\ell a_\ell,\qquad0\le\sum_\ell a_\ell-(1-r_j)\le\frac12\Bigl(\sum_\ell a_\ell\Bigr)^2,
\]
and
\[
 \Bigl|\sum_\ell a_\ell-\Bigl(\frac1h-\frac1m\Bigr)\sum_\ell\ell\Bigr|\le\frac2{hm}\sum_{\ell<j}\ell^2
 \le\frac{2j^3}{3h^2}.
\]
These imply \eqref{dir:eq:ratios}. For $j>h$ the claims follow from $r_j=0$ and $j\ge h+1$, and the cases
$j=0,1$ are exact.

Insert the first inequality of \eqref{dir:eq:ratios} into \eqref{dir:eq:factorial}. Since
\[
 \sum_jj(j-1)\binom sj(2t)^j=4t^2s(s-1)(1+2t)^{s-2}\le4\lambda^2e^{2\lambda},
\]
with the cases $s<2$ read separately, this gives \eqref{dir:eq:fpcfirst}. For \eqref{dir:eq:fpcsecond}
use the second inequality of \eqref{dir:eq:ratios}, the identity $t^2\nu''(t)=\sum_jj(j-1)t^je_j$, and
$\binom sjt^j\le\lambda^j/j!$; the remainder is at most
$(2/h^2)\sum_{j\ge0}j^4(2\lambda)^j/j!=\mathcal R(\lambda)/h^2$. The bound for the correction follows from
its expression as a sum over ordered pairs of defects: each summand is
$t^2d_i*d_{i'}$ convolved with probability measures, of norm at most $4t^2$.
\end{proof}

\begin{corollary}[Critical compound Poisson law and its first correction; \stU]\label{dir:cor:cp}
Suppose $h\to\infty$, $t=h/m\to0$, and $\lambda_x^{(m)}=t\,\#\set{i:x_i=x}\to\lambda_x<\infty$ for every
$x\ne e$. With $\Lambda$ as in \eqref{dir:eq:cp},
\begin{equation}\label{dir:eq:cpcorrection}
 \mu_{m,h}\longrightarrow P_\Lambda,\qquad h\bigl(\mu_{m,h}-\nu(t)\bigr)\longrightarrow-\tfrac12\Lambda^{*2}*P_\Lambda
\end{equation}
in $\norm\cdot_1$. Moreover, if $\Lambda_m=\sum_x\lambda_x^{(m)}(\delta_x-\delta_e)$, then
$\norm{\mu_{m,h}-\exp_*(\Lambda_m)}_1\le4\lambda^2e^{2\lambda}/h+2t\lambda$.
\end{corollary}

\begin{proof}
The Bernoulli$(t)$ law and the Poisson$(t)$ law of the number of jumps of one label differ in $\ell^1$ by
$2t(1-e^{-t})\le2t^2$. Pushing both forward by $j\mapsto x_i^j$ and telescoping the $s$ convolutions gives
$\norm{\nu(t)-\exp_*(\Lambda_m)}_1\le2st^2=2t\lambda$; with \cref{dir:thm:fpc} and the convergence of the
finitely many rates this proves the first assertion and the last bound.

For the correction fix a character $\chi$ of $F$ and put $\alpha_i=\chi(x_i)-1$. Then
$\widehat{\nu(t)}(\chi)=\prod_i(1+t\alpha_i)\to\exp(\sum_x\lambda_x(\chi(x)-1))$. For small $t$ the
factors are nonzero, and
\[
 t^2\widehat{\nu''(t)}(\chi)=\widehat{\nu(t)}(\chi)\Bigl[\Bigl(\sum_i\frac{t\alpha_i}{1+t\alpha_i}\Bigr)^2
 -\sum_i\frac{t^2\alpha_i^2}{(1+t\alpha_i)^2}\Bigr].
\]
The first sum tends to $\widehat\Lambda(\chi)$ and the second is $O(t\lambda)\to0$. Finite Fourier
inversion gives $t^2\nu''(t)\to\Lambda^{*2}*P_\Lambda$ in norm. Multiply \eqref{dir:eq:fpcsecond} by $h$:
the coefficient $h(1/(2h)-1/(2m))$ tends to $1/2$ and the remainder tends to zero.
\end{proof}

The correction in \eqref{dir:eq:cpcorrection} is relative to the current Bernoulli law $\nu(t)$, not to
the limit $P_\Lambda$: the drift of the rates and the Poisson approximation error $O(t\lambda)$ can exceed
$1/h$. At $h\asymp L$ and $m\asymp L^2/\log L$, for example, $h^2/m\asymp\log L$; \eqref{dir:eq:fpcsecond}
avoids this by keeping $t$ and every defect count exactly. For $F=C_2$ with $s$ defects the nontrivial
Fourier coefficients of $\mu_{m,h}$ and $\nu(t)$ are $P_{m,s}(h)$ and $(1-2t)^s$, and $\widehat\Lambda=-2\lambda_0$,
so \cref{dir:cor:cp} gives
\begin{equation}\label{dir:eq:paritylimit}
 h\bigl(P_{m,s}(h)-(1-2t)^s\bigr)\longrightarrow-2\lambda_0^2e^{-2\lambda_0}
 \qquad\Bigl(h\to\infty,\ \frac hm\to0,\ \frac{hs}m\to\lambda_0\Bigr),
\end{equation}
the sharp correction to \cref{dir:thm:parity}. For one product, \eqref{dir:eq:fpcfirst} compares the
law with independent inclusion of the exceptions within $O_M(1/h)$ in total variation whenever
$hs/m\le M$, without the restriction $s=o(\sqrt m)$ of \cref{dir:prop:sparseexc}, and
\eqref{dir:eq:fpcsecond} improves the remainder to $O_M(h^{-2})$; \cref{dir:prop:sparseexc} is uniform in
the number of blocks and in $hs/m$.

\paragraph{The mixing criterion.}

\begin{theorem}[Mixing criterion at every fixed density; \stG]\label{dir:thm:allmix}
For every $0<a\le1$ and $A>0$ there are $B,C>0$ such that, along every sequence $q\to\infty$ with
$y=CL^2$, $m\ge aN$, $h\ge BL$ and $h/m\to0$,
\[
 D_q\longrightarrow0\quad\Longleftrightarrow\quad\frac{h\,s_D(S)}m\longrightarrow+\infty .
\]
For $D\ge2$ and all large $q$, without the assumption $h/m\to0$,
\begin{equation}\label{dir:eq:allmix}
 D_q\le20(D-1)\exp\Bigl(-c_D\frac{h\,s_D(S)}m\Bigr)+q^{-A},
\end{equation}
with $c_D$ as in \cref{dir:lem:splitpair}. For $D=1$, $D_q\le q^{-A}$.
\end{theorem}

\begin{proof}
Apply \cref{dir:thm:quotient} with $k=h$, the $h$-subsets being those defining $Y$, and $B$ enlarged in
terms of $A$: the proof gives $\abs{H_\tau}\le D$ and
$\sum_{\chi\notin H_\tau}\abs{\dirE\chi(Y)}\le q(m+1)e^{-\tau h/2}\le q^{-A}$. Every nonprincipal
character in $H_\tau$ has order at most $D$, so $s_\chi(S)\ge s_D(S)$, and \cref{dir:lem:splitpair}
bounds its coefficient by $20\exp(-c_Dhs_D(S)/m)$. There are at most $D-1$ such characters, and none for
$D=1$. Fourier inversion gives \eqref{dir:eq:allmix}, and with it the sufficiency of the criterion.

For necessity suppose that $\lambda=hs_D(S)/m\le M<\infty$ on a subsequence. The minimum is then over a
nonempty set; choose a minimizing character $\chi$ and a most frequent value $\zeta$ of $\chi$ on $S$, and
divide all values by $\zeta$. In the cyclic group $\chi(G)$ the population has $m-s$ identities and
$s=s_D(S)$ defects, with $s/m=\lambda/h\to0$, so $s\le m/2$ eventually. Let $Q$ be the $\chi$-coefficient
under independent inclusion of these defects. \Cref{dir:thm:fpc} gives
$\abs{\zeta^{-h}\dirE\chi(Y)-Q}\le4\lambda^2e^{2\lambda}/h$. For every defect value $z$ on the unit circle
and $t<1/2$, $\abs{1+t(z-1)}\ge1-2t$, so $\abs Q\ge(1-2t)^s\ge\exp(-2\lambda/(1-2t))$. As $t\to0$ and
$\lambda\le M$, $\abs{\dirE\chi(Y)}$ stays bounded away from zero. But for every nonprincipal $\chi$
\[
 \abs{\dirE\chi(Y)}=\Bigl|\frac1{\varphi(q)}\sum_g\chi(g)\bigl(\varphi(q)\dirP(Y=g)-1\bigr)\Bigr|\le D_q ,
\]
so $D_q$ does not tend to zero.
\end{proof}

For $a>1/4$ one has $D\le3$, and for $t\to0$ both $-\log B_2(t)$ and $-\log B_3(t)$ are of order $t$; so
for one product in the sparse range \cref{dir:thm:allmix} agrees with the criterion of
\cref{dir:thm:disjoint-mixing}, and for $a>1/3$ with that of \cref{dir:cor:fibre}. Below density one quarter
the criterion involves all character orders up to $D$.

\paragraph{Classification of the limiting profiles.}

\begin{theorem}[Sparse transition profiles; \stG]\label{dir:thm:profiles}
The constants $B,C$ of \cref{dir:thm:allmix} may be chosen so that the following holds. Every sequence
$q_n\to\infty$ with $m_n\ge aN_n$, $h_n\ge BL_n$ and $h_n/m_n\to0$ has a subsequence on which there are a
fixed finite abelian group $F$ with $\abs F\le D$, surjective homomorphisms $\pi_n\colon G_{q_n}\to F$,
phases $b_n\in F$ and finite rates $\lambda_x\ge0$, $x\in F\setminus\set e$, such that
\begin{equation}\label{dir:eq:profiles}
 \max_{g\in G_{q_n}}\Bigl|\varphi(q_n)\dirP(Y_n=g)-\abs F\,P_\Lambda\bigl(\pi_n(g)b_n^{-h_n}\bigr)\Bigr|\longrightarrow0 .
\end{equation}
Conversely, every compound Poisson law on every finite abelian group $F$ with $\abs F\le D$ occurs in this
way, for some sequence of moduli and retained sets of density at least $a$, with $h_n=\ceil{B\log q_n}$.
\end{theorem}

The normalization in \eqref{dir:eq:profiles} matters: each residue probability is of order $1/\varphi(q)$,
and the theorem determines the relative profile, which is more than convergence of a projected law or an
unnormalized error tending to zero. The quotient is allowed to vary with $n$ before the subsequence is
taken; no single quotient is claimed to serve every retained set. By \eqref{dir:eq:cpfourier} a limiting
profile on a nontrivial $F$ is never uniform, and it may have limiting holes even when a few exceptional
primes make the exact support of \cref{dir:thm:quotient} larger.

\begin{proof}[Proof of the forward direction]
Apply \cref{dir:thm:quotient} at each $n$ to $Y_n$, which has the law of $X_{h_n}$, obtaining quotients
$\pi_n^0\colon G_{q_n}\to A_n$ with $\abs{A_n}\le D$ and remainder at most $q_n^{-2}$ in
\eqref{dir:eq:profile}. There are finitely many abelian groups of order at most $D$;
pass to a subsequence and identify every $A_n$ with a fixed group $A$ of order $d$. Choose a most frequent
value $a_{0,n}\in A$ of the projected primes, whose count is at least $m_n/d$; after a further subsequence
$a_{0,n}=a_0$ is fixed. For $x\in A\setminus\set e$ put
\[
 s_x^{(n)}=\#\set{p\in S_n:\pi_n^0(p)=a_0x},\qquad\lambda_x^{(n)}=\frac{h_ns_x^{(n)}}{m_n}.
\]
Since $[0,+\infty]$ is compact and $A$ is finite, pass to a subsequence along which every $\lambda_x^{(n)}$
has a limit $\lambda_x\in[0,+\infty]$, and let $K=\langle x\in A:\lambda_x=+\infty\rangle$.

If a character $\chi$ of $A$ is nontrivial on $K$, there is a generator $x$ with $\lambda_x=+\infty$ and
$\chi(x)\ne1$. The labels with values $a_0$ and $a_0x$ lie in distinct $\chi$-fibres, so at least
$\min(m_n/d,s_x^{(n)})$ labels lie outside a largest $\chi$-fibre, and after multiplication by $h_n/m_n$
this is at least $\min(h_n/d,\lambda_x^{(n)})\to\infty$. By \cref{dir:lem:splitpair} the coefficient of
$\chi\circ\pi_n^0$ tends to zero; normalizing by the unit phase $\chi(a_0)^{h_n}$ does not change its
modulus.

Now let $F=A/K$, let $\pi_n$ be the composite quotient and $b_n=a_0K$. For $y\in F\setminus\set e$ every
$x$ mapping to $y$ has finite $\lambda_x$, so the scaled number of $y$-defects converges to
$\lambda'_y=\sum_{xK=y}\lambda_x<\infty$. In $F$, after division of each label by $b_n$, the total scaled
defect count is bounded, so the proportion of defects is $O(1/h_n)$, and \cref{dir:cor:cp} gives
convergence of the normalized projected product to the compound Poisson law with rates $\lambda'_y$.

Every Fourier coefficient on $A$ is now determined: it tends to zero for characters nontrivial on $K$, and
to the compound Poisson coefficient for characters factoring through $F$. By Fourier inversion the law on
$A$, after normalization by $a_0^{h_n}$, converges uniformly to the lift of that law that is uniform on each
coset of $K$, whose mass at $\alpha\in A$ is $P_\Lambda(\alpha K)/\abs K$; multiplication by $d=\abs A$
gives $\abs FP_\Lambda(\alpha K)$. The arithmetic remainder of \cref{dir:thm:quotient} tends to zero
uniformly in $g\in G_{q_n}$, which gives \eqref{dir:eq:profiles}.
\end{proof}

So two regimes coexist in one bounded quotient: a defect type whose expected selected count diverges forces
uniformity on the subgroup it generates, and the remaining finite rates form a compound Poisson law on the
quotient by that subgroup. A limit therefore need not retain the quotient produced by the large-bias
argument.

\begin{proof}[Proof of the converse]
Let $F$ have order $d\le D$, and fix finite rates $\lambda_x\ge0$. Write $F$ as a product of cyclic groups
$C_{n_i}$, choose by Dirichlet's theorem distinct fixed primes $\ell_i\equiv1\pmod{n_i}$, and put
$q_0=\prod_i\ell_i$ ($q_0=1$ for the trivial group), so that the unit group modulo $q_0$ surjects onto $F$.
For primes $r\to\infty$ with $r\nmid q_0$ let $q=q_0r$, with the induced surjection $\pi_q\colon G\to F$.
Each fibre of $\pi_q$ contains a positive proportion of the $N$ pool primes, by the prime number theorem in
the fixed progressions modulo $q_0$, or by \cref{dir:lem:bias} with small fixed $\eps$ and the character
expansion of a fibre indicator. Choose a most frequent fibre $a_0$ and retain all of its $m_0\ge N/d$
primes; put $h=\ceil{B\log q}$ and retain in addition $s_x=\floor{\lambda_xm_0/h}$ primes from the fibre
$a_0x$ for each $x\ne e$, which is possible for large $q$ since $s_x=o(N)$. Then
\[
 m=m_0+\sum_xs_x\ge N/d\ge aN,\qquad h/m\to0,\qquad hs_x/m\to\lambda_x,
\]
and the projected profile converges to $P_\Lambda$ by \cref{dir:cor:cp}.

It remains to exclude other large modes. If $\chi$ does not factor through $\pi_q$, then
$\chi\cdot(\psi\circ\pi_q)$ is nonprincipal for every $\psi\in\widehat F$, and the fibre-indicator expansion
gives
\[
 \Bigl|\sum_{p\in\mathcal P,\ \pi_q(p)=a_0}\chi(p)\Bigr|\le\frac1d\sum_{\psi\in\widehat F}
 \Bigl|\sum_{p\in\mathcal P}\chi(p)\psi(\pi_q(p))\Bigr|\le N\eps .
\]
Divide by $m_0\ge N/d$ and add the $o(m_0)$ defects: for $\eps$ small in terms of $D$ every such $\chi$ has
bias on $S$ bounded away from one by a fixed amount. \Cref{dir:lem:circle}, summed over at most $q$
characters, makes their total contribution $O(q^{-A})$ after enlarging $B$, and Fourier inversion gives
the full normalized profile with $b_n=a_0$. This realizes every permitted profile, including the density
boundary $a=1/d$.
\end{proof}

\begin{example}[Four fibres]\label{dir:ex:fourfibre}
For $F=C_4=\langle z\rangle$ with the single positive rate $\lambda_z=\lambda$, the corrected product is
$z^J$ with $J$ Poisson of mean $\lambda$, and Fourier inversion gives
\begin{equation}\label{dir:eq:c4}
 4\dirP(J\equiv c\pmod4)=1+(-1)^ce^{-2\lambda}+2e^{-\lambda}\cos(\lambda-\pi c/2),\qquad c=0,1,2,3 .
\end{equation}
The four normalized fibre densities oscillate with damping; at $\lambda=1$ they are approximately
$1.532868$, $1.483784$, $0.737803$, $0.245545$ for $c=0,1,2,3$. Retained density $1/4$ permits this
profile, and the quadratic quotient alone sees only the term $e^{-2\lambda}$. For $F=C_2\times C_2$ with
rates $\lambda_{10},\lambda_{01},\lambda_{11}$, the normalized density at $(\sigma,\tau)\in\set{0,1}^2$ is
\begin{equation}\label{dir:eq:v4}
 4P_\Lambda(\sigma,\tau)=1+(-1)^\sigma e^{-2(\lambda_{10}+\lambda_{11})}+(-1)^\tau e^{-2(\lambda_{01}+\lambda_{11})}
 +(-1)^{\sigma+\tau}e^{-2(\lambda_{10}+\lambda_{01})} .
\end{equation}
The rate $\lambda_{11}$ couples the two quadratic signs, since one selected defect changes both; when
$\lambda_{10}=\lambda_{01}=0<\lambda_{11}$ the limit is supported on the diagonal subgroup although each
sign separately is nonconstant. \Cref{dir:cor:cp} gives the first correction for every Fourier mode at
once; for $C_2$ it is \eqref{dir:eq:paritylimit}. Exact rational arithmetic checks the finite identities
of \cref{dir:thm:fpc} in $2\,026$ population and quota cases against $77\,166$ enumerated subsets, and
$44\,280$ coefficient-ratio inequalities \eqref{dir:eq:ratios}; these checks do not cover the asymptotic
statements.
\dirfiles{\artifact{grh_extensions_2026-11-27_3/code/verify_prime_products.py},
\artifact{grh_extensions_2026-11-27_3/data/prime_product_verification.json}}
\end{example}

The relation between hypergeometric and binomial laws is classical; Ehm's Stein--Chen bounds include
sampling with and without replacement \cite{Ehm1991Binomial}, L\'opez-Bl\'azquez and Salamanca Mi\~no give
Krawtchouk expansions of hypergeometric probabilities \cite{LopezBlazquezSalamancaMino2000Binomial}, and
Darroch, Ji\v rina and Speed give signed approximations to the full law of a sample without replacement
\cite{DarrochJirinaSpeed1988Sampling}; the scalar correction mechanism is therefore not new. Poisson laws
for the prime factors of a random integer \cite{Ford2021Joint} concern a different model from a uniform
fixed-size sample from an arbitrarily retained pool of small primes. The constants $B,C$ are not optimized;
the theorems require $h\ge B\log q$ at the cutoff $C\log^2q$ and make no claim for shorter quotas, for
vanishing retained density, or uniformly as $a\to0$.

\subsection{Rainbow routing with small prime multipliers}\label{dir:sec:routing}
\sources{grh_prime_routing_2026-11-27}

In the preceding subsections the admissible primes form one retained set. Here the set of allowed
prime factors may depend on the residue reached so far: a route
$s\xrightarrow{p_1}sp_1\xrightarrow{p_2}\cdots\xrightarrow{p_k}t$ modulo $q$ uses $p_j$ only if it is
allowed at $s\prod_{i<j}p_i$, and supplies a squarefree $y$-smooth integer $P=p_1\cdots p_k$ with
$sP\equiv t\pmod q$. Fourier inversion over one retained set does not apply, since the intersection of
the allowed sets over all residues may be empty.

Let $q\ge3$, $L=\log q$, $G=(\Z/q\Z)^\times$, $n=\varphi(q)$, $\mathcal P$ the primes $p\le y$ not
dividing $q$, $d=\abs{\mathcal P}$, and $y<q$, so that distinct primes give distinct elements of $G$. The
directed graph $\mathcal G(q,y)$ has an arc $x\xrightarrow{p}xp$ for every $x\in G$ and $p\in\mathcal P$. For
$0<\delta<1/2$ a spanning subgraph $\mathcal H$, with no invariance required, is \emph{admissible} if every
vertex keeps at least $(1/2+\delta)d$ outgoing and at least $(1/2+\delta)d$ incoming arcs. A \emph{rainbow}
path is a simple directed path with distinct labels; a \emph{rainbow linkage} of terminal pairs is a family of
rainbow paths that are pairwise vertex-disjoint and use no label twice in total.

\begin{theorem}[Prime routing and packing; \stG]\label{dir:thm:routing}
Fix $A>2$ and $0<\delta<1/2$, and let $y=L^A$. For all large $q$, every admissible $\mathcal H$ has the
following properties.
\begin{enumerate}[label=\textup{(\roman*)}]
\item Every ordered pair of distinct vertices is joined by a rainbow path of length at most
$(2/(A-2)+o(1))L/\log L$.
\item Every family of $r$ disjoint ordered terminal pairs with
$r\le(\delta(A-2)/(64A)-o(1))L^{A-1}$ has a rainbow linkage whose paths obey the same length bound.
\item The largest $r\le n/2$ for which every admissible $\mathcal H$ and every $r$ disjoint pairs admit a
rainbow linkage, with no length restriction, is $\Theta_{A,\delta}(L^{A-1})$; it is at most
$(1+o(1))L^{A-1}$ already in the full graph.
\end{enumerate}
The $o(1)$ terms are independent of $\mathcal H$ and of the pairs.
\end{theorem}

\begin{corollary}[At the logarithm-squared cutoff; \stG]\label{dir:cor:routing}
For each $0<\delta<1/2$ there are $C_\delta,c_\delta,K_\delta>0$ such that with $y=C_\delta L^2$ any two
vertices of an admissible $\mathcal H$ are joined by a rainbow path of length at most $K_\delta L$, and any
$r\le c_\delta L/\log L$ disjoint pairs have a rainbow linkage with the same length bound.
\end{corollary}

At $y=L^A$ the integers attached to the routes are at most $q^{2A/(A-2)+o(1)}$, and those of a linkage
are pairwise coprime; at the logarithm-squared cutoff they are $q^{O_\delta(\log\log q)}$. The proof
separates a finite routing statement from the analytic input. Let $G$ be any finite abelian group of
order $n\ge2$, $P\subseteq G$ a set of $d$ distinct labels, $\mathcal H$ a spanning subgraph of the arcs
$x\to xp$ satisfying the same degree conditions, and $\beta=\max_{\chi\ne1}\abs{d^{-1}\sum_{p\in P}\chi(p)}$.

\begin{theorem}[Finite local-resilience lemma; \stU]\label{dir:thm:routingfinite}
Let $0<\delta<1/2$ and put
\[
 t=\sqrt{2d\log(16n)},\quad\tau=t/d,\quad\eps=2\beta+4\tau,\quad B=\Bigl\lfloor\frac{\delta d}{16}\Bigr\rfloor,
 \quad R=\Bigl(\frac\delta{2\eps}\Bigr)^2,\quad k=\Bigl\lfloor\frac{\log(n/2)}{\log R}\Bigr\rfloor+1 .
\]
Assume $\tau\le\delta/16$, $\eps<\delta/2$ and $k\le B$. There is one partition $P=P_+\sqcup P_-$, depending
only on $\mathcal H$, such that for all sets $F\subseteq P$ of forbidden labels and $W\subseteq G$ of
forbidden vertices with $\abs F+\abs W\le B$, and all $s,t_0\in G\setminus W$, there is a rainbow path from
$s$ to $t_0$ of length at most $2k$ avoiding $F$ and $W$, obtained from a walk of length at most $k$ with
labels in $P_+$ followed by one of length at most $k$ with labels in $P_-$. In particular every family of
$r$ disjoint ordered pairs with $(4k+2)r\le B$ has a rainbow linkage with paths of length at most $2k$.
\end{theorem}

For $Q\subseteq P$ let $E_Q(U,V)$ be the number of arcs of the full Cayley graph from $U$ to $V$ with label
in $Q$, $d_Q=\abs Q$ and $\beta_Q=\max_{\chi\ne1}\abs{d_Q^{-1}\sum_{p\in Q}\chi(p)}$.

\begin{lemma}[Directed expander mixing; \stU]\label{dir:lem:mix}
For all $U,V\subseteq G$, $\bigl|E_Q(U,V)/d_Q-\abs U\abs V/n\bigr|\le\beta_Q\sqrt{\abs U\abs V}$, and likewise for
reversed arcs.
\end{lemma}

\begin{proof}
The operator $f\mapsto d_Q^{-1}\sum_{p\in Q}f(\cdot\,p)$ is diagonal in the character basis, with norm at most
$\beta_Q$ on the complement of the constants, although its eigenvalues need not be real. Subtract the
constant components of $\mathbf 1_U$ and $\mathbf 1_V$ and apply Cauchy--Schwarz; the adjoint has the
conjugate eigenvalues.
\end{proof}

\begin{lemma}[Simultaneous split; \stU]\label{dir:lem:routesplit}
If $\tau\le\delta/16$, there is a partition $P=P_+\sqcup P_-$ with $(d-t)/2\le d_\pm=\abs{P_\pm}\le(d+t)/2$,
$\beta_{P_\pm}\le\eps$, and outgoing and incoming degrees in $\mathcal H$ restricted to labels in $P_\pm$ at
least $(1/2+7\delta/8)d_\pm$ at every vertex.
\end{lemma}

\begin{proof}
Give the labels independent random signs. Require that the signed sum of the indicator of each of the
$2n$ local label sets of $\mathcal H$, and of the real and imaginary parts of each of the $n$ characters,
have modulus at most $t$; each fails with probability at most $2e^{-t^2/(2d)}=1/(8n)$ by the
moment-generating bound $\cosh(x)\le e^{x^2/2}$, and there are at most $4n$ requirements. The principal
character gives the sizes. For $\chi\ne1$ the full sum is at most $\beta d$ and the signed sum at most
$\sqrt2t$, so $\beta_{P_\pm}\le(\beta+\sqrt2\tau)/(1-\tau)\le\eps$. Each local set has at least
$(1/2+\delta)d$ labels, hence at least $((1/2+\delta)d-t)/2$ in each pool, and
$(1/2+\delta-\tau)/(1+\tau)\ge1/2+7\delta/8$.
\end{proof}

\begin{proof}[Proof of \cref{dir:thm:routingfinite}]
Fix the partition of \cref{dir:lem:routesplit} and admissible $F,W$. Explore from $s$ with labels in $P_+$:
$U_0=\set s$, and $U_{i+1}$ adds, for each $x\in U_i$ with its stored witness walk, the endpoints of the
arcs of $\mathcal H$ from $x$ whose label lies in $P_+$, avoids $F$ and the witness's labels, and whose
endpoint avoids $W$, storing one witness per new endpoint. For $i<k$ each vertex loses at most $B+i\le2B$
of its at least $(1/2+7\delta/8)d_+$ arcs, and $d_+\ge d/3$ gives at least $(1/2+\delta/2)d_+$ eligible arcs,
all ending in $U_{i+1}$. \Cref{dir:lem:mix} gives
$(1/2+\delta/2)\abs{U_i}\le\abs{U_i}\abs{U_{i+1}}/n+\eps\sqrt{\abs{U_i}\abs{U_{i+1}}}$, so
$\abs{U_{i+1}}\ge R\abs{U_i}$ while $\abs{U_{i+1}}\le n/2$; hence $\abs{U_k}>n/2$. The same exploration
backwards from $t_0$ with incoming arcs and labels in $P_-$ gives a set of more than $n/2$ vertices
with reverse witnesses; excluding $W$ removes at most $\abs W$ incoming arcs per vertex since the
predecessors $xp^{-1}$ are distinct. The two witnesses meeting at a common vertex form a walk of length
at most $2k$ whose labels are distinct, the pools being disjoint; deleting closed subwalks gives a
rainbow path. For the linkage route the pairs one at a time with the same partition, forbidding the
labels already used and the internal vertices of earlier paths and all other terminals; then
$\abs F+\abs W\le(4k+2)r\le B$ at each stage.
\end{proof}

\begin{proof}[Proof of \cref{dir:thm:routing}\textup{(i)}, \textup{(ii)} and \cref{dir:cor:routing}]
Since $q$ has $O(L/\log L)$ prime factors, $\log n=(1+o(1))L$. For $y=L^A$, \cref{dir:lem:bias} gives
$d\sim L^A/(A\log L)$ and $\beta\le L^{1-A/2+o(1)}$, so $\tau=O(L^{(1-A)/2}\sqrt{\log L})$,
$\eps\le L^{1-A/2+o(1)}$ and $\log R\ge(A-2-o(1))\log L$; hence
$k\le(1/(A-2)+o(1))L/\log L=o(d)$, and
$\floor{B/(4k+2)}\ge(\delta(A-2)/(64A)-o(1))L^{A-1}$. At $y=C_\delta L^2$,
\cref{dir:lem:bias} with bias $\delta/64$ gives $\tau\to0$, $\eps\le\delta/8$, $R\ge16$, $k\le L/\log16+1$ and $B/(4k+2)\gg_\delta L/\log L$.
\end{proof}

\begin{proposition}[A counting obstruction to larger linkages; \stU]\label{dir:prop:routingupper}
Let $n\ge4$, $d\ge2$ and $h=1+\floor{\log(n/4)/\log d}$. If $r=\floor{d/h}+1$ satisfies $2r\le n/2$, some $r$
disjoint ordered pairs have no linkage with globally distinct labels, even in the full Cayley graph.
\end{proposition}

\begin{proof}
Fewer than $h$ steps reach at most $1+d+\dots+d^{h-1}\le n/2$ vertices from any source. Choose sources and
targets greedily, each target outside the short-distance ball of its source and away from the at most
$2r-1<n/2$ terminals already chosen. Each path then has at least $h$ labels, and $rh>d$.
\end{proof}

At $y=L^A$, $h\sim L/(A\log L)$ and $d/h\sim L^{A-1}$, and a failing family extends to larger sizes by
adding pairs; this is the upper bound in \cref{dir:thm:routing}(iii). The same count shows that the
worst-case length of one route is $\Omega_A(L/\log L)$ even in the full graph, so both orders are best
possible above the logarithm-squared scale. At $y=CL^2$ this count bounds the capacity only by $O_\delta(L)$.

\begin{proposition}[The deletion threshold is sharp; \stU]\label{dir:prop:routingsharp}
If $\beta=o(1)$, the prime Cayley graph can be disconnected while every vertex keeps $(1/2+o(1))d$
outgoing and incoming arcs. Keeping nearly all outgoing arcs alone does not guarantee reachability, and
neither does keeping nearly all incoming arcs alone.
\end{proposition}

\begin{proof}
Since $n$ is even there is a character $\chi$ of order two; keep the arcs with $\chi(p)=1$, which never
leave a coset of $\ker\chi$ and give every vertex in- and out-degree $(d+\sum_p\chi(p))/2=d/2+O(\beta d)$. At
$y=CL^2$, \cref{dir:lem:bias} makes this fraction arbitrarily close to $1/2$ by enlarging $C$. Deleting all
incoming arcs of one vertex removes at most one outgoing arc elsewhere but makes that vertex unreachable;
deleting outgoing arcs instead gives the converse.
\end{proof}

\begin{example}\label{dir:ex:routing}
Locally damaged graphs with vertex-dependent retained label sets and ten routed terminal pairs each,
checked by exact integer verification of every multiplication, retained arc and disjointness condition:
\begin{center}\small
\begin{tabular}{rrrrrr}
\toprule
$q$ & $n$ & $d$ & least in/out degree & routes & longest route\\
\midrule
$1009$ & $1008$ & $109$ & $82/81$ & $10$ & $2$\\
$1155$ & $480$ & $105$ & $78/76$ & $10$ & $2$\\
$5003$ & $5002$ & $154$ & $114/115$ & $10$ & $3$\\
\bottomrule
\end{tabular}
\end{center}
The finite inequalities of \cref{dir:thm:routingfinite} need not hold at these sizes.
\dirfiles{\artifact{grh_prime_routing_2026-11-27/examples.json}}
\end{example}

Small-prime Cayley expansion of $(\Z/q\Z)^\times$ under GRH is due to Jao, Miller and Venkatesan
\cite{JaoMillerVenkatesan2009Expander}; the results above add local resilience, rainbow paths and
linkages with globally distinct primes.

\subsection{Primes with a prescribed majority of quadratic signs}\label{dir:sec:signs}
\sources{hermite_thresholds_and_quadratic_primes; hermite_threshold_research; prime_cubature_grh_2026-11-27}

Let $q_1,\dots,q_k\le Y$ be distinct odd primes and prescribe a sign for each Legendre symbol
$\legendre p{q_i}$. How small a prime realizes more than half of the prescribed signs? Under
\cref{src:thm:dgrh}, a prime $p\ll k^3\log^2(2Y)$ realizes more than $k/2+\sqrt k/4$ of them,
and $p\ll k^5\log^2(2Y)$ realizes more than $k/2+3\sqrt k/4$. The powers of $k$ come from an
extremal problem on the cube with no arithmetic input, which is solved first, including the exact
lattice oscillation of the leading constant at the Hermite thresholds. Conversely, a rational law on the primes
of one dyadic interval can make $r$ quadratic signs exactly $k$-wise independent, robustly under
adversarial deletions, at a height $O_k(r^{2k}\log^2(2Y))$ (\cref{dir:thm:cubature}).

\paragraph{A missing Hamming ball.}
For $u\in\set{-1,1}^k$ put $X(u)=k^{-1/2}\sum_iu_i$ and $u_A=\prod_{i\in A}u_i$ for
$A\subseteq[k]$. The parity bias of a probability law $\mu$ on the cube is
$\epsilon(\mu)=\max_{A\ne\varnothing}\abs{\dirE_\mu u_A}$. For real $a$ let
\[
 \mathcal B_k(a)=\set{X\le a},\qquad\mathcal D_k(a)=\set{X>a},\qquad
 b_k(a)=\min_{\supp\mu\subseteq\mathcal B_k(a)}\epsilon(\mu),
\]
the least bias of a law that misses the cap $\mathcal D_k(a)$; a law of smaller bias meets every
translate of the cap by coordinate sign changes. Let $H_0=1$, $H_1=x$,
$H_{j+1}=xH_j-jH_{j-1}$ be the monic probabilists' Hermite polynomials, and $h_j$ the largest zero
of $H_j$ for $j\ge1$; thus $h_1,h_2,h_3,h_4=0,1,\sqrt3,\sqrt{3+\sqrt6}$.

\begin{lemma}[Sharp parity bias for a missing Hamming ball; \stU]\label{dir:lem:hamming}
\begin{enumerate}[label=\textup{(\alph*)}]
\item Fix $d\ge1$ and $h_d<a<h_{d+1}$, and put $m=2d+1$. Then
\[
 \lim_{k\to\infty}k^{m/2}b_k(a)=C_d(a),\qquad C_d(a)=-d!\,\frac{H_{d+1}(a)}{H_d(a)}>0 .
\]
\item Fix $j\ge1$ and $R>0$. Uniformly in integers $s\equiv k\pmod2$ with
$r=s-h_j\sqrt k\in[-R,0]$,
\[
 b_k(s/\sqrt k)=k^{-j}\bigl(j!(1-r)+O_{j,R}(k^{-1/2})\bigr),
\]
and this is attained to the same accuracy by a permutation-invariant law on at most $4j-2$ layers. In
particular, if $\delta_{k,j}\in[0,2)$ is the distance from $h_j\sqrt k$ down to the nearest integer of the
parity of $k$, then $b_k(h_j)=j!k^{-j}(1+\delta_{k,j}+O_j(k^{-1/2}))$; and for fixed $\lambda\le0$, with
$\delta_k\in[0,2)$ the distance from $h_j\sqrt k+\lambda$ down to the nearest integer of the parity of $k$,
\[
 k^jb_k\bigl(h_j+\lambda/\sqrt k\bigr)=j!(1-\lambda+\delta_k)+O_{j,\lambda}(k^{-1/2}),
\]
uniformly for $\lambda$ in compact subsets of $(-\infty,0]$.
\item For an integer $-k\le s\le0$ with $s\equiv k\pmod2$, the least bias of a law supported on
$\set{\sum_iu_i\le s}$ is $(1-s)/(k+1)$. In particular $b_k(0)=1/(k+1)$ for even $k$ and
$2/(k+1)$ for odd $k$, and $b_k(a)=\abs a/\sqrt k+O_a(1/k)$ for fixed $a<0$.
\end{enumerate}
\end{lemma}

For $0<a<1$ the constant is $1/a-a$, and for $1<a<\sqrt3$ it is $2a(3-a^2)/(a^2-1)$. Between two
consecutive roots $C_d(a)$ tends to zero at the upper root and to infinity at the lower one, and
at a root the exponent lies halfway between those of the neighbouring intervals. At a root there is no
single leading constant.

\begin{corollary}[All subsequential constants at a root; \stU]\label{dir:cor:hermiteconst}
For $j\ge2$ and fixed $\lambda\le0$ the empirical distribution of $k^jb_k(h_j+\lambda/\sqrt k)$, $1\le k\le N$,
converges as $N\to\infty$ to the uniform law on $[j!(1-\lambda),j!(3-\lambda)]$. In particular the
subsequential limits of $k^jb_k(h_j)$ fill $[j!,3j!]$; for $j=1$ the two parities of $k$ give the limits
$1$ and $2$.
\end{corollary}

\begin{proof}
For fixed parity $e$ and $k=2n+e$, $\delta_k=2\{(h_j\sqrt{2n+e}+\lambda-e)/2\}$, and for $h_j>0$ the
fractional parts of $a\sqrt{2n+e}+c$ are uniformly distributed, with discrepancy $O(N^{-1/2})$ by
inverting the inequalities over the $O(\sqrt N)$ blocks of consecutive integer parts. Apply part (b).
For $j=1$, $h_1=0$ and $\delta_k\in\set{0,1}$ according to parity.
\end{proof}

For $j=2$ the cluster interval is $[2,6]$: along $k=n^2$ and $k=n^2+2$ the constant tends to $2$, and
along $k=n^2-2$ to $6$; for $j=3,4$ the intervals are $[6,18]$ and $[24,72]$.

\emph{A quadrature rule with a prescribed extreme node.} Let $G$ be a standard normal variable.
The polynomials $H_j/\sqrt{j!}$ are orthonormal for $G$, by integration by parts from
$H_j(x)e^{-x^2/2}=(-1)^j(d/dx)^je^{-x^2/2}$, and multiplication by $x$ has off-diagonal entries
$\sqrt j$ between positions $j-1$ and $j$. Let $J_d(c)$ be its truncation to positions
$0,\dots,d$ with the last diagonal entry changed from $0$ to $c$, and $\nu_c$ the spectral law of
$J_d(c)$ at the first coordinate.

\begin{lemma}[A Gauss--Radau rule; \stU]\label{dir:lem:radau}
For $t>h_d$ put $c(t)=H_{d+1}(t)/H_d(t)$. The largest eigenvalue of $J_d(c(t))$ is $t$, all
nodes of $\nu_{c(t)}$ are simple with positive masses, the moments of $\nu_c$ through degree $2d$
are those of $G$, and its moment of degree $m=2d+1$ is $d!c$. For $d=0$ the rule is the unit mass
at $c$. As $t\to h_{d+1}$ the nodes and masses of $\nu_{c(t)}$ differ from those of Gaussian quadrature
$\nu_0$ by $O_d(\abs{t-h_{d+1}})$, and the mass $w_*$ of $\nu_0$ at $h_{d+1}$ satisfies
$w_*Q_{d+1}(h_{d+1})^2=(d+1)!$, where $Q_{d+1}(x)=H_{d+1}(x)/(x-h_{d+1})$.
\end{lemma}

\begin{proof}
The determinant recurrence gives the characteristic polynomial $H_{d+1}-cH_d$. An eigenvector
is determined by its first entry since all off-diagonal entries are nonzero, so eigenspaces are
lines and all masses are positive. The leading principal block of size $d$ has characteristic
polynomial $H_d$, and interlacing is strict because consecutive Hermite polynomials have no
common zero; so the eigenvalue $t>h_d$ is the largest, and $c(t)<0$ for $h_d<t<h_{d+1}$. A moment
at $e_0$ is the total weight of closed walks from position $0$ in the tridiagonal matrix; a walk
reaching the omitted position $d+1$ has length at least $2d+2$, and one using the changed
diagonal entry has length at least $2d+1$, with the single such walk of that length going up to
$d$, staying once and returning, of weight $d!c$. The odd Gaussian moment is zero. Near $c=0$ the
simple eigenvalues and their spectral projections depend smoothly on $c$, and $c(t)=O_d(\abs{t-h_{d+1}})$.
Finally $Q_{d+1}$ vanishes at every node of $\nu_0$ except $h_{d+1}$, quadrature is exact for
$Q_{d+1}^2$ of degree $2d$, and $\dirE Q_{d+1}(G)^2=(d+1)!$: from
$H_{d+1}'=(d+1)H_d=Q_{d+1}+(x-h_{d+1})Q_{d+1}'$, orthogonality and monicity give
$\dirE H_{d+1}'(G)Q_{d+1}(G)=(d+1)!$ and $\dirE H_{d+1}(G)Q_{d+1}'(G)=0$.
\end{proof}

\begin{proof}[Proof of \cref{dir:lem:hamming}\textup{(a)}]
\emph{Lower bound.} Let $y_1,\dots,y_d$ be the other nodes for $t=a$ and put
\[
 Q(x)=\frac{H_{d+1}(x)-c(a)H_d(x)}{x-a},\qquad F(x)=(x-a)Q(x)^2 .
\]
This monic polynomial of degree $m$ vanishes at every node and is nonpositive for $x\le a$, and
since its lower moments agree under $G$ and $\nu_{c(a)}$,
\begin{equation}\label{dir:eq:witness}
 \dirE F(G)=-d!\,c(a)=C_d(a);
\end{equation}
for example $a=1/2$ gives $F=(x-1/2)(x+2)^2$ and $\dirE F(G)=3/2$. Write $\mathcal U_k$ for the
uniform law on the cube, $\widehat g(A)=\dirE_{\mathcal U_k}[gu_A]$, $L(g)=\sum_A\abs{\widehat g(A)}$
and $L_*(g)=\sum_{A\ne\varnothing}\abs{\widehat g(A)}$. Then $L(gh)\le L(g)L(h)$ and
$L(X)=\sqrt k$, so
\begin{equation}\label{dir:eq:wcost}
 L(F(X))\le k^{m/2}+O_F\bigl(k^{(m-1)/2}\bigr);
\end{equation}
the degree-$m$ terms have coefficients $m!k^{-m/2}$ and total mass $(k)_mk^{-m/2}$, where
$(k)_m=k(k-1)\cdots(k-m+1)$, and products with repeated indices have mass $O_m(k^{m/2-1})$. For
fixed $j$, $\dirE_{\mathcal U_k}X^j=\dirE G^j+O_j(k^{-1})$: odd moments vanish, and in a moment
of order $2r$ the tuples of $r$ distinct pairs give $(2r-1)!!(k)_r/k^r$. Any law $\mu$ on
$\mathcal B_k(a)$ has $\dirE_\mu F(X)\le0$, so
\[
 C_d(a)+O_a(k^{-1})\le\dirE_{\mathcal U_k}F(X)-\dirE_\mu F(X)\le\epsilon(\mu)L_*(F(X)),
\]
and \eqref{dir:eq:wcost} gives $\liminf k^{m/2}b_k(a)\ge C_d(a)$.

\emph{Upper bound.} Fix $h_d<t_1<t_2<a$ and mix the two rules of \cref{dir:lem:radau} with largest
nodes $t_1,t_2$ in equal proportions. The $2d+2=m+1$ nodes are distinct, since the two
characteristic polynomials differ by a nonzero multiple of $H_d$, the masses are positive, the
moments through degree $m-1$ are Gaussian, and the moment of degree $m$ is
$M=\frac{d!}2(c(t_1)+c(t_2))$. Round every node to the nearest attainable value of $X$ on the
layers of even weight $w=\#\set{i:u_i=-1}$, and a second copy on the layers of odd weight; each
mesh is $4/\sqrt k$, and for large $k$ the rounded nodes remain distinct and below $a$. On each
list solve for the $m+1$ weights whose moments are $\dirE_{\mathcal U_k}X^j$ for $j<m$ and $M$ for
$j=m$. Both Vandermonde systems converge to the same invertible one, so the weights converge to
the positive original weights. Choose a parity with probability $1/2$, a layer by its weights,
and a uniform word in the layer; this gives a law $\mu_k$ on at most $4d+4$ layers.

Let $R_j(w)$ be the mean of a product of $j$ distinct coordinates on the layer of weight $w$, and
$s=k-2w=\sqrt kX$. Then
\begin{equation}\label{dir:eq:layer}
 \sum_{j=0}^k\binom kjR_j(w)z^j=(1-z)^w(1+z)^{k-w},\qquad R_{k-j}(w)=(-1)^wR_j(w),
\end{equation}
and multiplication by the coordinate sum gives $sR_j=jR_{j-1}+(k-j)R_{j+1}$, so $(k)_jR_j$ is
monic of degree $j$ in $s$. Hence on each parity the level-$j$ bias vanishes for $1\le j<m$ and
equals $Mk^{m/2}/(k)_m$ at $j=m$; as the two parities have the same moments through $m$, the
complementary levels $k-m,\dots,k$ vanish in the mixture. For the middle levels, orthogonality of
the Walsh characters gives $\sum_w2^{-k}\binom kwR_j(w)^2=\binom kj^{-1}$, and on the chosen
layers $\abs X$ is bounded, so $2^{-k}\binom kw\gg k^{-1/2}$ and, uniformly for
$m+1\le j\le k-m-1$, $\abs{R_j(w)}\ll k^{1/4}\binom k{m+1}^{-1/2}=O(k^{-m/2-1/4})$. So
$\limsup k^{m/2}b_k(a)\le\abs M$, and letting $t_1,t_2\to a$ gives $\abs M\to C_d(a)$.
\end{proof}

\begin{proof}[Proof of \cref{dir:lem:hamming}\textup{(b)}]
Put $h=h_j$, $J=j!$, $Q=H_j/(x-h)$, monic of degree $j-1$, and $F=(x-h)Q^2=H_jQ$. Then $\dirE F(G)=0$ by
orthogonality and $\dirE Q(G)^2=J$ by \cref{dir:lem:radau} with $d=j-1$.

\emph{Lower bound: a parity-sensitive witness.} Let $w_0=(k-s)/2$ and
$\eta(u)=(-1)^{w-w_0}=(-1)^{w_0}\prod_iu_i$. On the permitted set $S=\sum_iu_i\le s$ one has
$S\le s-1+\eta(u)$, since $S\equiv s\pmod 4$ fails exactly when $\eta=-1$, and then $S\le s-2$. As
$h=(s-r)/\sqrt k$, the function
\[
 g(u)=\Bigl(X-h+\frac{1-r-\eta(u)}{\sqrt k}\Bigr)Q(X)^2=\frac{S-s+1-\eta(u)}{\sqrt k}\,Q(X)^2
\]
is nonpositive on the permitted set. For $k>2j-2$ the uniform mean of $\eta Q(X)^2$ vanishes, because
$Q(X)^2$ has Walsh degree at most $2j-2$, so the comparison of cube and Gaussian moments gives
$\dirE_{\mathcal U_k}g=J(1-r)/\sqrt k+O_{j,R}(k^{-1})$. Since $L(X)=\sqrt k$, $L(\eta)=1$ and $F$ is monic,
$L_*(g)\le k^{j-1/2}+O_{j,R}(k^{j-1})$. As in (a), every law on the permitted set has bias at least
$\dirE_{\mathcal U_k}g/L_*(g)$, which is the lower bound with error $O_{j,R}(k^{-j-1/2})$.

\emph{Upper bound: two exact top nodes.} Split the permitted layers into the classes $\eta=\pm1$, with
largest nodes $a_+=s/\sqrt k=h+r/\sqrt k$ and $a_-=(s-2)/\sqrt k$; put $r_+=r$ and $r_-=r-2$. In each class
take the rule $\nu_{c(a_\pm)}$ of \cref{dir:lem:radau} with $d=j-1$ and largest node $a_\pm$. Keep the top
node; split every other node between two grid points of the class at distances in
$[4/\sqrt k,8/\sqrt k]$ by linear interpolation, each fraction being at least $1/3$, so that the moments
through degree $2j-2$ change by $O_{j,R}(k^{-1})$. On the $2j-1$ resulting nodes correct the weights to
match the cube moments of degrees $0,\dots,2j-2$ exactly: the inverse of the square Vandermonde matrix
has entries $O_{j,R}(\sqrt k)$, each Lagrange denominator having one factor of order $k^{-1/2}$, so the
weights change by $O_{j,R}(k^{-1/2})$ and stay positive. Spread each weight uniformly on its layer to get
$\mu_\pm$, and put $\mu=(\mu_++\mu_-)/2$.

With $R_\ell(w)$ as in \eqref{dir:eq:layer}, $W_\ell=(k)_\ell k^{-\ell/2}R_\ell$ is a monic polynomial of degree
$\ell$ in $X$ with $W_{\ell+1}=XW_\ell-\ell(1-(\ell-1)/k)W_{\ell-1}$, so $W_\ell=H_\ell+O_\ell(k^{-1})$
coefficientwise and $\dirE_{\mathcal U_k}W_\ell=0$ for $\ell\ge1$. Levels $1\le\ell\le2j-2$ vanish on both
classes. Since $W_{2j-1}-F$ has degree at most $2j-2$, and $F$ has double zeros at the interior nodes of
$\nu_0$ while the top node contributes $(r_\pm/\sqrt k)(w_*Q(h)^2+O(k^{-1/2}))$,
\[
 \dirE_{\mu_\pm}W_{2j-1}=\dirE_{\mu_\pm}F-\dirE_{\mathcal U_k}F=\frac{Jr_\pm}{\sqrt k}+O_{j,R}(k^{-1}).
\]
Since $H_{2j}-H_j^2$ has degree at most $2j-2$ and Gaussian mean $-J$, and $H_j^2$ vanishes on the nodes
of $\nu_0$, $\dirE_{\mu_\pm}W_{2j}=-J+O_{j,R}(k^{-1/2})$. By the reflection in \eqref{dir:eq:layer}, with
$(-1)^w=(-1)^{w_0}\eta$, $k^j$ times the bias of $\mu$ at the levels $2j-1$, $k-2j+1$, $2j$ and $k-2j$ is
$J(r-1)$, $(-1)^{w_0}J$, $-J$ and $0$, each up to $O_{j,R}(k^{-1/2})$. The top level vanishes since the
classes have equal mass, the low levels and their complements vanish, and for
$2j+1\le\ell\le k-2j-1$ the middle-level bound in the proof of (a) gives $O_j(k^{-j-1/4})$, all nodes lying
in a fixed compact interval. As $r\le0$ the largest of these is $Jk^{-j}(1-r)+O(k^{-j-1/2})$. The window
formula follows by substituting $r=\lambda-\delta_k$.
\end{proof}

\begin{proof}[Proof of \cref{dir:lem:hamming}\textup{(c)}]
Let $S=\sum_iu_i$ and $w_0=(k-s)/2$ be the top allowed layer; every allowed word has
$S\le s-1+(-1)^{w-w_0}$. If the bias is $\epsilon$, then $-k\epsilon\le\dirE S\le s-1+\epsilon$,
which gives the lower bound. For the upper bound put $N=k+1$ and $t=1-s$, take a uniform word of
the $N$-cube with coordinate sum $-t$ and discard its last coordinate; the remaining sum lies in
$\set{s,s-2}$. On such a slice, with $t\ge1$, every product of $j$ coordinates, $1\le j\le N-1$,
has mean of absolute value at most $t/N$. For $j=1$ this is the average of the sum and for
$j=N-1$ it follows by complementation. For $t\ge2$ and $2\le j\le N-2$, condition on whether a
fixed negative coordinate is sampled: the rest is uniform in a population of size $N-1$ and
imbalance $t-1$, and induction gives the bound $(t-1)/(N-1)\le t/N$. For $t=1$, $N=2l+1$, the
generating polynomial \eqref{dir:eq:layer} is $(1-z)(1-z^2)^l$, and by complementation it
suffices to consider even degrees $2i$, $1\le i\le l$, where the mean is
$\binom li/\binom{2l+1}{2i}$; its reciprocal equals $N$ at $i=1,l$ and the ratio of consecutive
reciprocals, $(2l-2i+1)/(2i+1)$, decreases in $i$, so the reciprocal is smallest at an endpoint.
This is a Krawtchouk eigenvalue bound \cite[Theorem~1.2(i)]{BrouwerCioabaIhringerMcGinnis2017Smallest}.
Taking $s=0$ or $s=-1$, and for fixed $a<0$ the largest attainable $s\le a\sqrt k$, gives the
remaining assertions.
\end{proof}

The extremal problem has a dual description. For a nonempty proper subset $D$ of the cube let
$\phi(D)$ be Bella\"iche's optimized Littlewood complexity in the abelian case
\cite{Bellaiche2016Chebotarev}: the infimum of $L(g)/\dirE_{\mathcal U_k}g$ over real $g$ with
positive mean and $g\le0$ off $D$.

\begin{proposition}[Duality; \stU]\label{dir:prop:hammingdual}
If $\mathcal D_k(a)$ is nonempty and proper, then $\phi(\mathcal D_k(a))=1+1/b_k(a)$.
\end{proposition}

\begin{proof}
With $v(u)=(u_A)_{A\ne\varnothing}$, $b_k(a)$ is the $\ell^\infty$ distance of $0$ from the convex
hull of the $v(u)$ with $u\in\mathcal B_k(a)$, so finite-dimensional minimax gives
$b_k(a)=\max_{\norm z_1\le1}\min_{u\in\mathcal B_k(a)}z\cdot v(u)$, which is positive because a law
with all nonconstant coefficients zero is uniform. An optimal $z$ gives $g=b_k(a)-z\cdot v$.
Conversely, if $g$ has mean $c>0$ and nonconstant coefficients $s$, then $-s\cdot v\ge c$ off
$\mathcal D_k(a)$, and the minimax identity gives $c\le b_k(a)\norm s_1$.
\end{proof}

So the degree-$m$ polynomial $F$ is asymptotically optimal among all real functions that detect
the cap, and in the setting of \cref{dir:lem:hamming}(b) the complexity of the cap $\set{S>s}$ is
$1+k^j(1+O_{j,R}(k^{-1/2}))/(j!(1-r))$, attained asymptotically by the parity-sensitive witness of its
proof. Near-optimal laws are also rigid.

\begin{proposition}[Rigidity of near minimizers; \stU]\label{dir:prop:hammingrigid}
Let $h_d<a<h_{d+1}$, $m=2d+1$ and $C=C_d(a)$. If $\supp\mu_k\subseteq\mathcal B_k(a)$ and
$k^{m/2}\epsilon(\mu_k)\le C+o(1)$, then the law of $X$ under $\mu_k$ converges weakly to
$\nu_{c(a)}$, and
\[
 \binom km^{-1}\sum_{\abs A=m}\bigl|k^{m/2}\dirE_{\mu_k}u_A+C\bigr|\longrightarrow0 .
\]
No coordinate symmetry of the laws is assumed.
\end{proposition}

\begin{proof}
By \eqref{dir:eq:witness} and \eqref{dir:eq:wcost},
$0\le\dirE_{\mu_k}[-F(X)]\le\epsilon(\mu_k)L_*(F(X))-\dirE_{\mathcal U_k}F(X)=o(1)$. On $(-\infty,a]$
the function $-F\ge0$ vanishes exactly at the $d+1$ nodes of $\nu_{c(a)}$ and grows like
$\abs x^m$, so the mass away from the nodes tends to zero, the laws are tight, and their absolute
moments of order $m$ are bounded. For $j<m$,
$\abs{\dirE_{\mu_k}X^j-\dirE_{\mathcal U_k}X^j}\le\epsilon(\mu_k)k^{j/2}=o(1)$, so every subsequential
limit lies on the nodes with Gaussian moments of orders $0,\dots,d$, which determine the weights
of \cref{dir:lem:radau}. In $F(X)$ the nonconstant Fourier mass below degree $m$ is
$O_F(k^{(m-1)/2})$ and contributes $o(1)$; the top coefficients all equal $m!k^{-m/2}$, and
$\dirE_{\mu_k}F(X)-\dirE_{\mathcal U_k}F(X)\to-C$, so the mean of $k^{m/2}\dirE_{\mu_k}u_A$ over
$\abs A=m$ tends to $-C$, while each term is at least $-C-o(1)$.
\end{proof}

For $0<a<1$ the limiting law has mass $1/(1+a^2)$ at $a$ and $a^2/(1+a^2)$ at $-1/a$; at $a=1/2$
four fifths of the mass lies at $1/2$, and $k^{3/2}$ times the mean three-coordinate bias tends to
$-3/2$, even when individual words in one layer carry different weights.

\begin{proposition}[Top-layer rigidity at a root; \stU]\label{dir:prop:hammingrigidroot}
In the setting of \cref{dir:lem:hamming}(b), let $\mu_k$ be arbitrary laws on $\set{S\le s}$ with
$k^j\epsilon(\mu_k)\le j!(1-r)+o(1)$, and let $w_*$ be the mass of Gaussian quadrature $\nu_0$ at $h_j$.
Then each parity class $\eta=\pm1$ has mass $1/2+o(1)$, the conditional law of $X$ on each class tends to
$\nu_0$, and $\mu_k(S=s)\to w_*/2$ and $\mu_k(S=s-2)\to w_*/2$. No coordinate symmetry is assumed.
\end{proposition}

\begin{proof}
For the witness $g$ of the proof of (b), $0\le\dirE_{\mu_k}[-g]\le\epsilon(\mu_k)L_*(g)-\dirE_{\mathcal U_k}g=o(k^{-1/2})$.
For $t\le2j$, $L(X^t)\le k^{t/2}$ and $L(\eta X^t)\le k^{t/2}$ with zero uniform mean for $t<k$, so on
each class, whose mass tends to $1/2$ by the bias bound for $\eta$, the moments of $X$ through degree
$2j-1$ tend to the Gaussian ones and the moment of degree $2j$ stays bounded. On compact sets $g=F+O(k^{-1/2})$,
and $-F\ge0$ on $(-\infty,h_j]$ vanishes only at the zeros of $H_j$, so every subsequential limit of
either class is supported on these zeros with Gaussian moments of orders below $j$, hence equals
$\nu_0$. Finally $\sqrt k(-g)=(s-1+\eta-S)Q(X)^2$, where $Q(X)^2$ is bounded below near $h_j$ and the
first factor is at least four off the top layer of the class; so all but $o(1)$ of the limiting mass
$w_*$ near $h_j$ lies on that layer.
\end{proof}

\paragraph{The prime application.}
For a sign word $v\in\set{-1,1}^k$ and a prime $p\nmid q_1\cdots q_k$ let
$M_v(p)=\#\set{i:\legendre p{q_i}=v_i}$, and for fixed $a$ with $h_d<a<h_{d+1}$, $d\ge1$, put
$D(a)=2d+1$. The roots and the windows to their left are treated in \cref{dir:thm:signswindow}.

\begin{theorem}[Primes realizing most prescribed quadratic signs; \stG]\label{dir:thm:signs}
For fixed $a>0$ that is not a zero of any $H_j$, every sufficiently large $k$, all distinct odd primes $q_1,\dots,q_k\le Y$ and
every $v\in\set{-1,1}^k$, there is a prime $p\nmid\prod_iq_i$ with
\[
 M_v(p)>\frac{k+a\sqrt k}2,\qquad p\ll_ak^{D(a)}\log^2(2Y),
\]
uniformly in the $q_i$ and in $v$.
\end{theorem}

\begin{proof}
Let $D=D(a)$ and $\mathcal Q=\prod_iq_i$. Take the witness $F(X)$ of \eqref{dir:eq:witness} and
translate it by $u_i\mapsto v_iu_i$ to obtain $g$; translation preserves the uniform mean and the absolute
Fourier coefficients. So $g$ has mean $g_0\gg_a1$, $L(g)\ll_ak^{D/2}$ and degree at most $D$, and
$g(u)>0$ forces $\sum_iv_iu_i>a\sqrt k$. For $A\subseteq[k]$ the character
$\chi_A=\prod_{i\in A}\legendre\cdot{q_i}$ is primitive of conductor $f_A=\prod_{i\in A}q_i$, by
the Chinese remainder theorem, so $\log(2f_A)\le D\log(2Y)$ for $\abs A\le D$. The smoothed
explicit formula \eqref{rec:eq:triangle} of \cref{rec:lem:triangle}, with kernel $x^s/(s(s+1))$ (the degree-one
case of \cite[Lemma~6]{Bellaiche2016Chebotarev}), gives for $x\ge3$
\[
 \sum_{p<x,\ p\nmid\mathcal Q}\Bigl(1-\frac px\Bigr)\log p\,\chi_A(p)
 =\mathbf 1_{A=\varnothing}\frac x2+O\bigl(\sqrt x\log(2f_A)+\log\operatorname{rad}\mathcal Q\bigr):
\]
by \cref{src:thm:dgrh} and $\sum_\rho\abs{\rho(\rho+1)}^{-1}\ll\log(2f_A)$ the zero sum is
$O(\sqrt x\log(2f_A))$, higher prime powers contribute
$\sum_{j\ge2}\vartheta(x^{1/j})\ll\sqrt x$, and removing the primes dividing $\mathcal Q$ costs at
most $\log\operatorname{rad}\mathcal Q\le k\log Y$. Multiplying by $\widehat g(A)$ and summing gives
$\frac12xg_0+O_a\bigl(k^{D/2}[\sqrt x\log(2Y)+k\log(2Y)]\bigr)$. With
$x=K_ak^D\log^2(2Y)$ and $K_a$ large, the main term dominates since $D\ge2$, and a positive
smoothed sum contains a prime with $g>0$.
\end{proof}

At a root the parity-sensitive witness has a high-degree part, $-\eta Q(X)^2/\sqrt k$ after translation,
involving characters of conductor as large as $\mathcal Q$; its Fourier mass is small enough to keep the
exponent.

\begin{theorem}[Critical-window prime patterns; \stG]\label{dir:thm:signswindow}
Fix $j\ge1$ and $R>0$. For every sufficiently large $k$, all distinct odd primes $q_1,\dots,q_k\le Y$,
every $v\in\set{-1,1}^k$ and every integer $s\equiv k\pmod2$ with $r=s-h_j\sqrt k\in[-R,0]$, there is a
prime $p\nmid\prod_iq_i$ with
\[
 M_v(p)>\frac{k+s}2,\qquad p\ll_{j,R}\frac{k^{2j}\log^2(2Y)}{(1-r)^2}.
\]
In particular, for fixed $\lambda\le0$ there is such a prime with $M_v(p)>(k+h_j\sqrt k+\lambda)/2$ and
$p\ll_{j,\lambda}k^{2j}\log^2(2Y)$; for $\lambda=0$ this is the threshold $a=h_j$.
\end{theorem}

\begin{proof}
Take the witness $g$ of the proof of \cref{dir:lem:hamming}(b), translated by the signs $v_i$. Its part
of degree at most $2j-1$ has Fourier mass $O_{j,R}(k^{j-1/2})$ and conductors with $\log(2f_A)\ll_j\log(2Y)$,
and the rest, $-\eta Q(X)^2/\sqrt k$, has Fourier mass $O_j(k^{j-3/2})$ and $\log(2f_A)\le k\log(2Y)$; so
$\sum_A\abs{\widehat g(A)}\log(2f_A)\ll_{j,R}k^{j-1/2}\log(2Y)$. The smoothed formula in the proof of
\cref{dir:thm:signs} gives
\[
 \sum_{p<x,\ p\nmid\mathcal Q}\Bigl(1-\frac px\Bigr)\log p\,g(u(p))=\frac x2\dirE_{\mathcal U_k}g
 +O_{j,R}\bigl(k^{j-1/2}(\sqrt x+k)\log(2Y)\bigr),
\]
and $\dirE_{\mathcal U_k}g\ge j!(1-r)/(2\sqrt k)$ for large $k$. With
$x=C_{j,R}k^{2j}\log^2(2Y)/(1-r)^2$ the sum is positive, so some prime has $g(u(p))>0$, that is $S>s$. Since
$S\equiv k\pmod2$, in the window exceeding $s_k$ is the same as exceeding $h_j\sqrt k+\lambda$.
\end{proof}

The exponent $2j$ matches the optimized unweighted Littlewood complexity of the cap; no arithmetic
lower bound for the least prime is asserted.

The argument uses only the characters occurring in the witness. For primitive quadratic
characters of conductors $f_i\le Y$, with all primes dividing a multiple $\mathcal Q$ of their
conductors excluded, the same proof gives a prime with $p\ll_ak^D\log^2(2Y)+k^{D/2}\log
\operatorname{rad}\mathcal Q$ whenever every nonempty product of at most $D$ of the characters is
nonprincipal; relations among more than $D$ of them are allowed. \Cref{dir:lem:hamming} shows what
a uniform bound on individual parity biases can give; a lower bound for the least prime for a
worst choice of moduli would need an arithmetic construction, and the constants in
\cref{dir:thm:signs} are not evaluated.

\begin{example}\label{dir:ex:signs}
For the twelve odd primes $3,\dots,41$ and all $4096$ sign words, the largest distance from a
target word to the nearest word of Legendre symbols of a prime up to $x$ first falls to $5$ at
$x=79$, to $4$ at $229$ and to $3$ at $691$; so at least seven, respectively nine, matching signs
are guaranteed exactly from $79$, respectively $691$, on. Exact values certified by rational
primal and dual certificates are
\begin{center}\small
\begin{tabular}{rrrl}
\toprule
$k$ & $a$ & largest distance in $\mathcal D_k(a)$ & $b_k(a)$\\
\midrule
$8$ & $1/2$ & $3$ & $1/9$\\
$8$ & $3/2$ & $1$ & $1/69$\\
$12$ & $1/2$ & $5$ & $1/13$\\
$12$ & $3/2$ & $3$ & $1/104$\\
$17$ & $1/2$ & $7$ & $1/33$\\
$17$ & $3/2$ & $5$ & $53/13717$\\
\bottomrule
\end{tabular}
\end{center}
and an eight-layer construction with every Fourier level bounded at $k=10^{16}$ gives
$1.49999994<k^{3/2}b_k(1/2)<1.50100041$, with its largest bias at degree three. The exact
equatorial formula of \cref{dir:lem:hamming}(c) was checked for $1088$ balls. At the roots $h_2,h_3$
exact rational primal and dual certificates bracket $k^jb_k(s/\sqrt k)$ around the prediction
$j!(1-r)$ of \cref{dir:lem:hamming}(b): for $j=2$ by $[1.98437,2.26482]$ around $2$ at $k=16384$, and by
$[5.88412,6.31162]$ around $5.98438$ at $k=16382$; for $j=3$ by $[15.91509,17.94418]$ around $16.21502$
at $k=16384$.
\dirfiles{\artifact{hermite_thresholds_and_quadratic_primes/verification.txt},
\artifact{hermite_thresholds_and_quadratic_primes/radau_output.txt},
\artifact{hermite_threshold_research/checks.txt}}
\end{example}

The dual small-bias problem appears for read-once formulas in \cite{BazziNahas2013Small},
small-bias laws on few central layers in \cite{DerksenIvanovLeeViola2025Pseudorandomness}, and
integers with prescribed Legendre symbols in \cite{HansonVaughanZhang2016Least}.

\paragraph{Exactly independent prime laws after erasures.}
A single prime realizes a sign pattern only approximately; a probability law on the primes of one dyadic
interval can make the quadratic signs exactly $k$-wise independent, with weights comparable to a smooth
baseline on every prime that survives an adversarial deletion, at a height polynomial in the number of
signs for each fixed $k$. Here the number of signs is $r$ and the degree is $k$. Let $q_1,\dots,q_r\le Y$ be
distinct odd primes, $1\le k\le r$, and put
\[
 V(n,k)=\sum_{j=0}^k\binom nj\quad(0\le k\le n),\quad V(n,-1)=0,\qquad
 b(t)=\begin{cases}30(t-1)^2(2-t)^2,&1<t<2,\\0,&\text{otherwise},\end{cases}
\]
so that $b\ge0$, $\int b=1$ and $\norm b_\infty=15/8$. For an integer $X\ge2$ let $\mathcal P_X$ be the set of
primes $X<p<2X$ other than the $q_i$, $u(p)=(\legendre{q_1}p,\dots,\legendre{q_r}p)$ and
$\mu(p)=b(p/X)/\sum_{l\in\mathcal P_X}b(l/X)$; for $S\subseteq\mathcal P_X$ of positive mass put
$\mu_S=\mu|_S/\mu(S)$. Let $\mathcal U$ be the uniform law on $\set{-1,1}^r$ and $u_A=\prod_{i\in A}u_i$,
$u_\varnothing=1$.

\begin{theorem}[Prime cubature after erasures; \stG]\label{dir:thm:cubature}
For each $k\ge1$ there is an effective $C_k>0$ such that for all $r\ge k$, all distinct odd primes
$q_1,\dots,q_r\le Y$, and all
\[
 0\le\delta\le c_k:=\frac1{1024\cdot9^k},\qquad X\ge C_k\bigl(V+\delta V^2\bigr)^2\log^2(2Y),\qquad V=V(r,k),
\]
the set $\mathcal P_X$ is nonempty and, simultaneously for every $S\subseteq\mathcal P_X$ with
$\mu(\mathcal P_X\setminus S)\le\delta$, there are positive rational weights $w_p$ on $S$ with
\begin{equation}\label{dir:eq:cubature}
 \tfrac34\mu_S(p)\le w_p\le\tfrac54\mu_S(p),\qquad
 \sum_{p\in S}w_pu_A(p)=\begin{cases}1,&A=\varnothing,\\0,&1\le\abs A\le k.\end{cases}
\end{equation}
In particular, for every $B\subseteq[r]$ with $\abs B\le k$ and every $v\in\set{-1,1}^B$, the primes $p\in S$ with
$\legendre{q_i}p=v_i$ for $i\in B$ have total weight exactly $2^{-\abs B}$. The same holds for the signs
$\legendre p{q_i}$, at a cutoff of the same order.
\end{theorem}

For fixed $k$ the height is $O_k(r^{2k}\log^2(2Y))$ without deletions and $O_k(r^{4k}\log^2(2Y))$ at a fixed
admissible deletion fraction, the crossover being at $\delta\asymp1/V$; finding a prime in every full sign
cell would need a height exponential in $r$. The constant $C_k$ is not evaluated.

\begin{corollary}[Sparse designs, entropy and deletion by count; \stG]\label{dir:cor:cubature}
Under the hypotheses of \cref{dir:thm:cubature}:
\begin{enumerate}[label=\textup{(\alph*)}]
\item some exact rational design, that is nonnegative weights satisfying the moment equations of
\eqref{dir:eq:cubature}, is supported on at most $V(r,k)$ surviving primes, and every exact design has at least
$V(r,\floor{k/2})$ distinct sign vectors in its support \textup{(}\stU\textup{)};
\item the weights of \eqref{dir:eq:cubature} satisfy $\max_pw_p\ll\log X/X$, so their min-entropy is at least
$\log_2(X/\log X)-O(1)$;
\item there is an absolute $K\ge1$ such that the theorem holds for the deletion of any $\tau\abs{\mathcal P_X}$
primes, chosen adversarially, with $\tau\le c_k/K$, on replacing $\delta$ in the cutoff by $K\tau$.
\end{enumerate}
\end{corollary}

The arithmetic enters only through a small truncated parity bias, and the rest is a finite theorem. For a
positive probability law $\mu$ on a finite labelled set $\Omega$ with a map $u\colon\Omega\to\set{-1,1}^r$
(repeated values allowed) put $\epsilon=\max_{1\le\abs A\le\min(4k,r)}\abs{\dirE_\mu u_A}$.

\begin{lemma}[Three norm bounds; \stU]\label{dir:lem:cubnorms}
If $f=\sum_{\abs A\le k}a_Au_A$ is real with $\dirE_{\mathcal U}f^2=1$ and $V=V(r,k)$, then
\[
 \abs{\dirE_\mu f-\dirE_{\mathcal U}f}\le\epsilon\sqrt V,\qquad 1-\epsilon V\le\dirE_\mu f^2\le1+\epsilon V,\qquad
 \dirE_\mu f^4\le9^k+\epsilon V^2 .
\]
\end{lemma}

\begin{proof}
The $u_A$ are orthonormal for $\mathcal U$, so $\sum_Aa_A^2=1$, and with $L(f)=\sum_A\abs{a_A}$ one has
$L(f)\le\sqrt V$ and $L(fg)\le L(f)L(g)$. For a polynomial $g$ of degree at most $4k$,
$\abs{\dirE_\mu g-\dirE_{\mathcal U}g}\le\epsilon L(g)$, which gives the first two bounds and bounds the
fourth-moment discrepancy by $\epsilon V^2$. Finally $(\dirE_{\mathcal U}f^4)^{1/2}\le\sum_A3^{\abs A}a_A^2$ by
induction on the coordinates: writing $f=g+u_rh$ with weighted coefficient sums $G,H'$ for $g,h$, the odd
terms in $u_r$ vanish and
$\dirE f^4=\dirE g^4+6\dirE g^2h^2+\dirE h^4\le G^2+6GH'+H'^2\le(G+3H')^2$ by Cauchy--Schwarz. As the degree
is at most $k$, $\dirE_{\mathcal U}f^4\le9^k$.
\end{proof}

\begin{lemma}[Box criterion; \stU]\label{dir:lem:box}
Let $v_x=(u_A(x))_{\abs A\le k}\in\R^{V}$ and $e_0=(1,0,\dots,0)$. For a probability law $\nu$ on $S$ and
$0<\eta<1$ there are weights $w_x\in[(1-\eta)\nu(x),(1+\eta)\nu(x)]$ with $\sum_xw_xv_x=e_0$ if and only if
$\abs{\dirE_{\mathcal U}f-\dirE_\nu f}\le\eta\dirE_\nu\abs f$ for every real $f$ of degree at most $k$.
\end{lemma}

\begin{proof}
The correction $e_0-\sum_x\nu(x)v_x$ must lie in the compact convex set
$Z=\set{\sum_xt_x\nu(x)v_x:\abs{t_x}\le\eta}$, whose support function in the direction $a$ is
$\eta\sum_x\nu(x)\abs{a\cdot v_x}$. With $f=a\cdot v$ one has $a\cdot e_0=\dirE_{\mathcal U}f$ and
$a\cdot\sum_x\nu(x)v_x=\dirE_\nu f$, so the separation criterion in the directions $\pm a$ is the stated
inequality. The constant coordinate makes the total weight one.
\end{proof}

\begin{theorem}[Bounded-density repair; \stU]\label{dir:thm:cubrepair}
Let $V=V(r,k)$ and $H=9^k+\epsilon V^2$. If
\[
 \epsilon\le\frac1{32\cdot3^kV}\qquad\text{and}\qquad\mu(\Omega\setminus S)\le\frac1{512H},
\]
then some probability law $w$ on $S$ satisfies \eqref{dir:eq:cubature} with points of $\Omega$ in place of
primes, and $w$ can be taken rational when $\mu$ is. It suffices that
$\mu(\Omega\setminus S)\le\delta$ with $\delta\le1/(1024\cdot9^k)$, $\epsilon\le1/(32\cdot3^kV)$ and
$\delta\epsilon V^2\le1/1024$.
\end{theorem}

\begin{proof}
Put $R=\Omega\setminus S$ and $d=\mu(R)$. By \cref{dir:lem:box} with $\nu=\mu_S$ and $\eta=1/4$ it suffices,
by homogeneity, to treat $\norm f_{L^2(\mathcal U)}=1$. With $g=f-\dirE_{\mathcal U}f$, so that
$\norm g_{L^2(\mathcal U)}\le1$, \cref{dir:lem:cubnorms} and Cauchy--Schwarz give
\[
 (1-d)\abs{\dirE_{\mu_S}f-\dirE_{\mathcal U}f}=\abs{\dirE_\mu g-\dirE_\mu(\mathbf 1_Rg)}\le\epsilon\sqrt V+\sqrt{2d},
\]
and, since $\dirE f^2\le(\dirE\abs f)^{2/3}(\dirE f^4)^{1/3}$,
\[
 (1-d)\dirE_{\mu_S}\abs f\ge\frac{(1-\epsilon V)^{3/2}}{\sqrt H}-\sqrt{2d}.
\]
Under the hypotheses $\epsilon V\le1/8$, $\epsilon^2VH=\epsilon^2V9^k+\epsilon^3V^3<1/256$ and
$\sqrt{2dH}\le1/16$, so the first quantity is at most $1/(8\sqrt H)$ and the second at least $1/(2\sqrt H)$, as
$(7/8)^{3/2}-1/16>1/2$; their ratio is at most $1/4$. For rational $\mu$ the feasible set is a nonempty
bounded rational polytope and has a rational vertex. The last sentence follows from $\delta H\le1/512$.
\end{proof}

The proof controls signed averages and fourth moments, not the total variation distance from $\mathcal U$:
the law can live on far fewer than $2^r$ points, no prime in each full sign cell is needed, and positivity on
every surviving point is stronger than membership of the target in the convex hull.

\begin{lemma}[A rational bump estimate; \stG]\label{dir:lem:bump}
For $X\ge2$ and a primitive real character $\chi$ of conductor $f$, the zeta case $f=1$ included,
\[
 \sum_pb(p/X)\chi(p)=\frac1{\log X}\Bigl(\mathbf 1_{f=1}XI_X+O\bigl(\sqrt X\log(2f)\bigr)\Bigr),\qquad
 I_X=\int_1^2b(t)\frac{\log X}{\log(Xt)}\,dt\in[\tfrac12,1],
\]
with an absolute implied constant.
\end{lemma}

\begin{proof}
Let $h_X(t)=b(t)\log X/\log(Xt)$ on $(1,2)$, zero elsewhere. Its derivatives through order three are bounded
on the interior uniformly in $X\ge2$, and $h_X$ and $h_X'$ vanish at both endpoints, so integration by parts in
$\log t$ gives $\abs{\widehat h_X(\sigma+it)}\ll(1+\abs t)^{-3}$ for $-1/2\le\sigma\le2$, uniformly in $X$, for
the Mellin transform $\widehat h_X(s)=\int_1^2h_X(t)t^{s-1}\,dt$. Shift the integral of
$-(L'/L)(s,\chi)X^s\widehat h_X(s)$ from $\Re s=2$ to $\Re s=-1/2$: the pole gives $XI_X$ for $f=1$; by
\cref{src:thm:dgrh} and the local zero count $O(\log(2f(\abs t+2)))$ the zeros contribute $O(\sqrt X\log(2f))$; a
trivial zero at $0$ contributes $O(1)$; and the functional equation bounds the new vertical integral by
$O(X^{-1/2}\log(2f))$. Higher prime powers contribute $\sum_{j\ge2}\vartheta((2X)^{1/j})\ll\sqrt X$, and on
primes $\log p\,h_X(p/X)=\log X\,b(p/X)$. The bounds on $I_X$ follow from $\log X\le\log(Xt)\le2\log X$.
\end{proof}

\begin{lemma}[Truncated bias of the prime pool; \stG]\label{dir:lem:primebias}
For each $k\ge1$ there is an effective $B_k$ such that, if $X\ge B_kr^2\log^2(2Y)$, then $\mathcal P_X$ is
nonempty, the law $\mu$ has truncated parity bias $\epsilon\le B_k\log(2Y)/\sqrt X$, and
$Z_X=\sum_{p\in\mathcal P_X}b(p/X)\gg X/\log X$ and $\abs{\mathcal P_X}\ll X/\log X$, the last two with absolute
constants after enlarging the lower bound on $X$ by an absolute factor.
\end{lemma}

\begin{proof}
For odd $p\nmid q_i$, $p\mapsto\legendre{q_i}p$ is induced by the primitive Kronecker character of
discriminant $q_i$ or $4q_i$ as $q_i\equiv1$ or $3\pmod4$. For $A\ne\varnothing$ the product $u_A$ is induced by
the nonprincipal primitive character of conductor $f_A=\prod_{i\in A}q_i$ or $4f_A$; for $\legendre p{q_i}$ the
Chinese remainder theorem gives conductor $\prod_{i\in A}q_i$. For $\abs A\le4k$, $\log(2f_A)\ll k\log(2Y)$.
Excluding the $q_i$ costs $O(r\log(2Y))$ before division by $\log X$, since an excluded prime in $(X,2X)$
satisfies $X<q_i\le Y$. By \cref{dir:lem:bump} the principal sum is $XI_X+O(\sqrt X+r\log(2Y))$, at least
$X/4$ at this height, and each nonempty parity sum is $O_k(\sqrt X\log(2Y)+r\log(2Y))$; divide. Chebyshev's
bound gives the count.
\end{proof}

\begin{proof}[Proof of \cref{dir:thm:cubature}]
As $V\ge r+1$, a large $C_k$ gives the height of \cref{dir:lem:primebias} and
$\epsilon\le1/(32\cdot3^kV+1024\delta V^2)$, which implies the conditions of the last sentence of
\cref{dir:thm:cubrepair}, also for $\delta=0$. The baseline is rational for integral $X$ and positive on every
prime of the open interval. The cell weights follow by expanding $2^{-\abs B}\prod_{i\in B}(1+v_iu_i)$.
\end{proof}

\begin{proof}[Proof of \cref{dir:cor:cubature}]
(a) The exact designs form a rational polytope $\set{w\ge0:\sum_pw_pv_p=e_0}$; at a vertex the columns in
the support are linearly independent, as a dependence would allow a perturbation in both directions, so there
are at most $V(r,k)$ of them. The Walsh monomials of degree at most $\floor{k/2}$ are orthonormal for any exact
design, as their pairwise products have degree at most $k$, so their number is at most the number of distinct
sign vectors in the support. (b) By \cref{dir:lem:primebias}, $\mu(p)\ll\log X/X$; renormalization after a
deletion multiplies it by at most $(1-c_k)^{-1}$, and \eqref{dir:eq:cubature} by at most $5/4$. (c) A set of
$\tau\abs{\mathcal P_X}$ primes has mass at most $\tau\abs{\mathcal P_X}\norm b_\infty/Z_X\le K\tau$.
\end{proof}

The deletion scale is sharp up to constants in degree one, for laws on the cube controlled only by their biases
through degree four.

\begin{theorem}[Three-layer obstruction; \stU]\label{dir:thm:threelayer}
Let $r\ge7$ be odd and $\ell\le r$ an odd integer with $\ell^2\ge6r$. There is a rational exchangeable law $\mu$
on $\set{-1,1}^r$ whose one- and two-coordinate parity moments vanish and whose largest parity bias through
degree four is $\asymp\ell^2/r^3$, such that deleting a set of mass $d\asymp r/\ell^2$ leaves a support carrying
no probability law with vanishing first moments. In particular $d\,\epsilon r^2\asymp1$, with absolute
constants.
\end{theorem}

\begin{proof}
Write $T=\sum_iu_i$ and let $\mu$ be uniform on each of the layers $T=1,\ell,-\ell$, with masses
$1-\lambda,\alpha,\beta$, where
\[
 \lambda=\frac{r-1}{\ell^2-1},\qquad\alpha=\frac12\Bigl(\lambda-\frac{1-\lambda}\ell\Bigr),\qquad
 \beta=\frac12\Bigl(\lambda+\frac{1-\lambda}\ell\Bigr);
\]
$\alpha\ge0$ is equivalent to $\ell\le r$. Then $\dirE T=0$, $\dirE T^2=r$, $\dirE T^3=r-\ell^2$ and
$\dirE T^4=r+(r-1)\ell^2$. The common parity moment $m_j$ of $j$ distinct coordinates is the coefficient of $z^j$
in $(1+z)^{(r+T)/2}(1-z)^{(r-T)/2}$ divided by $\binom rj$, so
\begin{gather*}
 m_1=\frac{\dirE T}r,\qquad m_2=\frac{\dirE T^2-r}{r(r-1)},\qquad m_3=\frac{\dirE T^3-(3r-2)\dirE T}{r(r-1)(r-2)},\\
 m_4=\frac{\dirE T^4-(6r-8)\dirE T^2+3r(r-2)}{r(r-1)(r-2)(r-3)},
\end{gather*}
which gives $m_1=m_2=0$, $m_3=(r-\ell^2)/(r(r-1)(r-2))$ and $m_4=(\ell^2-3r)/(r(r-2)(r-3))$, both of order
$\ell^2/r^3$ since $6r\le\ell^2\le r^2$. Delete the layer $T=-\ell$, of mass $\beta\in[\lambda/2,\lambda]$. On the
remaining layers $T\ge1$, so every probability law there has $\dirE T>0$.
\end{proof}

As $\ell$ grows from order $\sqrt r$ to $r$ the obstruction passes from a constant deletion fraction at bias of
order $r^{-2}$ to deletion of order $r^{-1}$ at bias of order $r^{-1}$, matching the scale
$\delta\epsilon V^2\asymp1$ of \cref{dir:thm:cubrepair} for $k=1$, $V=r+1$. It does not show that the
prime-height exponent of \cref{dir:thm:cubature} is sharp.

\begin{example}\label{dir:ex:cubature}
For $r=8$, $k=2$, $(q_i)=(3,5,7,11,13,17,19,23)$ and $X=100000$, delete primes whose first two signs are $+1$,
largest bump weights first, until the deleted mass reaches $0.02$, which exceeds $c_2$. Of the $8392$ primes of
the pool $94$ are deleted, and the $37$ moment equations of \eqref{dir:eq:cubature} hold exactly on the $8298$
survivors with weights between $0.8758$ and $1.1929$ times the renormalized baseline, computed from the
integer weights $(p-X)^2(2X-p)^2$ in rational arithmetic; the layer-moment formulas of the proof of
\cref{dir:thm:threelayer} were checked in $906$ cases.
\dirfiles{\artifact{prime_cubature_grh_2026-11-27/certificate.json},
\artifact{prime_cubature_grh_2026-11-27/verification.txt}}
\end{example}

Small-bias spaces and almost $k$-wise independence go back to \cite{NaorNaor1993SmallBias}, Fourier complexity in
effective Chebotarev bounds to \cite{Bellaiche2016Chebotarev}, and exact positive cubature and stable randomized
cubature are studied in \cite{Glaubitz2023Positive,MiglioratiNobile2022Stable}.

\subsection{Prime samples as codes for sparse character sums}\label{dir:sec:codes}
\sources{prime_character_codes_2026-11-27; prime_sparse_recovery_2026-11-27; prime_sparse_recovery_2026-11-27_2; grh_prime_fingerprints_2026-11-27; prime_legendre_codes_2026-11-27}

\Cref{dir:thm:signs} finds one prime realizing a prescribed sign word. Conversely, the values of
characters at one fixed short list of small primes can identify every sparse combination of
characters, even after an adversary corrupts a fraction of the values; \cref{rec:cor:dict} does this
for single characters. The optimal corruption fraction is set by subgroups of characters, and the
sampling input is the almost-orthogonality of characters on small primes given by
\cref{src:thm:dgrh}. The subsection ends with binary linear codes whose generator matrices are
Legendre symbols between two lists of primes, with both distances near the Gilbert--Varshamov
thresholds at a prime height $O(n^2\log^2n)$ (\cref{dir:thm:legendrecodes}).

For $q\ge2$ let $G_q=(\Z/q\Z)^\times$, with characters extended by zero off the units. For a finite set
$S$ of distinct primes put $S_q=\set{p\in S:p\nmid q}$ and $n_q=\abs{S_q}$. Let $\mathcal X_Q$ be the set of
primitive Dirichlet characters of conductor at most $Q$, the character of conductor $1$ included, each
evaluated at all integers with value zero at primes dividing its conductor; $\abs{\mathcal X_Q}\le Q^2$. For a
coefficient vector $c$ indexed by characters put $f_c=\sum_\chi c_\chi\chi$ and
$\sigma_k(c)_1=\min_{\abs T\le k}\sum_{\chi\notin T}\abs{c_\chi}$; a vector is $k$-sparse if it has at most $k$
nonzero entries.

\paragraph{Samples.}

\begin{lemma}[Universal small-bias prime samples; \stG]\label{dir:lem:codesample}
\begin{enumerate}[label=\textup{(\alph*)}]
\item There are absolute $B,K>0$ such that for every $Q\ge2$, with $L=\log(2Q)$, and every
$(1000Q^3)^{-1}\le b\le1/4$ there is a set $S$ of $n$ distinct primes in $[X,2X]$, $X=Kb^{-2}L^2$, with
$m/2\le n\le2m$ for $m=Bb^{-2}L$, and $\abs{\sum_{p\in S}\xi(p)}\le bn/8$ for every nonprincipal $\xi$
modulo every $2\le h\le Q^4$. If $M$ is the product or the least common multiple of at most $s\le Q$
integers at most $Q$ and $sb\le1/4$, then $\abs{S\setminus S_M}\le bn/8$, and every nonprincipal character
modulo $h\mid M$ with $h\le Q^4$, pulled back to $G_M$, has mean of modulus at most $b$ on $S_M$.
\item For $Q\ge3$ and $0<\mu\le1/2$ there is a set $P$ of $m\asymp\mu^{-2}\log(2Q)$ distinct primes, all
$\ll\mu^{-2}\log^2(2Q/\mu)$, with
$\abs{m^{-1}\sum_{p\in P}\overline{\chi(p)}\psi(p)-\mathbf 1_{\chi=\psi}}\le\mu$ for all $\chi,\psi\in\mathcal X_Q$,
the diagonal included; the constants are absolute.
\item For $q\ge2$ and $0<\delta<1/2$ there is a set $S$ of $\ll\delta^{-2}\log(2q)$ distinct primes not dividing
$q$, all $\ll\delta^{-2}\log(2q)\log(2q/\delta)$, on which every nonprincipal character modulo $q$ has mean
of modulus at most $\delta$; for rational $\delta$ it is found by a deterministic finite procedure. Also, for
every integer $X\ge C\delta^{-2}\log^2(2q)$ the primes $p\nmid q$ in $(X,2X)$ with weights proportional to
$(p-X)^3(2X-p)^3$ have bias at most $\delta$ and largest weight $\ll\log X/X$.
\end{enumerate}
\end{lemma}

\begin{proof}
(a) Let $W$ be smooth, nonnegative, supported in $(1,2)$ with $\int W=1$, and $w_p=(\log p)W(p/X)$.
\Cref{rec:lem:gram}, applied to the primitive character inducing $\xi$ (the primes dividing $h$ cost
$O_W(\log h)$), gives $\sum_pw_p\xi(p)=\mathbf 1_{\xi\text{ principal}}X+O_W(\sqrt X\log(2h))$ for every character $\xi$
modulo $h$; so
$W_0=\sum_pw_p\ge X/2$ and the weighted mean of every admissible $\xi$ is at most $b/32$ once $K$ is large,
since $\log(2h)\le4L$. Select each prime independently with probability $mw_p/W_0$, which is at most one
for large $K$ because $b^{-2}\le10^6Q^6$. Bernstein's inequality for the real and imaginary parts and a
union bound over at most $Q^8$ characters, with $B$ large, give a sample with $m/2\le n\le3m/2$ and all
sums at most $bm/16\le bn/8$. At most $s\log Q/\log X$ sampled primes divide $M$, a fraction at most
$2sb^2/(B\log X)\le b/8$; removing them changes each sum by at most their number, so the mean on $S_M$ is
at most $(b/8+b/8)/(1-b/8)<b$.

(b) Take $x=K\mu^{-2}\log^2(2Q/\mu)$ and select each prime of $[x,2x]$ with probability $M_0w_p/\sum w$,
$M_0=B\mu^{-2}\log(2Q)$. By \cref{rec:lem:gram} for $\chi,\psi\in\mathcal X_Q$, of conductors at most $Q$, the weighted
Gram deviation is at most $\mu/8$ once $K$ is large, each centred summand of
$\sum_p\mathbf 1_{p\in P}(\overline{\chi(p)}\psi(p)-\mathbf 1_{\chi=\psi})$ has modulus at most one, and Bernstein's inequality
over at most $Q^4$ pairs gives the sample.

(c) For the weight $w(t)=(t-1)^3(2-t)^3$ on $[1,2]$, which is $C^2$, argue as in the proof of \cref{dir:lem:bump}
with $g_X(t)=w(t)\log X/\log(Xt)$, whose Mellin transform satisfies $\abs{H_X(1/2+i\gamma)}\ll(1+\abs\gamma)^{-2}$ by
two integrations by parts; this gives
\[
 \sum_{p\nmid q}w(p/X)\chi(p)=\mathbf 1_{\chi=1}X\int_1^2\frac{w(t)\,dt}{\log(Xt)}+O\Bigl(\frac{\sqrt X\log(2q)}{\log X}\Bigr),
\]
hence the weighted statement. For the thinned set take $X=\ceil{B\delta^{-2}\log(2q)\log(2q/\delta)}$ and
$m=\ceil{256\delta^{-2}\log(16\varphi(q))}$, and retain each prime with probability $mw(p/X)/\sum_pw(p/X)$.
For each of the $2\varphi(q)$ real and imaginary parts $a_p$ of the characters and each sign, the exponential
moment of $\lambda(\pm(\sum_pa_p(B_p-\pi_p))-\delta m/8)$ with $\lambda=\delta/16$ is at most
$e^{-\delta^2m/256}$, so the sum $\Phi$ of these moments has mean at most $1/4$; revealing the retention
decisions one by one and never increasing the conditional mean of $\Phi$ (computed to accuracy $1/(8N)$
for $N$ candidate primes) ends with $\Phi<1$. Then the number of retained primes lies in $(m/2,2m)$, and
every nonprincipal mean is at most $\delta$.
\end{proof}

\paragraph{Rigidity of near-minimum support.}
Let $\mu$ be a probability law on a finite abelian group $G$, $\widehat\mu(\gamma)=\dirE_\mu\gamma$, and
$F=\sum_{\chi\in A}c_\chi\chi$ with all $c_\chi\ne0$ and $t=\abs A$.

\begin{lemma}[Gram and Bessel bounds; \stU]\label{dir:lem:codegram}
If $\gamma_1,\dots,\gamma_r$ are distinct characters with $\abs{\widehat\mu(\gamma_i\gamma_j^{-1})}\le b$ for $i\ne j$,
then $(1-(r-1)b)\norm a_2^2\le\dirE_\mu\abs{\sum_ia_i\gamma_i}^2\le(1+(r-1)b)\norm a_2^2$ for $a\in\C^r$, and
$\alpha\sum_i\abs{\dirE_{\mu_U}\gamma_i}^2\le1+(r-1)b$ for every $U$ of measure $\alpha>0$, with $\mu_U$ the
conditional law.
\end{lemma}

\begin{proof}
The off-diagonal terms contribute at most $b((\sum_i\abs{a_i})^2-\norm a_2^2)\le(r-1)b\norm a_2^2$. For the second
bound apply the adjoint of $a\mapsto\sum_ia_i\gamma_i$, of squared norm at most $1+(r-1)b$, to $\mathbf 1_U$.
\end{proof}

Fix $s\ge2$ and put
\[
 \eta_s=\frac1{1000s^3},\qquad\eps_s=\frac1{100s^2},\qquad\tau_s=\frac1{100s},\qquad
 d_s(\tau)=\frac{1-(s-1)\eta_s-\tau^2}s .
\]

\begin{theorem}[Support rigidity under local small bias; \stU]\label{dir:thm:rigid}
\begin{enumerate}[label=\textup{(\alph*)}]
\item If $\abs{\widehat\mu(\gamma)}\le b$ for $1\ne\gamma\in AA^{-1}$, then
$\mu\set{\abs F>\tau\norm c_2}\ge(1-(t-1)b-\tau^2)/t$ for every $\tau\ge0$.
\item Let $t\le s$ and suppose $\abs{\widehat\mu(\gamma)}\le\eta_s$ for every $1\ne\gamma\in DD^{-1}$, $D=AA^{-1}$.
For $0\le\tau\le\tau_s$ let $U=\set{\abs F>\tau\norm c_2}$ and $\alpha=\mu(U)$. Then $\alpha\ge d_s(\tau)$. If
$\alpha\le(1+\eps_s)/s$, then $t=s$ and $A=\chi_0H$ for a subgroup $H\le\widehat G$ of order $s$, and with
$H^\perp=\set{u:h(u)=1\ (h\in H)}$ the set $U\subseteq G$ is a single coset $C$ of $H^\perp$, and for some
$\lambda\ne0$ one has $F=\lambda\chi_0$ on $C$ and
\[
 \norm{F-\lambda\chi_0\mathbf 1_C}_{\infty}\le\tau\norm c_2,\qquad\norm{c-c_*}_2\le\tau\sqrt{1-1/s}\,\norm c_2,
\]
$c_*$ being the coefficient vector of $\lambda\chi_0\mathbf 1_C$. Conversely, for $\tau=0$ every nonzero
modulated indicator $\lambda\chi_0\mathbf 1_C$ of a coset of $H^\perp$ has Fourier support $\chi_0H$ and
satisfies $\alpha\le(1+\eps_s)/s$.
\end{enumerate}
\end{theorem}

\begin{proof}
Normalize $\norm c_2=1$. (a) By \cref{dir:lem:codegram}, $\dirE_\mu\abs F^2\ge1-(t-1)b$, while $\abs F^2\le t$
everywhere and $\abs F^2\le\tau^2$ off $U$.

(b) Part (a) with $b=\eta_s$ gives $\alpha\ge d_s(\tau)$. Put $a_0=(s-1)\eta_s+\tau^2\le1/(500s^2)$. If $t\le s-1$,
then $(1-a_0)/(s-1)\le(1+\eps_s)/s$, which is false; so $t=s$. Let $u_i=c_i\chi_i$ and
$\rho=s-\dirE_{\mu_U}\abs F^2$. Since $\dirE_\mu\abs F^2\le\alpha\dirE_{\mu_U}\abs F^2+\tau^2$, the energy bound and
the upper threshold give $\rho\le s(\eps_s+a_0)\le1/(80s)$. The identity
$s\sum_i\abs{u_i-F/s}^2=s-\abs F^2$ gives $\dirE_{\mu_U}\abs{u_i-u_j}^2\le2\rho/s$, hence $\abs{c_i}\ge1/(2\sqrt s)$
for all $i$, and then every ratio in $D$ has conditional bias at least $1-4\rho\ge1-1/(20s)$. If
$\abs D\ge s+1$, the Bessel bound for $s+1$ of these ratios gives
$1+s\eta_s\ge(s+1)\alpha(1-1/(20s))^2\ge(1+1/s)(1-1/(500s^2))(1-1/(20s))^2\ge1+5/(6s)$, a contradiction; so
$\abs D\le s$, $D=A\chi_0^{-1}$ for any $\chi_0\in A$, and $D$ is closed under ratios, a subgroup $H$ of order
$s$. For a coset $C=gH^\perp$, $\mathbf 1_C=s^{-1}\sum_{h\in H}\overline{h(g)}h$, and $H\subseteq DD^{-1}$ gives
$(1-(s-1)\eta_s)/s\le\mu(C)\le(1+(s-1)\eta_s)/s$. Since $F/\chi_0$ is constant on cosets of $H^\perp$, $U$ is a
union of such cosets, nonempty, and two of them would have measure above $(1+\eps_s)/s$. So the large set
is one coset $C$; with $\lambda=(F/\chi_0)|_C$ the difference $F-\lambda\chi_0\mathbf 1_C$ vanishes on $C$ and is at
most $\tau$ elsewhere on all of $G$, and Parseval for the uniform measure gives the coefficient bound. The
converse follows from the coset identity and $(s-1)\eta_s\le\eps_s$.
\end{proof}

For the uniform measure and $\tau=0$, part (a) is the Donoho--Stark support inequality on a finite
abelian group, whose equality cases are known \cite{DonohoStark1989Uncertainty,MatusiakOzaydinPrzebinda2004DonohoStark,Meshulam2006Uncertainty}; the content of (b) is the exact algebraic
structure forced by a nonuniform small-bias sample and an approximate-zero threshold. In the same way, if
$F_1,\dots,F_r$ come from an orthonormal basis of an $r$-dimensional space of coefficient vectors on $A$, then
$\mu\set{\sum_i\abs{F_i}^2>u^2}\ge(r(1-(t-1)b)-u^2)/t$, since $\sum_i\abs{F_i}^2\le t$ pointwise by Bessel.

\begin{theorem}[Prime codes: rigidity and exact distance; \stG]\label{dir:thm:codes}
There are absolute $B,K>0$ such that for every $Q\ge2$ and $2\le s\le Q$ there is a set $S$ of distinct
primes in $[X,2X]$, $X=K\eta_s^{-2}\log^2(2Q)$, with
$B\eta_s^{-2}\log(2Q)/2\le\abs S\le2B\eta_s^{-2}\log(2Q)$, for which \cref{dir:thm:rigid}(b) holds for the
empirical law of the residues of $S_q$, for every $2\le q\le Q$ and every nonzero $s$-sparse character sum
on $G_q$. For $s=2k$ the code
$\mathcal C_{q,k}=\set{(F(p))_{p\in S_q}:F\text{ $k$-sparse}}$ has positive minimum Hamming distance $D_{q,k}$.
If $s\mid\varphi(q)$, then
\[
 D_{q,k}=\min_{\substack{H\le\widehat G_q\\\abs H=s}}\ \min_{g\in G_q}\ \#\set{p\in S_q:p\bmod q\in gH^\perp}\le\Bigl\lfloor\frac{n_q}s\Bigr\rfloor,
\]
every difference attaining it is a modulated coset indicator with a minimizing coset, and every such
indicator is a difference of two $k$-sparse sums; if $s\nmid\varphi(q)$, then $D_{q,k}>(1+\eps_s)n_q/s$. With
$u$ known erasures and $e$ errors, exact recovery is guaranteed in the worst case precisely when
$2e+u<D_{q,k}$. If a received vector agrees to accuracy $v\ge0$ with each of two $k$-sparse evaluation
vectors outside at most $e$ positions and $2e<d_s(\tau_s)n_q$, their coefficient vectors satisfy
$\norm{c-c'}_2\le200s\,v$; if only $2e\le(1+\eps_s)n_q/s$ and $\norm{c-c'}_2>200s\,v$, the difference has the
structure of \cref{dir:thm:rigid}(b), and the two supports are disjoint with exactly $k$ elements each.
\end{theorem}

\begin{proof}
Take $b=\eta_s$ in \cref{dir:lem:codesample}(a) with $M=h=q$: the residues of $S_q$ have bias at most
$\eta_s$ on every nonprincipal character, so \cref{dir:thm:rigid}(b) applies to every support. A difference of
two $k$-sparse sums is $s$-sparse and conversely, so $\tau=0$ gives injectivity and the distance bounds; if
$s\mid\varphi(q)$, $\widehat G_q$ has a subgroup of order $s$, the modulated coset indicators have weight
below $(1+\eps_s)n_q/s$, and the $s$ cosets of one subgroup partition $S_q$. Two Hamming balls of radius $e$
are disjoint when $2e<D_{q,k}$, and a word copying each member of a minimum-distance pair on half of
their differing positions lies in both otherwise; erasures are handled in the same way. For noisy data the
difference $F$ has $\abs{F(p)}\le2v$ outside at most $2e$ positions; if $\norm{c-c'}_2>2v/\tau_s=200sv$, the set
where $\abs F>\tau_s\norm{c-c'}_2$ has at most $2e$ elements, and \cref{dir:thm:rigid}(b) applies.
\end{proof}

\begin{theorem}[Unknown and mixed conductors; \stG]\label{dir:thm:codesmixed}
For the set $S$ of \cref{dir:thm:codes}, every nonzero $F=\sum_{\chi\in A}c_\chi\chi$ with
$A\subseteq\mathcal X_Q$ and $1\le\abs A\le s$ satisfies
$\#\set{p\in S:\abs{F(p)}>\tau\norm c_2}\ge(1-\eta_s/8)d_s(\tau)\abs S$ for $0\le\tau\le\tau_s$. If the left side is
at most $(1-\eta_s/8)(1+\eps_s)\abs S/s$, then the characters of $A$, inflated to $G_M$ with $M$ the least
common multiple of their conductors, form a coset of a subgroup of order $s$, and $F$ has on $G_M$ the
structure of \cref{dir:thm:rigid}(b). For $s=2k$ all $k$-sparse mixtures are identifiable from their values
on $S$, without knowledge of the conductors, and the two-candidate bound $\norm{c-c'}_2\le200s\,v$ holds when
$2e+u<(1-\eta_s/8)d_s(\tau_s)\abs S$, with $u$ known erasures.
\end{theorem}

\begin{proof}
The inflated characters are distinct. Each nontrivial $\gamma\in DD^{-1}$ is a product of at most four of them or
their inverses, so it factors through $G_h$ with $h\mid M$ and $h\le Q^4$, nonprincipally; by
\cref{dir:lem:codesample}(a) the residues of $S_M$ satisfy the local bias condition with $\eta_s$, and
$\abs{S_M}\ge(1-\eta_s/8)\abs S$. Apply \cref{dir:thm:rigid} on $S_M$; for recovery apply the lower bound to
differences.
\end{proof}

With $b=\delta/s$ in \cref{dir:lem:codesample}(a), $0<\delta<1/4$, part (a) of \cref{dir:thm:rigid} with
$\tau=\sqrt\delta$ gives a cheaper sample, of $O_\delta(s^2\log Q)$ primes of size $O_\delta(s^2\log^2Q)$, on which
every nonzero $s$-sparse sum, modulo any $q\le Q$ or from $\mathcal X_Q$, exceeds $\sqrt\delta\norm c_2$ on a
fraction at least $(1-\delta/(8s))(1-2\delta)/s$ of the sample, the factor $1-\delta/(8s)$ being needed only for
mixed conductors; this gives recovery near the optimal radius, with coefficient error $2v/\sqrt\delta$, but
no classification of extremizers.

\paragraph{Convex decoding.}
The distance alone gives a nearest-codeword decoder; a weighted $\ell^1$ program reaches the same
corruption fraction and is stable. Let $A\in\C^{m\times D}$ and data $y_j=(Ac)_j+b_j+w_j$ for $j\notin E$, where $E$
is a known set of $e$ erased rows, $b$ has at most $t$ nonzero entries outside $E$ (unknown positions and
values), and $w$ is arbitrary noise.

\begin{theorem}[Robust convex decoding; \stU]\label{dir:thm:decoder}
\begin{enumerate}[label=\textup{(\alph*)}]
\item Suppose that every entry of $A$ has modulus at most one and that every entry of $A^*A/m-I$,
the diagonal included, has modulus at most $\mu$. If $d_0=1-2k\mu-2k(e+2t)/m>0$, a minimizer of
\[
 J(z)=\norm z_1+\frac{2k}m\sum_{j\notin E}\abs{y_j-(Az)_j}
\]
exists, and every minimizer satisfies
\[
 \norm{\widehat c-c}_1\le\frac1{d_0}\Bigl(2\sigma_k(c)_1+\frac{4k}m\sum_{j\notin E}\abs{w_j}\Bigr);
\]
for $k$-sparse $c$ and $w=0$ it is unique and equals $c$.
\item Let $A_{p,\chi}=\chi(p)$ for the characters of a finite abelian group and a weighted sample with
probability weights $\mu_p$, so that $A^*D_\mu A$ has unit diagonal, and suppose its off-diagonal entries are at
most $\delta$. With observations $y=Ac+e'+\nu$ where the corruption $e'$ is supported on a set of weight at
most $\alpha$ and $\sum_p\mu_p\abs{\nu_p}\le\eta'$, if $\Gamma=1-(2k-1)\delta-4k\alpha>0$, every minimizer of
$\norm z_1+\frac{2k}{1+\delta}\sum_p\mu_p\abs{(Az-y)_p}$ satisfies
$\norm{\widehat c-c}_1\le(2(1+\delta)\sigma_k(c)_1+4k\eta')/\Gamma$.
\end{enumerate}
\end{theorem}

\begin{proof}
(a) Coercivity gives a minimizer. Let $U=E^c$, $T\subseteq U$ the support of $b$ and $S_0$ a set of $k$
largest entries of $c$. For $h\in\C^D$, entrywise control of the Gram error and $\abs{A_{ji}}\le1$ give
$\norm h_\infty\le\norm{(Ah)_U}_1/m+(\mu+e/m)\norm h_1$, so
$\norm{h_{S_0}}_1\le(k/m)\norm{(Ah)_U}_1+k(\mu+e/m)\norm h_1$. Also
$\norm{c+h}_1-\norm c_1\ge\norm h_1-2\norm{h_{S_0}}_1-2\sigma_k(c)_1$, and the residual term increases by at least
$\norm{(Ah)_U}_1-2t\norm h_1-2\norm{w_U}_1$. At the weight $2k/m$ the terms in $\norm{(Ah)_U}_1$ cancel, giving
$J(c+h)-J(c)\ge d_0\norm h_1-2\sigma_k(c)_1-(4k/m)\norm{w_U}_1$; take $h=\widehat c-c$.
(b) For every $\chi$, $\abs{h_\chi}\le\norm{Ah}_{1,\mu}+\delta(\norm h_1-\abs{h_\chi})$, so summing over a set of $k$
largest entries gives $(1+\delta)\norm{h_{S_0}}_1\le k\norm{Ah}_{1,\mu}+k\delta\norm h_1$; with the same two
triangle inequalities the terms in $\norm{Ah}_{1,\mu}$ cancel at the weight $2k/(1+\delta)$.
\end{proof}

\begin{theorem}[Prime decoders; \stG]\label{dir:thm:primedecoder}
Fix $0<\eps<1$.
\begin{enumerate}[label=\textup{(\alph*)}]
\item For $Q\ge3$ and $1\le k\le\abs{\mathcal X_Q}$ there is a set $P$ of $m\asymp_\eps k^2\log(2Q)$ distinct primes,
all $\ll_\eps k^2\log^2(2Q)$, such that for every $c\in\C^{\mathcal X_Q}$ and every pattern with
$\Delta=1-\eps-2k(e+2t)/m>0$ every minimizer of $J$ for $A=(\chi(p))_{p\in P,\chi\in\mathcal X_Q}$ satisfies
\[
 \norm{\widehat c-c}_1\le\frac1\Delta\Bigl(2\sigma_k(c)_1+\frac{4k}m\sum_{j\notin E}\abs{w_j}\Bigr).
\]
In particular $k$-sparse mixtures of primitive characters are recovered exactly whenever
$e+2t<(1-\eps)m/(2k)$.
\item Let $q\ge2$, $1\le k\le\varphi(q)$ and $\delta=\eps/(4k)$. The weighted window of
\cref{dir:lem:codesample}(c) at some height $X\ll k^2\eps^{-2}\log^2(2q)$, and also a suitable set of
$\ll k^2\eps^{-2}\log(2q)$ distinct primes, all of size $\ll k^2\eps^{-2}\log(2q)\log(2qk/\eps)$, have the
following property. For corruption weight $\alpha\le(1-\eps)/(4k)$, every minimizer of the program in
\cref{dir:thm:decoder}(b) satisfies
\[
 \norm{\widehat c-c}_1\le\frac{4(1+\delta)}\eps\sigma_k(c)_1+\frac{8k}\eps\eta'.
\]
\end{enumerate}
\end{theorem}

\begin{proof}
(a) Use \cref{dir:lem:codesample}(b) with $\mu=\eps/(2k)$, so $\log(2Q/\mu)\ll_\eps\log(2Q)$ as
$k\le Q^2$, and \cref{dir:thm:decoder}(a), where $d_0\ge\Delta$. (b) Use \cref{dir:lem:codesample}(c) with
bias $\delta$; the off-diagonal Gram entries are means of nonprincipal characters, and $\Gamma\ge\eps/2$.
\end{proof}

\begin{theorem}[Obstructions; \stU]\label{dir:thm:codeobstruction}
\begin{enumerate}[label=\textup{(\alph*)}]
\item If $\widehat G$ has a subgroup of order $2k$, then for every fixed design of units of $G$ with
probability weights $\mu_p$ there are two distinct $k$-sparse signals and one observation consistent with
each after corruption of weight at most $1/(4k)+\frac12\max_p\mu_p$; for $M$ unweighted units, after at most
$\ceil{\floor{M/(2k)}/2}$ corruptions.
\item If some prime $q\le Q$ has $q\equiv1\pmod{4k}$, then for every list of $m$ primes, of any size and with
repetitions allowed, two distinct exactly $k$-sparse vectors in $\C^{\mathcal X_Q}$ have measurements
differing in at most $\floor{m/(2k)}$ positions; so no decoder guarantees exact recovery for all $k$-sparse
vectors once $e+2t\ge\floor{m/(2k)}$. Under \cref{src:thm:dgrh} \textup{(\stG)} such a prime exists for
$Q\ge Ck^2\log^2(2k)$, with $C$ absolute.
\item If $\widehat G$ has a subgroup $H$ of order $2k$, then for every $0\le\delta<1/(2k-1)$ some weighted design
of bias exactly $\delta$ attains the bound of \cref{dir:thm:rigid}(a) with $\tau=0$ and $t=2k$ by the difference
of two $k$-sparse signals, which are confused at corruption weight $(1-(2k-1)\delta)/(4k)$ each.
\end{enumerate}
\end{theorem}

\begin{proof}
(a) Let $K=H^\perp$ and choose a coset $aK$ of sample weight at most $1/(2k)$. Split $H=H_1\sqcup H_2$ into
halves and put $c^{(1)}_\chi=\overline{\chi(a)}$ on $H_1$, $c^{(2)}_\chi=-\overline{\chi(a)}$ on $H_2$; the difference of
the evaluations is $2k\mathbf 1_{aK}$. Assign the differing coordinates one by one to the side of smaller
accumulated weight, so the two sides differ by at most $\max_p\mu_p$, and let the observation follow
$c^{(1)}$ on one side and $c^{(2)}$ on the other. (b) Let $\theta$ generate the characters modulo $q$ and
$\omega=\theta^{(q-1)/(2k)}$; the $2k$ characters $\chi_j=\theta\omega^j$ are distinct, nonprincipal, hence
primitive, and vanish at $q$. For a $2k$th root of unity $v$, $\sum_jv^{-j}\chi_j(n)$ equals $2k\theta(n)$ when
$\omega(n)=v$ and $0$ otherwise. Choose the least populous of the $2k$ classes $\omega(p)=v$ among the
sample and split the sum into two halves of $k$ terms with opposite signs. Under \cref{src:thm:dgrh}
the smoothed prime sum over $p\equiv1\pmod{4k}$ at height $X=K_0k^2\log^2(2k)$ is
$X/\varphi(4k)+O(\sqrt X\log(8k))>0$. (c) Give the elements of $K=H^\perp$ weight $(1-(2k-1)\delta)/\abs G$ and
the others $(1+\delta)/\abs G$; the nonprincipal means are $-\delta$ on $H\setminus\set1$ and zero
otherwise, and the signals $\mathbf 1_{H_1}$, $-\mathbf 1_{H_2}$ differ by $2k\mathbf 1_K$; split each residue
into two coordinates of half its weight.
\end{proof}

Since a finite abelian group has subgroups of every order dividing its order, the limiting fraction
$1/(4k)$ of \cref{dir:thm:codes,dir:thm:primedecoder} cannot be improved for any modulus with
$2k\mid\varphi(q)$, whatever the sampling set, and for $t=2k$ the bound of \cref{dir:thm:rigid}(a) with
$\tau=0$ is sharp at every bias $\delta<1/(2k-1)$.

\begin{corollary}[Sparse packets of $L$-functions; \stG]\label{dir:cor:packets}
If at most $k$ of the $c_\chi$, $\chi\in\mathcal X_Q$, are nonzero, the coefficients of
$-\sum_\chi c_\chi(L'/L)(s,\chi)$ at the primes of \cref{dir:thm:primedecoder}(a), divided by $\log p$, determine
$c$ with the same adversarial guarantees; for nonnegative integers $c_\chi$ this recovers the factors and
multiplicities of $\prod_\chi L(s,\chi)^{c_\chi}$. If the $c_\chi$ are integers and the stability bound is below
$1/2$, rounding gives exact recovery in the presence of noise.
\end{corollary}

\begin{proof}
For $\Re s>1$, $-(L'/L)(s,\chi)=\sum_n\Lambda(n)\chi(n)n^{-s}$, so the coefficient at $p$ is $(\log p)f_c(p)$.
\end{proof}

The primitive convention matters: two imprimitive characters induced by one primitive character agree
at all primes not dividing their moduli, so such aliases cannot be separated by samples avoiding those
divisors.

\paragraph{List decoding of characters of prime order.}
For a fixed prime $\ell$ let $\mathcal X_{Q,\ell}\subseteq\mathcal X_Q$ consist of the characters of order dividing
$\ell$, and $M_\ell=\abs{\mathcal X_{Q,\ell}}$. A list $P=(p_1,\dots,p_n)$ of primes gives each $\chi$ a fingerprint:
at $p_i\nmid f_\chi$ the symbol $a\in\F_\ell$ with $\chi(p_i)=e^{2\pi ia/\ell}$, and a wildcard at $p_i\mid f_\chi$. For
fixed $0\le e<1$ and $0\le b<1-1/\ell$, a received word is a set $I$ of observed positions with symbols
$y\in\F_\ell^I$, and $\chi$ is compatible with it if $\abs I\ge(1-e)n$ and
$\#\set{i\in I:p_i\nmid f_\chi,\ \chi(p_i)\ne e^{2\pi iy_i/\ell}}\le b\abs I$; ramified positions are free. Put
\[
 \mathsf H_\ell(u)=-u\log_\ell u-(1-u)\log_\ell(1-u)+u\log_\ell(\ell-1),\qquad
 K_\ell(e,b)=(1-e)\bigl(1-\mathsf H_\ell(b)\bigr).
\]

\begin{theorem}[Universal prime fingerprints; \stG]\label{dir:thm:fingerprint}
Fix $\ell,e,b$ and $0<R<K_\ell(e,b)$, and an integer $r$ with $r(K_\ell(e,b)-R)>\log_\ell(\ell+1)$. There are
$B,c>0$ such that for all large $Q$ some list of $n=\ceil{\log_\ell M_\ell/R}$ distinct primes in
$(B\log^2Q,2B\log^2Q)$ has at most $\ell^{r-1}$ compatible characters in $\mathcal X_{Q,\ell}$ for every received
word; a weighted random choice without replacement fails with probability at most $\ell^{-cn}$.
\end{theorem}

\begin{theorem}[Converse; \stU]\label{dir:thm:fingerprintconverse}
If lists of lengths $n(Q)$, of primes of any size, have at most $L_0$ compatible characters for every
received word, with $L_0$ fixed, then
\[
 \limsup_{Q\to\infty}\frac{\log_\ell M_\ell}{n(Q)}\le K_\ell(e,b);
\]
the same holds for every alphabet-$\ell$ fingerprint scheme.
\end{theorem}

\begin{corollary}[Quadratic characters at the optimal length; \stG]\label{dir:cor:fingerprint}
For $e<1$, $b<1/2$ and $\eps>0$ there are $L_0$ and, for large $Q$, lists of at most
$(1+\eps)\log_2Q/((1-e)(1-\mathsf H_2(b)))+O(1)$ distinct primes of size $O(\log^2Q)$ such that every allowed
observation of Legendre symbols $\legendre Dp$ is compatible with at most $L_0$ fundamental discriminants
$\abs D\le Q$, positive or negative, without knowledge of the conductor; the leading coefficient is optimal
for bounded lists.
\end{corollary}

\begin{proof}[Proof of \cref{dir:thm:fingerprint,dir:thm:fingerprintconverse,dir:cor:fingerprint}]
\emph{Counting.} The number of characters of order dividing $\ell$ at modulus $f$ is at most
$\ell^{\omega(f)+1}=f^{o(1)}$, and primes $p\equiv1\pmod\ell$ up to $Q$ give $\gg_\ell Q/\log Q$ of them, so
$M_\ell=Q^{1+o(1)}$; for $\ell=2$, fundamental discriminants give $M_2\asymp Q$. Hence $n\asymp\log Q$.

\emph{Bounded-rank equidistribution.} Let $w$ be smooth in $(1,2)$ with $\int w=1$ and
$\mu_X(p)\propto w(p/X)\log p$ with $X=B\log^2Q$. Characters of order dividing $\ell$ form an $\F_\ell$-vector
space in the direct limit over moduli, and a primitive character is the unique representative of its
class. For independent $\chi_1,\dots,\chi_r\in\mathcal X_{Q,\ell}$, pad wildcards by zero; every nonzero Fourier
index $a\in\F_\ell^r$ corresponds to the nonprincipal character $\prod_j\chi_j^{a_j}$ of modulus at most $Q^r$, so
\cref{rec:lem:gram}, applied to the primitive character inducing it (padding and ramified primes cost
$o(1)$ after normalization), bounds each nontrivial coefficient of the law of the padded column by
$O(r\log(2Q)/\sqrt X)+o(1)=O(r/\sqrt B)+o(1)$.
Fourier inversion on $\F_\ell^r$ gives $\dirP(\text{column}=v)\le(1+\eta)\ell^{-r}$ for every $v$, for any fixed
$\eta>0$ once $B$ is large. Draw the $n$ primes successively from $\mu_X$ restricted to unused primes; the
removed mass is $O(\log\log Q/\log Q)$, so any prescribed $r\times m$ matrix of symbols on $m$ fixed positions
has probability at most $(1+\eta)^m\ell^{-rm}$, by conditioning in order.

\emph{Rank witnesses.} Choose $b'>b$ with $r(K_\ell-(1-e)(1-\mathsf H_\ell(b')))<\Delta/3$, where
$\Delta=r(K_\ell-R)-\log_\ell(\ell+1)$, and $\log_\ell(1+\eta)<\Delta/3$. A character has at most
$\log Q/\log X=o(n)$ ramified positions, so for large $Q$ each compatible candidate has a padded fingerprint
within Hamming distance $b'\abs I$ of $y$ on $I$. More than $\ell^{r-1}$ candidates span a space of dimension at
least $r$ and contain $r$ independent ones. Hamming balls satisfy $V_\ell(m,\floor{b'm})\le\ell^{m\mathsf H_\ell(b')}$;
summing over at most $M_\ell^r$ tuples and $(\ell+1)^n$ received words, with $m\ge(1-e)n$ and
$\log_\ell M_\ell\le Rn$, the failure probability is at most $\ell^{-\Delta n/3}$.

\emph{Converse.} Fix $I$ of size $m=\ceil{(1-e)n}$, pad wildcards by zero, and double count pairs
$(\chi,y)$ with $y$ in the Hamming ball of radius $\floor{bm}$ about the fingerprint of $\chi$: this gives
$\log_\ell M_\ell\le m-\log_\ell V_\ell(m,\floor{bm})+\log_\ell L_0$, and $\log_\ell V_\ell(m,\floor{bm})=m\mathsf H_\ell(b)+O(\log m)$.
The corollary takes $\ell=2$ and $R$ between $K_2/(1+\eps)$ and $K_2$.
\end{proof}

For $e=b=0.1$ one has $K_2\simeq0.4779$, so the optimal leading length is about $2.092\log_2Q$; at
$R=0.4$ the theorem allows $r=21$, that is lists of at most $2^{20}$ candidates. Since a list of
$\asymp\log Q$ distinct primes reaches size $\gg\log Q\log\log Q$, the height $\log^2Q$ is not known to be
necessary.

\begin{example}\label{dir:ex:codes}
\begin{enumerate}[label=\textup{(\arabic*)}]
\item For $Q=31$ and $k=2$, $1616$ distinct primes in $[40000,80000]$ have complex Gram deviation at most
$61/1000$ over all $197$ primitive characters, certified in rational arithmetic; with $61$ erasures and $91$
corruptions the margin of \cref{dir:thm:decoder}(a) is $7803/50500>0$, and a cyclic packet modulo $17$
gives a collision with $196$ errors per candidate.
\item For $q=24$, with unit group $C_2^3$ and $k=2$, the $464$ primes in $(4096,8192)$ with the weights of
\cref{dir:lem:codesample}(c) have bias $0.0360834780$, and a thinned set of $202$ of them has bias $7/101$;
with $25$, respectively $18$, corrupted samples the certified margins are positive.
\item For $Q=128$, the sixteen primes $103,113,\dots,293$, at most two erasures and two errors, an exhaustive
search over all $79$ fundamental discriminants and all received words gives largest list size five,
attained for instance by $-83,24,61,97,104$.
\end{enumerate}
\dirfiles{\artifact{prime_sparse_recovery_2026-11-27/certificate_output.txt},
\artifact{prime_sparse_recovery_2026-11-27/sample_primes.txt},
\artifact{prime_sparse_recovery_2026-11-27_2/example_output.txt},
\artifact{prime_sparse_recovery_2026-11-27_2/prime_pools.csv},
\artifact{grh_prime_fingerprints_2026-11-27/example.json}}
\end{example}

Sparse recovery from corrupted measurements and coherence bounds are treated in
\cite{StuderKuppingerPopeBolcskei2012Recovery,IosevichMayeli2023Uncertainty}, and list decoding of random
linear codes in \cite{GuruswamiHastadKopparty2011List}; here the messages are characters of bounded
conductor and each coordinate is an evaluation at a distinct small prime.

\paragraph{Binary codes from Legendre symbols between two prime lists.}
For disjoint lists $Q=(q_1,\dots,q_r)$ and $P=(p_1,\dots,p_n)$ of odd primes, without repetitions within either
list, put
\[
 A=A(Q,P)\in\F_2^{r\times n},\qquad A_{ij}=\frac{1-\legendre{p_j}{q_i}}2\bmod2,
\]
and let $C=\set{uA:u\in\F_2^r}$ and $C^\perp=\set{v\in\F_2^n:Av^t=0}$, with minimum nonzero Hamming weights
$d(C)$ and $d(C^\perp)$. Let $\mathsf H_2$ be the binary entropy, inverted on $[0,1/2]$, and for $I\subseteq[r]$,
$J\subseteq[n]$ put $\Delta_A(I,J)=\min(\abs I,\abs J)-\rank A_{I,J}$ and $\Delta(A)=\max_{I,J}\Delta_A(I,J)$, the
rank of an empty matrix being zero; the rank condition is uniform over all $2^{r+n}$ rectangles. For $S\subseteq[r]$
the sign word $\sigma_S(j)=\prod_{i\in S}\legendre{p_j}{q_i}$ is the codeword of the indicator of $S$ with each
bit $e$ replaced by $(-1)^e$.

\begin{theorem}[Legendre codes at a near-quadratic prime cutoff; \stG]\label{dir:thm:legendrecodes}
Fix $0<R<1$, $0<\delta,\delta^\perp<1/2$ and $\tau>0$ with $\mathsf H_2(\delta)<1-R$ and
$\mathsf H_2(\delta^\perp)<R$. There is $K$, depending only on these parameters, such that for every large $n$,
with $r=\floor{Rn}$ and $X=\ceil{Kn^2\log^2n}$, there are distinct primes $q_i\in(X,2X)$ and $p_j\in(3X,6X)$
with
\[
 \rank A=r,\qquad d(C)\ge\ceil{\delta n},\qquad d(C^\perp)\ge\ceil{\delta^\perp n},\qquad\Delta(A)<\ceil{\tau n};
\]
under the sampling law of \cref{dir:prop:legendrefinite}, conditioned on no repeated prime, the failure
probability is exponentially small in $n$.
\end{theorem}

\begin{theorem}[Approaching both Gilbert--Varshamov thresholds; \stG]\label{dir:thm:legendrelimit}
Fix $0<R<1$ and $\eta>0$, and let $X=\ceil{n^{2+\eta}}$ and $r=\floor{Rn}$. Primes in the intervals of
\cref{dir:thm:legendrecodes} can be chosen with $\rank A=r$,
\[
 \frac{d(C)}n\ge\mathsf H_2^{-1}(1-R)-o(1),\qquad\frac{d(C^\perp)}n\ge\mathsf H_2^{-1}(R)-o(1),\qquad
 \Delta(A)<\ceil{cn/\log_2n}
\]
for any fixed $c>2(R+\log_23)/\eta$; the success probability of the no-repeat sampling law tends to one.
\end{theorem}

\begin{theorem}[Fixed row primes; \stG]\label{dir:thm:legendreprefix}
Fix $0<R<1$, let $r=\floor{Rn}$ and let $q_i$ be the $i$th odd prime. For every large $n$ there are $n$ distinct
primes $p_j\in(X,2X)$, $X=\ceil{r^3(\log r)^4}$, such that $\rank A=r$, the two distances satisfy the bounds of
\cref{dir:thm:legendrelimit}, and $\Delta(A)<\ceil{cn/\log_2n}$ for any fixed $c>2(R+\log_23)$; the success
probability of $n$ independent draws from $\mu_X$ conditioned on distinctness tends to one.
\end{theorem}

\begin{corollary}[Robust identification and exact local uniformity; \stG]\label{dir:cor:legendre}
For the primes of \cref{dir:thm:legendrecodes}, distinct subsets $S\subseteq[r]$ give sign words differing at
at least $\ceil{\delta n}$ column primes, so $S$ is determined by its sign word with at most
$\floor{(\ceil{\delta n}-1)/2}$ sign errors. For $J\subseteq[n]$ with $\abs J<\ceil{\delta^\perp n}$, every sign pattern
on $J$ occurs for exactly $2^{r-\abs J}$ subsets $S$. For arbitrary $I,J$ the subsets $S\subseteq I$ realize exactly
$2^{\rank A_{I,J}}$ patterns on $J$, each $2^{\abs I-\rank A_{I,J}}$ times, hence at least
$2^{\min(\abs I,\abs J)-\Delta(A)}$ patterns.
\end{corollary}

\begin{proof}
The codeword of a nonzero difference has weight at least $d(C)$, and Hamming balls of radius below
$d(C)/2$ are disjoint. The projection to the coordinates $J$ is onto exactly when the columns in $J$ are
independent, since otherwise a nonzero dual word is supported in $J$. The general count is rank--nullity for
$\F_2^I\to\F_2^J$.
\end{proof}

The entropy thresholds are the classical random-code bounds \cite{GuruswamiRudraSudan2022Essential}, and the
unrestricted Gilbert--Varshamov bound is not improved; the content is the prime height. Quadratic-residue cyclic
codes, and Paley matrices built from $\legendre{j-i}p$ for one fixed $p$, are different constructions: here rows
and columns are labelled by primes and the denominator of the symbol varies.

\begin{lemma}[A rational window; \stG]\label{dir:lem:window}
Let $w(t)=(t-1)^4(2-t)^4$ on $[1,2]$, zero elsewhere, $W_X=\sum_pw(p/X)$ and $\mu_X(p)=w(p/X)/W_X$. There are
absolute $B_0,B_1,X_0$ such that for $X\ge X_0$ and every primitive nonprincipal character $\chi$ of conductor
$D$, with no restriction on $D$,
\[
 \abs{\dirE_{p\sim\mu_X}\chi(p)}\le B_0\frac{\log(2DX)}{\sqrt X},\qquad\max_p\mu_X(p)\le B_1\frac{\log X}X .
\]
\end{lemma}

\begin{proof}
By the prime number theorem $W_X\sim(X/\log X)\int_1^2w=X/(630\log X)$, which gives the second bound. Argue as
in the proof of \cref{dir:lem:bump} with $g_X(t)=w(t)/\log(Xt)$ on $[1,2]$, whose Mellin transform $G_X$ satisfies
$\abs{G_X(\sigma+it)}\ll(1+\abs t)^{-3}/\log X$ for $-1/2\le\sigma\le2$ by three integrations by parts in $\log t$, as $w$
and its first three derivatives vanish at the endpoints. Since $\chi$ is nonprincipal there is no residue at $1$,
and the contour shift with \cref{src:thm:dgrh} gives $\sum_m\Lambda(m)\chi(m)g_X(m/X)\ll\sqrt X\log(2D)/\log X+1$. Prime
powers with exponent at least two cost $O(\sqrt X)$, and dividing by $W_X$ gives the first bound.
\end{proof}

\begin{lemma}[Small bias and affine subspaces; \stU]\label{dir:lem:affine}
If a random $Z\in\F_2^a$ has $\abs{\dirE(-1)^{u\cdot Z}}\le\eps$ for all $u\ne0$, then every affine subspace $U$ of
codimension $b$ has $\dirP(Z\in U)\le2^{-b}+(1-2^{-b})\eps$.
\end{lemma}

\begin{proof}
For $U=z+U_0$ expand $\mathbf 1_U=2^{-b}\sum_{u\in U_0^\perp}(-1)^{u\cdot(Z-z)}$ and isolate $u=0$.
\end{proof}

\begin{lemma}[All rectangles at once; \stU]\label{dir:lem:rectangle}
If an $r\times n$ random binary matrix has independent columns, each with bias at most $\eps$ on every nonzero
linear character, then $\dirP(\Delta(A)\ge h)\le2^r3^n(2^{-h}+\eps)^h$ for every $h\ge1$.
\end{lemma}

\begin{proof}
Fix $I,J$ with $a=\abs I$, $b=\abs J$, $\min(a,b)\ge h$. A defect at least $h$ means rank at most
$t=\min(a,b)-h$, so some set $T\subseteq J$ of $t$ columns spans the columns of $A_{I,J}$ (pad a basis). Given the
columns in $T$, their span in $\F_2^a$ has codimension at least $h$, and the restriction to $I$ of every other
column still has bias at most $\eps$; by \cref{dir:lem:affine} and independence all $b-t\ge h$ other columns of
$J$ lie in the span with probability at most $(2^{-h}+\eps)^h$. There are $2^r$ choices of $I$, and assigning
each column to $T$, $J\setminus T$ or the complement of $J$ gives at most $3^n$ pairs $(J,T)$.
\end{proof}

\begin{proposition}[Finite two-sided failure bound; \stG]\label{dir:prop:legendrefinite}
Fix $1\le r\le n$, $1\le a,b\le n$ and $h\ge1$. Draw $r$ independent primes from $\mu_X$ and $n$ independent
primes from $\mu_{3X}$, independently, and let $\mathcal D$ be the event of no repetition in either list. For
$X\ge X_0$ there is an absolute $B$ such that, with
\[
 \eps=B\frac{n\log(12X)}{\sqrt X},\qquad M=\binom r2B_1\frac{\log X}X+\binom n2B_1\frac{\log(3X)}{3X},
\]
the probability that one of $\mathcal D$, $\rank A=r$, $d(C)\ge a$, $d(C^\perp)\ge b$, $\Delta(A)<h$ fails is at most
\[
 M+2^{r-n}V(n,a-1)(1+\eps)^n+2^{-r}\bigl(V(n,b-1)-1\bigr)(1+\eps)^r+2^r3^n(2^{-h}+\eps)^h,
\]
and if $M<1$ the conditional probability of a matrix failure given $\mathcal D$ is at most the sum of the last three
terms divided by $1-M$.
\end{proposition}

\begin{proof}
A collision of two independent draws from $\mu$ has probability $\sum_p\mu(p)^2\le\max_p\mu(p)$, so
$\dirP(\mathcal D^c)\le M$ by \cref{dir:lem:window}. Conditionally on distinct row primes, for $u\ne0$ the column
character $(-1)^{u\cdot A_{*,j}}=\prod_{u_i=1}\legendre{p_j}{q_i}$ is primitive and nonprincipal of conductor
$\prod_{u_i=1}q_i$, with logarithm at most $r\log(2X)$, so each bit of $uA$ is $1$ with probability at most
$(1+\eps)/2$, independently over the columns, and
$\dirP(\exists u\ne0:\mathrm{wt}(uA)<a)\le(2^r-1)V(n,a-1)((1+\eps)/2)^n$; as $a\ge1$ this event contains rank
deficiency. \Cref{dir:lem:rectangle} gives the rectangle term. Conditionally on distinct column primes, for
$v\ne0$ let $m_v=\prod_{v_j=1}p_j$ and $D_v=m_v$ or $4m_v$ as $m_v\equiv1$ or $3\pmod4$; then
$(-1)^{(Av^t)_i}=\legendre{m_v}{q_i}=\chi_{D_v}(q_i)$ with $\chi_{D_v}$ primitive and nonprincipal, accounting
for the reciprocity sign, and $\log D_v\le\log4+n\log(6X)$, so $\dirP(Av^t=0)\le((1+\eps)/2)^r$; a union bound over
the $V(n,b-1)-1$ nonzero $v$ of weight below $b$ gives the dual term. On $\mathcal D$ each failure event intersected
with $\mathcal D$ is bounded by its bound under the corresponding conditioning; summing and adding
$\dirP(\mathcal D^c)$ gives the claim.
\end{proof}

\begin{proof}[Proof of \cref{dir:thm:legendrecodes,dir:thm:legendrelimit}]
For the first, take $a=\ceil{\delta n}$, $b=\ceil{\delta^\perp n}$, $h=\ceil{\tau n}$. At $X=\ceil{Kn^2\log^2n}$,
$M=O(1/\log n)$ and $\eps\le3B/\sqrt K$ for large $n$. With $g_1=1-R-\mathsf H_2(\delta)$ and
$g_2=R-\mathsf H_2(\delta^\perp)$, choose $K$ so large that $\eps\le e_0$, where $0<e_0<1/2$ satisfies
$\log_2(1+e_0)<g_1/2$, $R\log_2(1+e_0)<g_2/2$ and $\tau\log_2(1/(2e_0))>R+\log_23+1$. Since
$V(n,k)\le2^{n\mathsf H_2(k/n)}$ for $k\le n/2$, the primal and dual terms are $2^{-\Omega(n)}$, and the rectangle
term is at most $2^{r+n\log_23}(2e_0)^h\le2^{-n}$. For the second, $M=O(n^{-\eta}\log n)$ and
$\eps=O_\eta(n^{-\eta/2}\log n)$; take $\gamma_n=n^{-\min(\eta,1)/4}$,
$a=\ceil{n\mathsf H_2^{-1}(1-R-\gamma_n)}$, $b=\ceil{n\mathsf H_2^{-1}(R-\gamma_n)}$ and
$h=\ceil{cn/\log_2n}$. Then $\eps=o(\gamma_n)$, the primal and dual terms are $2^{-n\gamma_n(1+o(1))}$ up to a
constant factor, and $2^r3^n(2^{-h}+\eps)^h\le2^{n(R+\log_23-c\eta/2+o(1))}=o(1)$; continuity of
$\mathsf H_2^{-1}$ gives the distances.
\end{proof}

With the row primes fixed, a bias bound for single columns no longer controls the dual distance; short dual
relations are excluded deterministically by a second use of \cref{src:thm:dgrh}.

\begin{lemma}[No short dual relations; \stG]\label{dir:lem:shortrelations}
In the setting of \cref{dir:thm:legendreprefix} there is an absolute $c_0>0$ such that for all large $r$ and any
distinct column primes in $(X,2X)$, every set of between $1$ and $H=\floor{c_0\sqrt{r/\log r}}$ columns of $A$ has
nonzero sum.
\end{lemma}

\begin{proof}
Let $S$ be such a set, $s=\abs S$, $m=\prod_{j\in S}p_j$ and $D=m$ or $4m$ its fundamental discriminant, so that
$\chi_D$ is primitive and nonprincipal as $m$ is squarefree and not a square. A zero column sum means
$\chi_D(q_i)=1$ for $i\le r$. With $y=q_r\sim r\log r$, for large $r$ one has $X>y$ and every prime in $(y/2,y)$ is a
row prime, so the mean of $\chi_D$ under $\mu_{y/2}$ would be one; but \cref{dir:lem:window} bounds it by
$\ll\log(2Dy)/\sqrt y\ll(s+1)\log r/\sqrt{r\log r}$, which is below $1/2$ for $s\le H$ when $c_0$ is small.
\end{proof}

\begin{proof}[Proof of \cref{dir:thm:legendreprefix}]
Since $\log\prod_iq_i=O(r\log r)$, \cref{dir:lem:window} gives column bias $\eps=O(r\log r/\sqrt X)=O(r^{-1/2}/\log r)$,
and repetitions have probability $O(n^2\log X/X)=o(1)$. The primal distance and the rectangle ranks follow as in
\cref{dir:prop:legendrefinite} with $\gamma_n=n^{-1/4}$, $a=\ceil{n\mathsf H_2^{-1}(1-R-\gamma_n)}$ and
$h=\ceil{cn/\log_2n}$, the rectangle bound being $2^{n(R+\log_23-c/2+o(1))}$. For the dual, with
$b=\ceil{n\mathsf H_2^{-1}(R-\gamma_n)}$, \cref{dir:lem:shortrelations} excludes on the no-repeat event every nonzero
$v$ of weight at most $H$; for a fixed set of $s>H$ columns, independence and Fourier inversion on $\F_2^r$ give
$\dirP(\sum_{j\in S}A_{*,j}=0)\le2^{-r}+\eps^s$, so relations with $H<s<b$ have probability at most
$2^{-r}V(n,b-1)+\sum_{s>H}(en\eps/s)^s$. The first term is $O(2^{-n\gamma_n})$, and with
$\alpha_n=en\eps/(H+1)=O_R(1/\sqrt{\log r})$ the sum is at most $\alpha_n^{H+1}/(1-\alpha_n)=o(1)$. Adding the
collision probability and dividing by $1-o(1)$ for the conditioning completes the proof.
\end{proof}

The proofs give a Monte Carlo construction: rejection sampling from the exact rational laws $\mu_X$, $\mu_{3X}$
produces a good matrix with probability $1-2^{-\Omega(n)}$ for fixed entropy gaps and $\tau$, but certifying the
distances or the worst rectangle may take time exponential in $n$, and no efficient decoder is claimed. Neither
the prime cutoff nor the rectangle defect is shown to be optimal.

\begin{example}\label{dir:ex:legendre}
Exact row reduction and complete weight enumerators, checked by the MacWilliams identities
\cite{MacWilliamsSloane1977Theory}, certify
\begin{center}\small
\begin{tabular}{@{}rrrrrl@{}}
\toprule
$n$ & $r$ & $d(C)$ & $d(C^\perp)$ & $X$ & row primes\\
\midrule
$12$ & $6$ & $3$ & $3$ & $300$ & sampled from $(X,2X)$\\
$24$ & $12$ & $5$ & $5$ & $1500$ & sampled from $(X,2X)$\\
$32$ & $16$ & $5$ & $5$ & $3000$ & sampled from $(X,2X)$\\
$24$ & $12$ & $5$ & $4$ & $2500$ & first $12$ odd primes\\
\bottomrule
\end{tabular}
\end{center}
with column primes in $(3X,6X)$ and every matrix of full rank $r$; for the first, all $2^{18}=262\,144$
rectangles were examined and the largest rank defect is $3$. These are selected witnesses, not estimates of
the success probability.
\dirfiles{\artifact{prime_legendre_codes_2026-11-27/examples.json},
\artifact{prime_legendre_codes_2026-11-27/examples_summary.csv},
\artifact{prime_legendre_codes_2026-11-27/checks.txt}}
\end{example}

\subsection{Open questions}\label{prd:sec:open}
{\raggedright\sources{squarefree_prime_products; orthogonal_nullity_and_prime_products; squarefree_factorizations; hermite_thresholds_and_quadratic_primes; disjoint-prime-products; entropy_optimal_prime_factorizations_2026-11-27; grh_entropy_products_2026-11-27; grh_growing_factorizations; grh_quadratic_fibre; grh_prime_routing_2026-11-27; prime_character_codes_2026-11-27; prime_sparse_recovery_2026-11-27; prime_sparse_recovery_2026-11-27_2; grh_prime_fingerprints_2026-11-27; hermite_threshold_research; grh_quarter_density; prime_cubature_grh_2026-11-27; prime_legendre_codes_2026-11-27; grh_extensions_2026-11-27_3}}

\begin{enumerate}
\item \emph{Short squarefree products at the logarithm-squared scale.} At $y=C(\log q)^2$, can
 every unit be reached after majority retention with $O(\log q/\log\log q)$ distinct prime
 factors? This would give representatives of polynomial size at the smoothness scale of
 Corollary~\ref{dir:cor:squarefree}. For $y=(\log q)^A$, $A>2$, what is the optimal constant in
 front of $\log q/\log\log q$ for pointwise relative counts? For typical residues it is $1/(A-1)$ for
 almost all prime moduli (Theorems~\ref{dir:thm:typical} and~\ref{dir:thm:typicalsharp}); pointwise,
 Theorem~\ref{dir:thm:short} and Corollary~\ref{dir:cor:almostshort} give $2(A-1)/(A(A-2))$ under
 Dirichlet GRH and $2/(A-1)$ for almost all prime moduli. Can the typical constant
 $(A-1)/(A(A-2))$ under Dirichlet GRH be lowered to $1/(A-1)$ for every modulus? For one product,
 what happens when $\min(k,m-k)=o(\log q)$? Below density one third and in the sparse range
 $B\log q\le h=o(m)$, the fibre distances of the characters of order at most $\floor{1/a}$ decide mixing
 (Theorem~\ref{dir:thm:allmix}) and compound Poisson laws on a bounded quotient give every limiting
 profile (Theorem~\ref{dir:thm:profiles}); what replaces them for dense quotas $h\asymp m$, and do they
 hold uniformly as $a\to0$?
\item \emph{Mixing at density one quarter.} The support at retained density $1/4$ is given by
 Theorem~\ref{dir:thm:quarter}, and for one product the sparse-quota profiles at every density are
 classified by Theorem~\ref{dir:thm:profiles}, with an $O(1/h)$ comparison in the critical window that
 needs no bound on the number of exceptional primes (Theorem~\ref{dir:thm:fpc}). For $r\ge2$ products,
 and in particular for growing $r$, which statistics determine the approach to uniformity across the
 occupied quotient classes, beyond the relative counts within each class, once the number of exceptional
 primes is not $o(\sqrt m)$ as in Proposition~\ref{dir:prop:sparseexc}?
\item \emph{Constants for factorizations.} For one tuple of moduli and coverage alone the leading
 constant is the entropy threshold of Theorem~\ref{dir:thm:entropy}. What is the second-order term,
 between the lower bound of order $\log\log M$ and the upper bound $O(\log M/\log B)$, and can a fixed
 height factor $B$ attain the leading constant exactly? Does an entropy law hold for a growing number
 of factors, with deletions, or with comparable numbers of representations? For a fixed relative
 error $Q^{-A}$ and a fixed deletion fraction or retained density, what is the best leading constant
 in the logarithmic quota of Theorems~\ref{dir:thm:factorization}, \ref{dir:thm:minority}
 and~\ref{dir:thm:disjoint-growing}, also for growing $r$? Can the largest prime be made smaller than
 $K\log^2Q$ for a fixed number of factors while keeping the pointwise bound?
\item \emph{Hermite thresholds.} Lemma~\ref{dir:lem:hamming}(b) determines $b_k(a)$ in the windows
 $a=h_j+\lambda/\sqrt k$ with $\lambda\le0$. What is the window function for $\lambda>0$, where rare
 extreme layers may be used, and what is the next term in Lemma~\ref{dir:lem:hamming}(b)? Can
 correlations among the prime character sums, or a witness optimized for the conductor-weighted
 Fourier norm, improve the powers of $k$ in Theorems~\ref{dir:thm:signs}
 and~\ref{dir:thm:signswindow}, and is there an arithmetic lower bound for the least prime?
\item \emph{Computing the quotient.} Can the bounded quotient of Theorem~\ref{dir:thm:quotient}
 be computed from the retained primes in time polynomial in $\log q$, without enumerating the
 characters modulo $q$? Once it is known, Corollary~\ref{dir:cor:anydensity} decides the support of
 disjoint products in polynomial time.
\item \emph{Routing capacity.} At the cutoff $C(\log q)^2$, Corollary~\ref{dir:cor:routing} routes
 $\gg\log q/\log\log q$ disjoint pairs while the counting bound allows $O(\log q)$; what is the
 capacity? Can the constant $2/(A-2)$ in the length bound of Theorem~\ref{dir:thm:routing} be lowered
 at the same deletion threshold?
\item \emph{Prime codes.} Can the number of small primes in Theorems~\ref{dir:thm:codes}
 and~\ref{dir:thm:primedecoder} be made nearly linear in the sparsity while keeping the sharp
 corruption fraction, and is the cost $s^6$ of exact rigidity necessary? For moduli without a
 subgroup of characters of order $2k$ the minimum distance may exceed the coherence bound; can a
 tractable decoder use this? Is prime height $(\log Q)^2$ necessary for the fingerprints of
 Theorem~\ref{dir:thm:fingerprint}, and how small can the list size be near capacity? Are the prime
 cutoff $O(n^2\log^2n)$ and the rectangle defect of Theorem~\ref{dir:thm:legendrecodes} optimal for
 Legendre codes?
\item \emph{Prime designs.} Can the height exponent of Theorem~\ref{dir:thm:cubature} be lowered at a fixed
 positive deletion fraction, is there a sharp obstruction in degree $k\ge2$ analogous to
 Theorem~\ref{dir:thm:threelayer}, and can the bounded reweighting be replaced by equal weights on distinct
 primes?
\end{enumerate}
